% defumat user guide.
%
% This is the SOURCE, not a conversion of anything -- build with:
%     xelatex docs/features.tex     (twice, for the table of contents)
% xelatex rather than pdflatex because the guide carries Unicode in a few
% places and because fontspec gives it a decent monospace for the code.
\documentclass[11pt,a4paper]{article}

\usepackage{fontspec}
\usepackage[margin=2.4cm]{geometry}
\usepackage{amsmath,amssymb}
\usepackage{booktabs}
\usepackage{array}
\usepackage{longtable}
\usepackage{xcolor}
\usepackage[most]{tcolorbox}
\tcbuselibrary{listings,skins,breakable}
\usepackage{hyperref}
\usepackage{microtype}

\setmainfont{DejaVu Serif}[Scale=0.94]
\setmonofont{DejaVu Sans Mono}[Scale=0.80]

\definecolor{codebg}{HTML}{F7F7F4}
\definecolor{coderule}{HTML}{D8D8D0}
\definecolor{kw}{HTML}{0B5FA5}
\definecolor{str}{HTML}{097A5A}
\definecolor{cmt}{HTML}{8A8A80}
\definecolor{link}{HTML}{1A4F8A}
\definecolor{warnbg}{HTML}{FDF6EC}
\definecolor{warnrule}{HTML}{E0C08A}

\hypersetup{colorlinks=true, linkcolor=link, urlcolor=link,
            pdftitle={defumat: what it computes, and how to ask for it}}

% ---- code boxes: wrapped, so nothing can run off the page ----
\lstdefinestyle{py}{
  language=Python,
  basicstyle=\ttfamily\small,
  keywordstyle=\color{kw}\bfseries,
  stringstyle=\color{str},
  commentstyle=\color{cmt}\itshape,
  showstringspaces=false,
  breaklines=true,                    % the whole point: no overflow
  breakatwhitespace=false,
  postbreak=\mbox{\textcolor{coderule}{$\hookrightarrow$}\space},
  columns=fullflexible,
  keepspaces=true,
  upquote=true,
}
\newtcblisting{pycode}{
  listing only, listing style=py, breakable,
  colback=codebg, colframe=coderule, boxrule=0.5pt, arc=2pt,
  left=5pt, right=5pt, top=4pt, bottom=4pt,
  before skip=8pt, after skip=8pt,
}
\lstdefinestyle{sh}{basicstyle=\ttfamily\small, breaklines=true,
  columns=fullflexible, keepspaces=true, upquote=true}
\newtcblisting{shcode}{
  listing only, listing style=sh, breakable,
  colback=codebg, colframe=coderule, boxrule=0.5pt, arc=2pt,
  left=5pt, right=5pt, top=4pt, bottom=4pt,
  before skip=8pt, after skip=8pt,
}

% ---- a box for the "Refuses" notes, which are load-bearing here ----
\newtcolorbox{refuses}{
  colback=warnbg, colframe=warnrule, boxrule=0.5pt, arc=2pt,
  left=6pt, right=6pt, top=4pt, bottom=4pt,
  before skip=7pt, after skip=7pt, breakable,
}

\newcommand{\qe}{Quantum ESPRESSO}
\newcommand{\code}[1]{\texttt{\small #1}}
\setlength{\parskip}{0.45em}
\setlength{\parindent}{0pt}

\title{\textbf{defumat}\\[4pt]
  \large what it computes, and how to ask for it}
\author{\normalsize Plane-wave DFT in Python and JAX, validated against \qe}
\date{\normalsize\today}

\begin{document}
\maketitle
\thispagestyle{empty}

\begin{abstract}
\noindent
User documentation for the feature set: every capability, the equation behind
it, a snippet that runs it, what it was validated against, and \emph{what it
refuses}. Two conventions run through the whole document.

\textbf{Everything is validated against \qe\ on the same input}, and the
agreement is quoted per feature rather than claimed in general. Where a number
is missing it is because the comparison does not exist, and that is said rather
than hidden.

\textbf{Anything not implemented is refused by name}, with an error saying what
is missing --- never silently approximated. The ``Refuses'' notes are as much of
the documentation as the features: they are the promise that a run which starts
is a run whose physics is there.
\end{abstract}

\vspace{4pt}
\hrule
\vspace{6pt}

\textbf{One equation is written down and the rest are derived from it.} Almost
everything below is a derivative of the total energy, taken by the compiler
rather than by hand:
\begin{equation*}
E[\{\psi\},\tau,\varepsilon] = T_s + E_{\mathrm H} + E_{xc}
  + E_{\mathrm{loc}} + E_{\mathrm{nl}} + E_{\mathrm{Ewald}}
\end{equation*}

\begin{center}
\small
\begin{tabular}{@{}lll@{}}
\toprule
\textbf{quantity} & \textbf{what it is} & \textbf{how it is obtained} \\
\midrule
$v_{xc}(\mathbf r)$ & $\delta E_{xc}/\delta n$ & \code{jax.grad} of the energy density \\
$\mathbf F_I$ & $-\partial E/\partial\boldsymbol\tau_I$ & \code{jax.grad} at frozen $\psi$ \\
$\sigma_{ij}$ & $-\Omega^{-1}\partial E/\partial\varepsilon_{ij}$ & \code{jax.grad} along a strain \\
$C_{I\alpha,J\beta}$ & $\partial^2 E/\partial\tau_{I\alpha}\partial\tau_{J\beta}$ & \code{jvp} of the force \\
$Z^{\ast}_{I,\alpha\beta}$ & $\partial F_{I\alpha}/\partial\mathcal E_\beta$ & \code{jvp} of the force along the field response \\
$\partial\epsilon_{ij}/\partial\tau$ & Raman & \code{jvp} of the second-order energy \\
\bottomrule
\end{tabular}
\end{center}

The ``frozen $\psi$'' is what makes this exact rather than approximate: at
self-consistency the energy is stationary with respect to the wavefunctions, so
the terms from $\partial\psi/\partial\lambda$ vanish. That is the
Hellmann--Feynman argument, and taking it one order up is the $2n{+}1$ theorem,
which is why a \emph{third} derivative needs only first-order wavefunctions.

\clearpage
\tableofcontents
\clearpage

% =====================================================================
\section{Input and invocation}

defumat reads \textbf{\code{pw.x} input files}. The same file that runs in
\qe\ runs here, which is what makes every comparison in this document possible.

\begin{pycode}
from defumat import Calculator

calc = Calculator.from_file("scf.in", pseudo_dir="pseudo", conv_thr=1.0e-8)
result = calc.get_scf()

print(result.total_energy, "Ry in", result.iterations, "iterations")
\end{pycode}

\code{conv\_thr} means what it means in a \code{pw.x} input: it is compared
against QE's \code{dr2}, the Hartree energy of the density residual.

It need not be given at all. \code{from\_file} and \code{from\_text}
\textbf{adopt the input's own \code{\&electrons} namelist} ---
\code{conv\_thr}, \code{mixing\_beta}, \code{mixing\_mode},
\code{mixing\_fixed\_ns} and \code{electron\_maxstep} --- so an input file
that states how to converge itself converges the same way here as under
\code{pw.x}. A keyword argument given to the constructor still wins over the
file, which is how the snippet above overrides it. One variable of that
namelist is read by \code{pw.x} and deliberately \emph{not} adopted:
\code{diagonalization}, because this package offers one eigensolver, so
honouring \code{diagonalization = 'cg'} would turn a run that works into an
error, and mapping it silently onto Davidson would be the kind of substitution
refused everywhere else here. Name a solver at construction instead.

\textbf{An adopted default yields to a workflow that chose its own}, and until
2026-09-20 it did not. A threshold means the same thing wherever it is named,
which is why it is forwarded at all, but meaning the same thing is not wanting
the same value: twenty entry points declare a \code{conv\_thr} between
$10^{-8}$ and $10^{-12}$ and say in their own documentation why --- a
magnetocrystalline anisotropy is a difference of band-energy sums in the fifth
decimal of an eV, an effective mass is a second difference of eigenvalues, a
Berry phase is a product of overlaps around a closed loop. An input file that
states \code{pw.x}'s own $10^{-6}$ was replacing every one of them. So a value
adopted from the file, or given at construction, reaches an entry point only
where that entry point repeats \code{run\_scf}'s own default, meaning it has no
opinion; where it states one of its own, the workflow's number stands. Naming
it at the call site --- \code{get\_anisotropy(soc, conv\_thr=1.0e-6)} ---
still reaches it, which is the difference between a default and a refusal.

\textbf{Tightening is always allowed}, and only loosening is withheld: a
threshold is a bound on an error, so a \code{conv\_thr} of $10^{-12}$ given to
the constructor still reaches a workflow whose own is $10^{-10}$. That
asymmetry belongs to thresholds alone --- \code{max\_iterations} of 40 against
100 is not better or worse, it is what the workflow wants --- and it is what
keeps an input stating $10^{-10}$ from being withheld from an ultracell run,
whose own $10^{-8}$ is a second driver's default for the same quantity rather
than a tightening.

On the committed tetragonal cobalt cell the substitution is worth little,
and the reason is worth knowing: both legs of the force theorem diagonalise the
\emph{same} frozen density on the same k-points at the same threshold, so the
eigenvalue error is common mode and the difference keeps only what does not
cancel. Measured, $10^{-6}$ in place of $10^{-10}$ moves the in-plane leg by
$2.06 \times 10^{-5}$~meV and the out-of-plane leg by $1.97 \times
10^{-5}$~meV, of which $9.0 \times 10^{-7}$~meV survives into an anisotropy of
1.235~meV.

\textbf{A band count does not cross between the two legs at all.} The force
theorem is two calculators, and \code{nbnd} is a property of the system whose
bands are being counted: a spinor band holds one electron where a collinear one
holds two, so the same cobalt atom needs 9 bands collinear and 16 as a spinor.
A \code{nbnd} given to the scalar-relativistic calculator used to be forwarded
into the noncollinear run, where it wins over the spinor input's own, and 9
spinor bands for 9 electrons is a filled-band insulator rather than a metal:
the anisotropy reads \textbf{7.926~meV against 1.235}, and the free-energy
anisotropy, which should differ from it by more than half, collapses onto the
same number because there are no empty states for the Fermi level to sit among
and the smearing entropy is zero. The spinor calculator's own \code{nbnd}
crosses in its place, so stating it where it belongs still works.

\textbf{The adopted namelist reaches a relaxation too}, and did not until
2026-09-11: \code{electron\_maxstep}, \code{diago\_full\_acc} and
\code{mixing\_fixed\_ns} were read off the file and then dropped on the way into
\code{get\_relax}, so every SCF inside a relaxation ran at the default hundred
iterations whatever the input asked for. On a DFT+U relaxation the dropped
\code{mixing\_fixed\_ns} decides which minimum of the $+U$ functional the run
lands in, so the relaxed \emph{geometry} could differ with nothing saying the
request had been ignored.

\begin{refuses}
\textbf{An input variable that changes the physics and is not implemented is
refused by name}, never read and ignored --- a run that quietly drops one is
wrong by far more than the agreement with \code{pw.x} this code is checked to.
That covers a charged cell (\code{tot\_charge}), the saw-tooth field and the
dipole correction (\code{tefield}, \code{dipfield}), the Berry-phase finite
field (\code{lelfield}), the isolated-system corrections
(\code{assume\_isolated}), two chemical potentials (\code{twochem}, in either
namelist), per-orbital atomic occupations
(\code{one\_atom\_occupations}), the constant-cutoff kinetic functional
(\code{qcutz}), a hand-pinned FFT grid (\code{nr1}--\code{nr3s}), and the four
\textbf{symmetry} switches \code{no\_t\_rev}, \code{force\_symmorphic},
\code{use\_all\_frac} and \code{nosym\_evc}.

Those last four were read by \emph{nothing} until 2026-09-11. Each says that
some part of the group is not a symmetry of the state being asked for, so
ignoring one left the k-mesh reduced with operations the input had excluded and
\code{sym\_rho} averaging over them --- a total energy that will not match a
benchmark generated from the same file, with nothing saying why. Anything QE
defaults to is silent, so an ordinary input never reaches the refusal.

\textbf{Two more are refused because they are \emph{redundant} rather than
unimplemented}, and both used to be dropped in silence. The first is the
DFT+Hubbard spelling the \code{\&system} namelist carried before QE 7.1 ---
\code{lda\_plus\_u}, \code{lda\_plus\_u\_kind}, \code{Hubbard\_U},
\code{Hubbard\_U\_back}, \code{Hubbard\_J0}, \code{Hubbard\_J},
\code{Hubbard\_V}, \code{U\_projection\_type},
\code{Hubbard\_parameters} and \code{backall}. DFT+U is implemented here and
is read from the \code{HUBBARD} card, so the old names parsed and reached
nothing: an input asking for a $U$ converged as plain LSDA or PBE and reported
success, and what was dropped is the whole $U$-on against $U$-off gap.
\code{pw.x} 7.5 stops on the same set, and on the \emph{value} rather than the
name, so \code{lda\_plus\_u = .false.} and \code{Hubbard\_U = 0.0} pass in
both codes. The second is \code{ibrav} $\neq 0$ together with a
\code{CELL\_PARAMETERS} card: the card was never looked at, and because
\code{ATOMIC\_POSITIONS alat} is scaled by the \code{ibrav} cell's own
\code{alat} the positions went with it, so pasting a relaxed cell back under an
input that still carries its \code{ibrav} --- the ordinary way a vc-relax is
restarted by hand --- ran the pre-relaxation geometry and reported a total
energy, a force and a stress for it. Set \code{ibrav = 0} to use the card.
Neither refusal is reached by any of the 1194 inputs in this repository, in
\code{benchmarks/} or in QE's own vendored test-suite.

Two more say the lattice parameter twice, and \code{pw.x} stops on both. Giving
\code{celldm} \emph{and} \code{A} is one lattice constant written in bohr and
in angstrom, which is what a half-finished conversion leaves behind, and taking
\code{celldm} silently converges a ground state for whichever of the two was
not meant. Giving a \code{CELL\_PARAMETERS bohr} or \code{angstrom} card
\emph{and} \code{celldm(1)} is the same thing one step along: the card already
carries the lattice parameter, so \code{alat} became \code{celldm(1)} rather
than the length the vectors set, and every \code{ATOMIC\_POSITIONS alat}
coordinate was scaled by a parameter the cell does not come from. Write the card
in \code{alat} units to keep \code{celldm(1)}, or drop \code{celldm(1)} to
keep the card. No \code{pw.x} input in this repository or in the vendored
test-suite carries either pair.

One parsing trap sits under all of this and is worth stating once: Fortran's
list-directed logical input accepts an optional \code{.}, then a \code{T} or
an \code{F}, then anything, so \code{lberry = F} is legal and is false. The
bare spellings were missing from the converter, so \code{F} survived as the
\emph{string} \code{'F'} and every truth test read it as true.
\end{refuses}

\subsection*{One object, bound methods}

A \code{Calculator} is a system together with its pseudopotentials, and every
calculation in this document is a method on it. The pseudopotentials are read
from the names the \code{ATOMIC\_SPECIES} card gave, resolved against
\code{pseudo\_dir} --- which \code{from\_file} defaults to the input file's own
directory --- so a script names an input and then names physics:

\begin{pycode}
calc = Calculator.from_file("scf.in", pseudo_dir="pseudo")

bands = calc.get_bands(kpoints=path)      # runs the SCF first if none is cached
dos = calc.get_dos(grid=(8, 8, 8))
forces = calc.get_forces()
epsilon = calc.get_dielectric_tensor()
phonons = calc.get_phonons()

bands.plot()                              # and each result draws itself
\end{pycode}

Three properties are worth knowing before relying on it.

\textbf{The ground state is cached, and computed on demand.} A method that needs
one and finds no cache runs an SCF, printing a line to stderr saying so;
\code{get\_scf()} is the explicit path and is the same code. The cache is a
\emph{single slot} holding the result together with the options that produced
it --- calling \code{get\_scf} again with different options reruns and replaces
it --- because what it holds is the wavefunctions, and a cache keyed by option
sets would quietly hold several sets of them. \code{calc.scf\_result} reads the
slot without triggering anything, and a state that did not converge is refused
rather than differentiated.

\textbf{Nothing mutates.} \code{with\_positions}, \code{with\_cell},
\code{with\_spin} and \code{with\_moments} return a \emph{new} calculator with an empty cache; the
converged state crosses as a starting guess where that is sound, which for
\code{with\_spin} is the spin-regime continuation described under
\emph{Spin, magnetism and spin-orbit} and converges in one iteration rather
than twenty-five. A cached result under a moved atom
would otherwise answer for the geometry it was converged at.

\textbf{The refusals pass through untouched.} Every restriction in this document
--- a metal's Raman tensor, a spiral with spin-orbit coupling, a phonon away from
$\Gamma$ on an augmented dataset --- raises through the bound method exactly as
through the function beneath it.

The functional entry points the rest of this document names --- \code{run\_scf},
\code{run\_bands}, \code{dielectric\_tensor} and the others --- are unchanged and
remain the way to drive a calculation whose state is being managed by hand:

\begin{pycode}
from defumat.io.pwin import read_pw_input
from defumat.pseudo import read_upf
from defumat.scf.driver import Calculation, run_scf
from defumat.system import build_system

system = build_system(read_pw_input("scf.in"))
pseudos = tuple(read_upf(f"pseudo/{s.pseudo_file}") for s in system.structure.species)
result = run_scf(system, pseudos, conv_thr=1.0e-8)
\end{pycode}

What such a call has to carry with the density is the rest of the converged
state --- \code{becsum} for a PAW dataset, \code{ns} under a Hubbard $U$,
\code{tau} under a meta-GGA. None of the three can be rebuilt from the density,
and the entry points refuse without them rather than computing something else.
The calculator passes them.

\subsection*{Every method, in one place}

The rest of this document introduces each quantity where its physics belongs,
and names the functional entry point beneath it. This is the index the other way
round --- what to \emph{call}, for a reader who knows the physics and wants the
line. It is complete: a method absent from it is a documentation bug, and a test
says so (\code{tests/unit/test\_facade\_documentation.py}).

\begin{center}
\small
\begin{longtable}{@{}ll@{}}
\toprule
\multicolumn{2}{@{}l}{\textit{The ground state}} \\
\midrule
\code{get\_scf} & the self-consistent field loop, or the cached result \\
\code{get\_elk\_seed} & Elk's converged density on this run's grid, as a guess \\
\addlinespace
\multicolumn{2}{@{}l}{\textit{Band structure and densities of states}} \\
\midrule
\code{get\_bands} & diagonalise on a k-path at the converged density \\
\code{get\_nscf} & diagonalise on a k-grid at the converged density, and occupy it \\
\code{get\_dos} & the density of states, on a denser grid if asked \\
\code{get\_pdos} & the projected density of states, with L\"owdin charges \\
\addlinespace
\multicolumn{2}{@{}l}{\textit{Derivatives of the energy}} \\
\midrule
\code{get\_forces} & the forces on the atoms \\
\code{get\_stress} & the stress tensor \\
\code{get\_relax} & relax the geometry, and the whole optimisation with it \\
\addlinespace
\multicolumn{2}{@{}l}{\textit{Response, vibrations and spectra}} \\
\midrule
\code{get\_band\_velocities} & $\mathrm{d}\varepsilon/\mathrm{d}\mathbf k$ for every band \\
\code{get\_effective\_mass} & $(1/m^{*})_{ab}$ at one k-point, in $1/m_e$ \\
\code{get\_angular\_momenta} & $\langle L\rangle$, $\langle S\rangle$, $\langle J\rangle$ on every atom \\
\code{get\_dielectric\_tensor} & $\epsilon^{\infty}$, with no empty states \\
\code{get\_born\_charges} & the Born effective charges $Z^{*}$ \\
\code{get\_phonons} & the dynamical matrix at $\Gamma$, and its modes \\
\code{get\_phonons\_at\_q} & the dynamical matrix at one wavevector \\
\code{get\_raman\_tensors} & $\mathrm{d}\epsilon/\mathrm{d}\tau$, per atom \\
\code{get\_vibrational\_spectrum} & per-mode Raman and infrared activities \\
\code{get\_strain\_response} & the first-order response to a homogeneous strain \\
\code{get\_elastic\_constants} & $C_{ijkl}$ \\
\code{get\_piezoelectric\_tensor} & $e_{(k)ij}$, clamped-ion, in C/m$^2$ \\
\code{get\_electrostriction} & $\mathrm{d}\chi/\mathrm{d}(\text{strain})$ and its four tensors \\
\code{get\_absorption} & an optical spectrum, with excitons \\
\code{get\_optical\_conductivity} & $\sigma_{ab}(\omega)$, the Kerr angle, the anomalous Hall \\
\code{get\_structure\_factors} & $F(\mathbf H)$, X-ray and magnetic \\
\code{get\_stm} & a Tersoff--Hamann scanning-tunnelling image \\
\code{get\_sts} & the same, over an energy axis: a tunnelling spectrum \\
\code{get\_vertical\_transport} & tunnelling \emph{through} a 2D material \\
\code{get\_momentum\_transport} & which k-points that current comes out of \\
\code{get\_nesting} & $N(\mathbf q)$, the Fermi-surface nesting function \\
\code{get\_spin\_susceptibility} & $\chi^{+-}(\mathbf q,\omega)$ \\
\code{get\_magnon\_dispersion} & $\omega(\mathbf q)$ along a path \\
\code{get\_shift\_current} & the bulk photovoltaic effect, in A/V$^2$ \\
\code{get\_shg} & $\chi^{(2)}(-2\omega;\omega,\omega)$, in pm/V \\
\addlinespace
\multicolumn{2}{@{}l}{\textit{Topology, polarization and magnetism}} \\
\midrule
\code{get\_berry\_curvature} & $\Omega(\mathbf k)$ on a plane mesh \\
\code{get\_chern} & the Chern number of one plane \\
\code{get\_z2} & the 2D $Z_2$ invariant, by Wilson loops or parity \\
\code{get\_z2\_3d} & the four $Z_2$ indices of a 3D crystal \\
\code{get\_polarization} & the Berry-phase polarization \\
\code{get\_orbital\_magnetization} & $M_{\rm orb}$, by the modern theory \\
\code{get\_magnetoelectric\_tensor} & $\alpha_{ij} = \mathrm{d}P_i/\mathrm{d}B_j$ \\
\code{get\_anisotropy} & the magnetocrystalline anisotropy \\
\code{get\_torque} & the magnetic torque, and $K$ from one angle \\
\code{get\_first\_order\_soc} & the spin-orbit expectation at coupling-free states \\
\code{get\_force\_theorem} & one direction of the anisotropy: its band energy alone \\
\addlinespace
\multicolumn{2}{@{}l}{\textit{Spin spirals}} \\
\midrule
\code{get\_spiral\_scan} & one SCF per spiral wavevector \\
\code{get\_spiral\_relaxation} & the spiral wavevector relaxed down $\mathrm{d}E/\mathrm{d}\mathbf q$ \\
\addlinespace
\multicolumn{2}{@{}l}{\textit{Modulations over many cells}} \\
\midrule
\code{get\_ultracell} & a density or potential modulated over $N$ unit cells \\
\code{get\_ultracell\_stm} & a Tersoff--Hamann image of that modulation \\
\code{get\_ultracell\_sts} & a tunnelling spectrum across that modulation \\
\code{get\_ultracell\_transport} & tunnelling through it, tip to substrate \\
\bottomrule
\end{longtable}
\end{center}

A command line covers the common workflows:

\begin{shcode}
python3 -m defumat.cli inspect <qe-output>   # summarise what the parser reads
python3 -m defumat.cli dos     <input>       # density of states
python3 -m defumat.cli pdos    <input>       # projected DOS
python3 -m defumat.cli relax   <input>       # structural relaxation
python3 -m defumat.cli stress  <input>       # stress tensor
python3 -m defumat.cli spiral  <input>       # spin-spiral scan
\end{shcode}

Standard \code{pw.x} variables --- \code{ecutwfc}, \code{ecutrho}, \code{ibrav},
\code{nbnd}, \code{occupations}, \code{degauss} and the rest --- mean what they
mean there and are not re-documented here. A handful of knobs are worth knowing
because they are specific to this code or to a feature below:

\begin{center}
\small
\begin{tabular}{@{}ll@{}}
\toprule
\code{mbj\_c} & fixes Tran--Blaha's $c$ instead of taking the cell average \\
\code{london\_s6}, \code{london\_rcut} & Grimme D2's scaling and real-space cutoff \\
\code{cell\_dofree} & which cell degrees of freedom a vc-relax may move \\
\code{starting\_ns\_eigenvalue} & seeds the DFT+U occupation matrix \\
\code{spiral\_q} & the spin-spiral wavevector, in lattice coordinates \\
\code{LOCAL\_MAGNETIC\_FIELDS} & a card with no \code{pw.x} counterpart \\
\bottomrule
\end{tabular}
\end{center}

\textbf{Units are Rydberg atomic units throughout} (Ry, bohr), matching QE. The
only layer that speaks anything else is \code{io/}: Hubbard $U$ in the
\code{HUBBARD} card is in eV and is converted at the boundary, as \code{pw.x}
does it.

\subsection*{What a run will cost, before it is run}

\code{Calculation} builds the G-vectors, the plane waves and the projectors on
the device, so an input too large for the machine dies in setup without saying
what was too large. \code{Calculator.estimate()} answers that from the cutoffs
and the datasets alone, allocating nothing:

\begin{lstlisting}[language=Python]
print(Calculator.from_file("scf.in").estimate().report())
\end{lstlisting}

\noindent or \code{python3 -m defumat.cli size scf.in}. \textbf{It sizes the
run the input describes}, which means it reads that input's own
\code{diago\_david\_ndim} and resolves \code{k\_batch} the way the SCF would ---
an estimate that quietly assumed a library default would report a run that does
not happen, and nobody relying on one is in a position to notice. The
arguments and the CLI flags are \emph{overrides} of the file, and the report's
header states which assumptions it used. It reports $n_{gm}$,
$n_{gms}$, $n_{pwx}$, $n_{bnd}$, $n_{kb}$ and both FFT grids --- \emph{exactly},
from the same predicates the setup selects with, which is asserted against a
real \code{Calculation} rather than trusted --- and a floor on the bytes for the
wavefunctions, the projectors, the eigensolver's subspace and the fields. The
G-vector enumeration runs in slabs, so an input destined for a GPU can be sized
on a laptop.

Two lines of that report are worth reading before believing a hand estimate.
The Davidson subspace is $2\,n_{vecx}\,n_{dim}$ for $\psi$ and $H\psi$ at QE's
$n_{vecx}=4\,n_{bnd}$, plus the Ritz block, and it is usually the largest
single array --- on a 157-atom slab at $E_{\mathrm{cut}}=60$~Ry it is 139~GB
against 23~GB of wavefunctions. And the eigensolver is \emph{not} doubled by
spin, because the channels are solved one after another, where the
wavefunctions are; the report separates the two.

\noindent The one lever that changes that floor without changing the physics is
\code{diago\_david\_ndim} --- QE's own name for it, read from
\code{\&electrons}, and settable per call or at construction:

\begin{lstlisting}[language=Python]
calc = Calculator.from_file("scf.in", david=2)   # or diago_david_ndim = 2
\end{lstlisting}

\noindent It halves $\psi$ and $H\psi$ exactly, for one more SCF iteration on
silicon and the same total energy to 7.3e-11~Ry --- the trade QE's
documentation describes. Note that \code{INPUT\_PW.txt} gives its default as 2
and is stale: \code{input\_parameters.f90} sets 4 and \code{input.f90} assigns
it straight through, so the number \code{pw.x} runs is \textbf{4}, and that is
the default here.

\noindent \textbf{The subspace is capped at the space it lives in}, which is
$n_{pol}\,\min_{\mathbf k} n_{pw}$ --- the \emph{smallest} sphere on the k-set, not
$n_{pwx}$, which is the padded maximum and sits above the k-points that go
singular. A small cell asked for many bands is where this bites: silicon at
$E_{\mathrm{cut}}=12$~Ry folded to a 32 k-point ultracell holds between 169 and
192 plane waves, so $n_{bnd}=48$ at the default asks for 192 vectors in a
169-dimensional space, and 25 of the 32 overlaps came back non-finite where the
capped solve leaves none. A cap removes the case where the subspace is
\emph{larger} than the space and not the case where it \emph{equals} it: at
$n_{bnd}=80$ on the same cell the cap is the whole space at the smallest
k-point, one overlap still goes singular, and the canonical-orthogonalisation
fallback carries it. Two cases no cap can rescue are refused by name instead ---
more bands than the space holds, which is QE's own
\code{c\_bands.f90:286}, and $2\,n_{bnd}$ not fitting, which leaves Davidson no
room to expand (\code{cegterg.f90:125}) and wants a direct diagonalisation of
the space that is not written here.

\subsection*{A converged state on disk}

\code{result.save(path)} writes the state and \code{SCFResult.load(path,
system=...)} reads it back for \code{run\_scf(starting\_from=...)}, so a job
that hits its wall clock does not cost the run. The system is supplied rather
than stored --- a resume already has the input file --- and is checked against
a fingerprint. \code{run\_relax(checkpoint\_dir=...)} does the same every ionic
step \emph{including the optimizer's own history}, which is the half that is
easy to omit and most of the value: BFGS earns its convergence rate from the
inverse Hessian and the trust radius, so a resume from positions alone takes
its next step as if it were the first. Stopping a six-step relaxation at step
two and resuming costs $2+4$ steps, not $2+6$.

\subsection*{Gamma-only storage: half the plane waves}

At $k=0$ a Kohn--Sham state can be chosen real, so $c(-G)=c^*(G)$ and one plane
wave of each $(G,-G)$ pair carries all of it. \code{K\_POINTS gamma} stores that
half, which halves $n_{pwx}$ and with it every array a band lives in --- the
wavefunctions, \code{vkb}, the Davidson subspace. On a 157-atom slab at
$E_{\mathrm{cut}}=60$~Ry that is \textbf{189~GB against 96}. It is a memory
feature, not a speed one: the two-bands-per-FFT packing of
\code{vloc\_psi\_gamma} is deliberately not implemented.

\textbf{Only the wavefunction sphere halves.} The dense $G$ set stays whole,
where \code{pw.x} halves both. The memory is entirely in the plane-wave-sized
arrays, so halving the dense set as well saves a fifth of a per cent and puts a
conjugate fill inside every consumer of a real field. The consequence to know is
that \code{ngm} here is not the one \code{pw.x} prints: the relation is
$n_{gm} = 2 n_{gm}^{\mathrm{QE}} - 1$, and the test suite asserts it.

Three things carry the trick: the field is rebuilt from both halves, every sum
over plane waves becomes $2\,\mathrm{Re}(\sum) - (G{=}0)$ --- that vector being
its own conjugate partner rather than half a pair --- and $\mathrm{Im}\,c(0)$
must stay zero. \textbf{The last is the silent one}: a complex $c(0)$ makes the
rebuilt field complex and the run converges to a plausible wrong answer.

\textbf{The validation needs no other code.} An explicit $k=0$ on the whole
sphere is the same calculation in the other storage, so the two agree to
round-off: totals to \textbf{1e-14}~Ry, eigenvalues to \textbf{3e-15} and forces
to \textbf{3e-16}~Ry/bohr, on LDA, PBE, LSDA and PBE+LSDA alike.

\begin{refuses}
\textbf{Refuses.} \code{K\_POINTS gamma} is \emph{substituted} by an explicit
$k=0$ on the full sphere --- the same physics at twice the storage, with a
warning --- for an ultrasoft or PAW dataset (\code{addusdens} and \code{newd}
need their own \code{fact = 2}), for a run that uses symmetry (a rotation
carries a stored $G$ onto an unstored $-G'$, so the permutation would need a
conjugation; \code{nosym = .true.} is the way out), for a spinor or spiral run,
and for a \textbf{Hubbard $U$}. The last is the only one of them that would go
wrong \emph{during} the run rather than after it: $n_{\sigma}$ is built from a
plain $\langle \phi | \psi \rangle$ over the stored list, so the Hubbard
potential and energy would follow occupations about a quarter of their true
size and the SCF would converge silently to a different ground state.
\code{estimate()} makes the same decision the driver makes, so it sizes
the run either way.

\textbf{Four things after the run refuse rather than substitute}, because by
then the states exist and the honest answer is to say the sum is not written:
the projected density of states, the sum-over-states $\chi_0$ (so an optical
spectrum of a molecule in a box, which is the archetypal gamma case), an STM
image and the vertical tunnelling transmission. Each is a plane-wave sum that
would have to be $2\,\mathrm{Re}(\sum) - (G{=}0)$ and is not; the escape is the
same explicit $k=0$, which is an \emph{exact} substitution rather than an
approximate one.

\textbf{The dynamical matrix works under gamma storage}, including a partial
one --- the combination a molecule on a fixed slab needs. The Sternheimer solve
is gamma-aware throughout (the projector, the CG's products, the
preconditioner, the response density), and a displacement's source term is one
$\mathrm{jvp}$ through \code{at\_positions}, so it inherits the corrected
\code{h\_psi} with no branch of its own. Against the whole sphere: frequencies
to \textbf{9.6e-10}~cm$^{-1}$ and the matrix to \textbf{1.5e-13}~Ry/bohr$^2$,
with the acoustic residue identical; a partial matrix to \textbf{5.9e-11}. And
$\chi_0$ against a central difference of the density --- which shares no
machinery with either storage --- to \textbf{7.0e-6}.

\textbf{A potential-only meta-GGA works under gamma storage and did not until
2026-09-18.} $\tau$ is built from $i(\mathbf k+\mathbf G)c_{\mathbf G}$, which is
\emph{odd} in $\mathbf G$ where $c_{\mathbf G}$ is Hermitian, so the density's
conjugate fill is the wrong fill and the corrected form is not
$2\,\mathrm{Re}(\sum)-(G{=}0)$ either. Pairing each $\mathbf G$ with $-\mathbf G$
gives $\nabla\psi = 2\,\mathrm{Re}\,h$ with $h=\sum_{\mathbf G\in\mathrm{half}}
i\mathbf G\,c_{\mathbf G}e^{i\mathbf G\cdot\mathbf r}$, the $G=0$ term dropping out
on its own because it carries a factor $\mathbf G$, so the half-sphere $\tau$ is
$4(\mathrm{Re}\,h)^2$ where the whole-sphere one is $|\nabla\psi|^2$. Taking
$|h|^2$ there, which is what this did, \textbf{loses the oscillation pointwise}
and is a clean factor of two only in the \emph{integral}: on two-atom silicon
under \code{tb09} the integrals were in the ratio \textbf{0.506}, the pointwise
disagreement was \textbf{7.7e-2} on a $\tau$ whose maximum is 1.2e-1, and the two
total energies were \textbf{16.9~mRy} apart --- in a functional where $\tau$ enters
$v_x$ directly rather than through an energy. Corrected, the two storages agree to
\textbf{1.3e-15} pointwise with the totals identical to every printed digit.

\textbf{The velocity operator refuses}, and it is the same oddness one step
further out. $\partial\hat H/\partial\mathbf k$ is odd in $\mathbf G$, so the
full-sphere sum is $2i\,\mathrm{Im}(\sum_{\mathrm{half}})$ and the diagonal of a
real state is \emph{exactly zero} --- time reversal saying a real state at
$\Gamma$ does not move. The half-sum returns an $O(1)$ number instead:
\textbf{0.863}~Ry\,bohr against \textbf{1.8e-15} at an explicit $k=0$, on the same
cell. That is the whole error rather than half of it, so unlike a density it is
not recoverable by the usual rule, which is why this refuses.
\code{Calculation.at\_kcart} is the one place it is said, because every consumer
--- the optical conductivity, the Kerr angle, second-harmonic generation, the
shift current, the nesting function and \code{band\_velocities} --- reaches the
operator through \code{VelocityOperator}, its only caller. A \textbf{topological
invariant is unaffected}: \code{at\_kpoints} rebuilds the sphere whole, so a
Chern number or a $Z_2$ index on a gamma ground state runs, and agrees with the
whole-sphere curvature to \textbf{7.4e-12} on a curvature of 2.2e-4.

A \textbf{field} response is still refused, and the distinction is the reason: a
field enters through the commutator $[H,\mathbf r]$, which builds plane-wave
sums of its own. The dielectric tensor, the Raman tensor, the strain response
and the Born charges all stand on it. It does not fail --- with the guard
lifted silicon's dielectric constant is $501.7/213.1/253.1$ against an isotropic
$190.8$ --- which is why it is a refusal. A checkpoint refuses rather than drops a
\emph{driven} magnetic field (a fixed-spin-moment scheme's, which no input file
determines); an applied field is saved normally, and so is a DFT+U state, whose
manifold is rebuilt from the \code{HUBBARD} card on load. \code{stress}, \code{solver} and
\code{history} are deliberately not stored, so what comes back is a state and
not a report.
\end{refuses}

% =====================================================================
\section{Ground state}

The Kohn--Sham problem, generalised because ultrasoft and PAW make the overlap
non-trivial:
\begin{align*}
\hat H\,|\psi_{n\mathbf k}\rangle &= \epsilon_{n\mathbf k}\,\hat S\,|\psi_{n\mathbf k}\rangle,
& \hat S &= 1 + \sum_{I,ij} q_{ij}\,|\beta_i^I\rangle\langle\beta_j^I| \\[2pt]
\hat H &= -\nabla^2 + v_{\mathrm{loc}}(\mathbf r) + v_{\mathrm H}[n] + v_{xc}[n]
  &&+ \sum_{I,ij} D_{ij}^I\,|\beta_i^I\rangle\langle\beta_j^I|
\end{align*}
with
$n(\mathbf r) = \sum_{n\mathbf k} w_{\mathbf k} f_{n\mathbf k}
\big(|\psi_{n\mathbf k}|^2 + \sum_{ij} Q_{ij}^I \langle\psi|\beta_i\rangle\langle\beta_j|\psi\rangle\big)$,
the augmentation term being what ultrasoft adds. $D_{ij}$ is rebuilt from the
potential every iteration, which is why the SCF is a fixed point in $(n,D)$ and
not in $n$ alone. Convergence is measured as QE measures it --- the Hartree
energy of the residual:
\begin{equation*}
dr^2 = \iint \frac{\delta n(\mathbf r)\,\delta n(\mathbf r')}{|\mathbf r-\mathbf r'|}
  \,d\mathbf r\,d\mathbf r' < \code{conv\_thr}
\end{equation*}

\begin{pycode}
result = run_scf(system, pseudos, conv_thr=1.0e-10, mixing_mode="anderson")
print(result.total_energy, result.energy_terms)   # QE's own term-by-term split
print(result.fermi_energy, result.homo, result.lumo)
\end{pycode}

\textbf{Metals} work with every smearing QE has (\code{gaussian},
\code{marzari-vanderbilt}, \code{methfessel-paxton}, \code{fermi-dirac}) and
with the tetrahedron family --- both linear and Bloechl-corrected --- usable as
an occupation scheme \emph{inside} the SCF, not only for a density of states.

\textbf{Continuation across a change of spin regime.} Give \code{run\_scf} a
previous result as \code{starting\_from} --- \code{System.with\_spin}
underneath --- and it promotes a converged state into a
target's variables: non-magnetic into collinear, collinear into noncollinear,
spin-orbit switched on. It reaches the \emph{same} solution as a fresh run
($\le$4e-8 Ry on six cases) in as little as one iteration. The magnetization is
\textbf{seeded} from the target's \code{starting\_magnetization} when the source
has none, because nothing in the SCF breaks spin symmetry on its own and an
unseeded promotion would converge back to the unpolarized solution and report
success.

A \code{HUBBARD} card comes along, and the occupation matrix promotes the way
the density does: it is decomposed into a charge and a Pauli vector, the vector
is laid along the \emph{target's} own axis --- the \code{angle1}/\code{angle2}
direction or the \code{STARTING\_MOMENTS} card, exactly what the charge and
\code{becsum} are turned onto --- and the four spin pairs are rebuilt from the
two. That closes the staged route into a hard magnet --- converge the collinear
antiferromagnet, promote it, let the moments cant --- and, because a
spinor-to-spinor promotion is the same code path, the \textbf{checkpoint
resume} of a noncollinear DFT+U run. Coming back the other way, the moment is
projected onto the axis the density is collinear along and the part of it that
lies off that axis is \textbf{refused} above $10^{-8}$ rather than dropped: the
diagonal of a canted occupation matrix is a different state, not a coarser one.

\textbf{The shell and the charge have to cross onto the same axis, and until
2026-09-20 only the charge did.} The pair was written into the two diagonal spin
blocks whatever was asked for, which is a moment along $z$, so a run whose
density crossed onto $x$ started with its correlated shell 90 degrees away from
it. On fcc nickel at $U = 4.0$~eV the density arrived with 0.491~$\mu_B$ along
$x$ and the shell with 0.383 along $z$, worth 30.7~mRy of Hubbard splitting on
the wrong axis. Left free the shell turns, because the carried density's own
exchange field pulls it, and what the mismatch costs is iterations --- a global
spin rotation is free on a scalar-relativistic dataset, so the continuation must
cost the same whatever axis is asked for, and it took 2, 7 and 8 iterations at
\code{angle1} of 0, 45 and 90 against 2, 2 and 2 now. Under
\code{mixing\_fixed\_ns} it is not a cost: the matrix is held at its starting
value and the residual of that block is zero while it is held, so the run met
\code{conv\_thr} before the freeze was released and reported success with the
shell on $z$, the density on $x$, and a total 4.11~mRy above the right answer on
a cell whose anisotropy is exactly zero.

\code{magnetization} is how the moment crosses, on \code{run\_scf} and on
\code{Calculator.with\_spin} and \code{with\_moments}. \code{'auto'} is the
default and carries the source's moment where it has one, seeding the target's
\code{starting\_magnetization} where it does not; \code{'carry'} insists on
the first and raises rather than starting on a symmetric solution;
\code{'seed'} keeps the converged \emph{charge} and takes the moment from the
seed regardless, which is how a \textbf{different} magnetic state is reached
from the same charge; \code{'none'} starts unpolarized. \code{with\_moments}
defaults to \code{'seed'} where \code{with\_spin} defaults to \code{'auto'},
because a new texture is a new \code{STARTING\_MOMENTS} card and carrying a
noncollinear source's moment applies no rotation, so the card would never be
read.

\begin{pycode}
turning = 0.6 * np.array([[np.sin(a), 0.0, np.cos(a)]
                          for a in np.deg2rad([0, 90, 180, 270])])
cycloid = Calculator.from_file("h4-chain-ferro.in").with_moments(turning)
cycloid.get_scf()                      # 0.427 mu_B per site, a cycloid
afm = np.array([[0.0, 0.0, s] for s in (0.6, -0.6, 0.6, -0.6)])
flipped = cycloid.with_moments(afm).get_scf()   # the converged charge, a new texture
\end{pycode}

On that cell the default is worth 8.9 mRy. Carrying the cycloid's moment into
the antiferromagnetic card gives a converged run at $-3.792834$ Ry whose site
moments are $(0.402, 0, -0.402, 0)$, which is the cycloid's projection on $z$
with two sites dead and is not what was asked for; seeding gives
$-3.801738$ Ry at $(\pm 0.589)$ alternating, the antiferromagnet, in one
iteration fewer. Neither run says anything: both report
\code{converged = True}.

\textbf{Under a Hubbard $U$ the occupation matrix follows the same decision},
and until 2026-09-18 it did not: $n_\sigma$ was the one member of the mixed
triple crossing a continuation unchanged whatever \code{magnetization} said,
while the density and \code{becsum} were reseeded or zeroed beside it. A DFT+U
cell converged ferromagnetic and then continued with
\code{magnetization='none'}, or handed a new texture by \code{with\_moments},
entered iteration~1 with a Hubbard potential split by order $U$ ---
$v = \alpha + \tfrac{U}{2}\delta - U n^\sigma$, so 0.3~Ry for a typical Ni
$U$ --- on a density that was exactly spin-degenerate or the new texture's, and
\code{mixing\_fixed\_ns} freezes it there for as many iterations as it is set
to. The run can relax back onto the source's magnetic configuration and report
its energy as the answer for the one that was asked for, which is the failure
\code{with\_moments}'s \code{'seed'} default exists to prevent for the
density. Now: \code{'carry'} takes the matrix over, \code{'none'} takes its
spin splitting out and keeps its converged charge, and \code{'seed'} drops it
so that \code{init\_ns} builds it from the \emph{target's} own
\code{starting\_magnetization} --- which is what a fresh run with that
magnetization would do, and is better than any mapping of a density texture
onto per-site occupations available here. A checkpoint resume is untouched:
\code{magnetization} other than \code{'auto'} is refused outright on one.

\begin{refuses}
\textbf{A spin spiral is a change of frame and not only of regime.} A spiral
keeps its density in the \emph{rotated} frame, the one that turns with
$\mathbf q$, so its charge and $m_z$ mean the same thing as an ordinary
noncollinear run's and its transverse pair $(m_x, m_y)$ does not. A source that
has a transverse pair to carry, meaning any source with
\code{nspin\_mag = 4}, is therefore \textbf{refused} in both directions
across a spiral boundary. A collinear source is not: what crosses then is one
scalar field, laid along the target's own \code{angle1}/\code{angle2}. Spiral
to spiral is allowed at any $\mathbf q$, which is what a sweep needs.

\textbf{A moment on the spiral's own axis is refused as the ferromagnet it is,
and the refusal is at the door.} The spiral's spin rotation is about $z$, so a
magnetization along $z$ is invariant under it: a cone of zero opening angle, and
a stationary point at every $\mathbf q$. A run in that state converges, reports
a moment, and has computed the ferromagnet with $\mathbf q$ playing no part,
which is the \emph{nothing breaks the symmetry on its own} rule of the paragraph
above one axis further out. A spiral whose starting moments are all on the axis,
or all zero, is therefore refused when the calculation is built, so a run
started from scratch is caught as well as a continued one, and \code{angle1} is
named as the way through: 90 degrees puts the moment in the plane, which is the
proper screw. Including at $\mathbf q = 0$, where the collapsed state is the
right answer and refusing it looks strict --- a scan builds every other point
from that same system, so the one place its moments can be looked at is there.

\textbf{The direction a carried scalar is laid along comes from the atoms.} A
collinear source carries one number per point, and where it points in a spinor
target is that target's own texture: \code{System.local\_moments}, which is the
\code{STARTING\_MOMENTS} card where there is one and \code{angle1}/\code{angle2}
where there is not. Reading the angles alone said \emph{along $z$} for a target
whose texture is stated only in the card, since the card leaves both angles at
zero, and the card the run was built to honour was never read. A target whose
atoms point along different axes is refused instead, because one scalar cannot
follow a helix, with \code{magnetization='seed'} named as the way through.

\textbf{The source density's own electron count is checked}, by integrating it,
rather than inferred from the atom list. What that catches is a \textbf{swapped
dataset}: a file with a semicore shell carries a different \code{z\_valence}
on the same atoms, and a continuation carries a density rather than
renormalising one.

\textbf{A magnetization held by a field or a constraint says so when it is
released.} A warning and not a refusal, because hold-converge-release is a
workflow: what must not happen is that the moment being carried is the field's
answer rather than the functional's without anyone noticing.

\textbf{\code{magnetization} is refused rather than ignored} where it has
nothing to act on: with no \code{starting\_from}, beside a
\code{ContinuedState} that has already resolved the question, and on a
\textbf{checkpoint resume}, where a resubmitted command line would otherwise
throw the converged moment away after every wall-clock kill.
\end{refuses}

\textbf{Convergence.} \code{mixing\_mode} selects \code{anderson} (default),
\code{linear}, \code{broyden}, Kerker/Thomas--Fermi preconditioning (\code{TF}),
\code{local-TF}, or \code{adaptive}; \code{scf\_solver} offers a Newton--Krylov
alternative that is slower but reaches solutions the mixer cannot hold.

\code{TF} screens the residual with one Thomas--Fermi length for the whole cell.
\code{local-TF} makes that length a function of $\rho(\mathbf r)$, which is what
a \textbf{slab} needs: the metal wants strong screening and the vacuum wants
none, and a single compromise value over-screens one and under-screens the other
--- charge sloshing between the surfaces. The screened residual is the solution
of
\begin{equation*}
  4\pi e^2 v(\mathbf G) + |\mathbf G|^2 (\alpha v)(\mathbf G)
    = |\mathbf G|^2 (\alpha\,\delta n)(\mathbf G),
  \qquad \alpha(\mathbf r) = 3\left(\tfrac{2\pi}{3}\right)^{5/3} r_s(\mathbf r),
\end{equation*}
with $(\alpha f)$ multiplied in \emph{real} space, so the operator is not
diagonal in $\mathbf G$ and is inverted iteratively at each step. On a cobalt
film it converges at \code{mixing\_beta = 0.7} where plain Anderson at the same
$\beta$ diverges and Kerker at $0.3$ is still far out after 250 iterations.

\code{adaptive} is Elk's \code{mixadapt} and answers the opposite failure. Every
other mode here damps an iteration that is \emph{overshooting}; this one lengthens
one that is crawling. It carries a separate step length for each component $j$ of
the mixed vector and moves it by the sign of that component's residual,
\begin{equation*}
  \beta_j \leftarrow \min(\beta_j + \beta_0,\ \beta_{\max})
  \quad\text{if } r_j \text{ kept its sign,}
  \qquad
  \beta_j \leftarrow \tfrac{1}{2}(\beta_j + \beta_0) \quad\text{if it changed,}
\end{equation*}
and then steps as $\beta_j\,\nu_j + (1-\beta_j)\,\mu_j$. A component pushed the
same way for many iterations is one the SCF is crawling along, so its step grows
until it moves; one that overshoots and comes back has its step cut. That makes it
a stall detector with no threshold in it, which is why it is worth reaching for on
a \textbf{magnet}: turning every moment in a cell together costs no energy without
spin--orbit coupling, so the residual has no component along that rotation and the
flat direction has to be crossed rather than descended, and twisting the moments
slowly costs only the spin-wave energy $D q^2$, which vanishes as the cell grows.
The whole packed vector adapts independently, so \code{becsum} and the Hubbard
\code{ns} get their own step lengths rather than the density's.

\textbf{\code{mixing\_beta} does not mean the same thing under \code{adaptive}},
and this is the one thing to know before using it. It is Elk's \code{beta0}: the
increment, the starting value and the floor the halving relaxes towards, all
three, not a step length. Elk's default is $0.05$ and it is small on purpose;
$0.7$ saturates at $\beta_{\max}$ within two iterations and leaves plain linear
mixing at $0.85$ under an adaptive name, which is what the warning above that
value is for. Each mode supplies its own default when \code{mixing\_beta} is
unset, $0.7$ for the QE family and $0.05$ here, so an input that names neither
gets the right one; an input that \emph{does} set \code{mixing\_beta} has that
number adopted as $\beta_0$, and the same figure therefore means two different
things depending on the mode.

\textbf{Elk's $0.05$ is a safe default rather than the best value}, and the
measurement says so: on bcc iron with its moment in the plane
(\code{tests/data/qe/fe-noncolin-pbe-stress.in}, ultrasoft, PBE, noncollinear)
$\beta_0 = 0.2$ converges in \textbf{16} iterations where $\beta_0 = 0.05$ takes
24, because the soft directions want the step to grow faster than Elk's default
lets it. Anderson on the same cell takes 43 at the input's own
\code{mixing\_beta} and 24 at the best one for it, and \code{pw.x} takes 19; all
of them reach the same energy within $3\times10^{-10}$ Ry and the same moment
\emph{length} within $3\times10^{-5}\,\mu_B$, with the transverse components
that symmetry makes zero coming out below $2.4\times10^{-4}\,\mu_B$. So raising $\beta_0$ is the knob to try on a
magnet that will not move, and it is the opposite of the reflex a slow SCF
usually gets. The saving is smaller than the iteration count suggests, because
the diagonalisation threshold is scheduled from the density error and a mixer
that drops it faster is handed a tighter eigenproblem: the Davidson work falls
by $1.6\times$ where the iteration count falls by $2.7\times$.

On an easy cell it loses, which is what a robustness mixer does: two-atom
silicon (\code{benchmarks/si-1k.in}) at \code{conv\_thr = 1e-10} takes 17
iterations against Anderson's 7, to the same total energy to ten digits.

\begin{refuses}
\textbf{A mixer can change which solution the run finds}, and this is the thing
to know before trusting an \code{adaptive} result on a magnet. On a four-cell
hydrogen ultracell under a field turning $90^\circ$ per cell --- the case where
rotating every moment together costs nothing --- that cell has \textbf{three}
self-consistent solutions, and every one of them reports
\code{converged = True} with an \code{accuracy} three orders below
\code{conv\_thr}. Anderson at $\beta = 0.1$ and $0.7$ finds one of them, agreeing
to eight digits of the band energy; Anderson at $\beta = 0.3$ finds a second;
\code{adaptive} at both its values finds the third, the one whose moments follow
the applied field's ladder. \textbf{Which is the ground state is not something
the run reports}: the moments differ by tens of degrees per cell, and nothing in
a converged result orders them.

\textbf{The step length does this too, which is the part not to file under
``adaptive''.} Anderson's own $\beta$ moves between solutions here, so this is
not a property of one mixer and running a second mixer is not a test that the
basin was kept. What a growing step buys is more of the same freedom, not a new
kind of it.

So on a cell with a soft or degenerate magnetic direction, vary
\textbf{both} the mixer and $\beta$, and compare the \textbf{moments} rather
than the iteration count: a converged run is not evidence that the basin was
kept, and ``fewer iterations'' is satisfied by a run that found something else.
\code{SCFResult.site\_moments} and \code{UltracellResult.cell\_moments()} are
what to compare.

\code{mixing\_mode = 'adaptive'} also refuses a \textbf{preconditioner} rather
than ignoring one. Its step is pointwise in real space with its own factor per point,
and Kerker or \code{local-TF} is a single scalar acting in $\mathbf G$-space; the
two do not compose, so a run asking for both raises. Choose the one that matches
the direction the cell is slow in: Kerker for long-wavelength charge, this for a
residual that is small and will not turn around.

It is registered under its own name and is \textbf{not} aliased to any
\code{pw.x} mode. QE has no adaptive mixer, so a \code{pw.x} input reaching one
of QE's names here still gets what QE would do.
\end{refuses}

\code{max\_iterations}
defaults to 100 --- \textbf{a hard SCF may need more}, and hitting the cap is
reported as not converged rather than as an answer: \code{run\_scf} issues a
\code{RuntimeWarning} naming the accuracy it reached, and \code{Calculator}
raises. Both matter, because a script driving the functional entry point --- a
spiral scan, a $\lambda$ ramp, a moment-versus-field sweep --- is exactly what
would otherwise read \code{total\_energy} off an unconverged density without
looking at \code{converged}.

\code{mixing\_ndim} sets how many previous iterations the extrapolation spans,
default 8 as in \code{pw.x}. Raising it is the standard first move on a cell
that will not converge, and \emph{on the magnetic benchmark here it does not
help}: \code{benchmarks/fe-mag-1k.in} at \code{conv\_thr = 1e-8} takes 27
iterations at 4, \textbf{25 at 8}, 33 at 12 and 39 at 20, so the default is
already the best value and the deeper histories are worse. The nonmagnetic twin
of the same cell (\code{fe-unstable-nonmagnetic.in}, same $\beta$) has the same
shape --- 26, \textbf{21}, 30, 26 --- which says the sensitivity is not about
magnetism.

\textbf{\code{dr2} is reported as its two halves} as well as its total, in
\code{scf.history} and on the console. \code{rho\_ddot} weights the charge
residual by $1/G^2$ and the magnetization residual by a constant --- a factor of
13.6 apart at $G_{\min}$ on \code{fe-mag-1k} --- so a single \code{accuracy}
below \code{conv\_thr} bounds the moment much more weakly than it bounds the
charge, and only the split can say whether it is the charge or the magnetization
that is still moving.

% =====================================================================
\section{Bands, densities of states, projections}

A band structure is one diagonalisation per $\mathbf k$ on a fixed density ---
\code{run\_bands}, or \code{run\_nscf} for a non-self-consistent run on a grid
rather than a path. A density of states is
\begin{equation*}
g(E) = \sum_{n\mathbf k} w_{\mathbf k}\,\delta(E - \epsilon_{n\mathbf k})
\end{equation*}
with $\delta$ replaced by a smearing function or evaluated by tetrahedron
interpolation. The projected DOS weights each term by
$\big|\langle\phi_\mu|\hat S|\psi_{n\mathbf k}\rangle\big|^2$, and the spilling
parameter is what is \emph{not} captured,
$\sum_{n\mathbf k} w_{\mathbf k} f_{n\mathbf k}\big(1-\sum_\mu|\langle\phi_\mu|\hat S|\psi\rangle|^2\big)$.

\begin{pycode}
from defumat.workflows.bands import run_bands
from defumat.workflows.dos import run_dos
from defumat.workflows.pdos import run_pdos

bands = run_bands(system, pseudos, scf.density, kpoints=path)   # a k-path

dos, _ = run_dos(system, pseudos, scf.density, grid=(12, 12, 12),
                 scheme="tetrahedra_opt")     # -> (DensityOfStates, NSCFResult)

pdos, _ = run_pdos(system, pseudos, scf)      # -> (ProjectedDOS, NSCFResult)
print(pdos.channels[0])                       # resolved by atom, l and m
print(pdos.charges.charges, pdos.charges.spilling)
\end{pycode}

The DOS schemes live behind a name registry (\code{DOS\_SCHEMES}), and the
projected DOS feeds \emph{the same} registry as a per-band weight, so a PDOS and
a DOS are the same integration. Projections are
$\langle\phi|\hat S|\psi\rangle$ on Löwdin-orthogonalised pseudo-atomic
orbitals. Validated against a purpose-built \code{projwfc.x} on seven cases:
\textbf{6.9e-4} on a projection, \textbf{4.7e-5} on a Löwdin charge, and the
spilling parameter to all four decimals it prints.

\subsection*{Resolved by $j$ and $m_j$ for a spin-orbit run}

With \code{noncolin} and \code{lspinorb} the orbitals projected onto are the
\textbf{spin-angle functions} $|l\,j\,m_j\rangle$ rather than $|l\,m\rangle$
times a spin, so a channel is labelled by $j$ and $m_j$ and the projection
separates the two branches spin--orbit coupling splits a shell into. For
platinum that is $5d_{3/2}$ against $5d_{5/2}$: one number, resolved, where an
$l$-resolved projection reports their sum and cannot see the splitting at all.
The functions are
\begin{equation*}
  |l\,j\,m_j\rangle \;=\;
  \sum_{\sigma=\uparrow,\downarrow}
  C\!\left(l,\tfrac12,j;\,m_j-\sigma,\sigma\right)\,
  Y_{l,\,m_j-\sigma}\;\chi_\sigma ,
\end{equation*}
Clebsch--Gordan coefficients against the harmonics, and the number of columns is
$\sum_\mathrm{shells}(2j+1)$ --- which is \emph{not} in general twice the scalar
count, because it is read off the dataset's own $j$ channels.

\begin{pycode}
from defumat import Calculator

calc = Calculator.from_file("pt-soc-paw-nosym.in",   # noncolin + lspinorb
                            pseudo_dir="../pseudo")
pdos = calc.get_pdos()

print(pdos.channels[0])          # -> #1 Pt1 s_j0.5 m_j=-0.5
print(pdos.charges.format())     # -> s = 0.5032, p = 0.9944, d = 8.4903
print(pdos.charges.charges_lm)   # -> None: a spin-angle function has no m
\end{pycode}

Validated on fully-relativistic \textbf{ultrasoft and PAW} platinum datasets
against \code{projwfc.x} runs generated for them, four ways. The \emph{columns} --- how
many, and each one's $(l, j, m_j)$ --- match QE's own label table exactly. The
two datasets differ there in a way worth seeing: the PAW file carries both $j$
partners of every $l$, so its 18 columns are exactly twice the $j$-averaged 9,
while the ultrasoft one has 12 against a scalar 11 --- the count is
$\sum (2j + 1)$ read off the file and not a doubling. The \emph{projections}
agree to
\textbf{6.3e-6}, compared as sums over each degenerate multiplet: a spinor band
structure is Kramers-degenerate everywhere, and an individual
$|\langle\phi|\hat S|\psi_n\rangle|^2$ is not invariant under the mixing inside
a multiplet that either code's eigensolver is free to choose, where the sum is.
The \textbf{Löwdin charges and the spilling parameter agree to every digit
\code{projwfc.x} prints}. And the \emph{curves} --- the total and each shell's,
on \code{projwfc.x}'s own energy grid and broadening --- agree to \textbf{0.16
per cent of the peak} everywhere below the lowest empty band, above which the
topmost states are the ones neither eigensolver converges. Timing, one core each from a converged ground state:
\code{get\_pdos} costs \textbf{1.65~s} against \code{projwfc.x}'s 0.62~s, and
on both sides that is mostly rebuilding from the pseudopotential files rather
than projecting --- the projection and its integration are 0.137~s of it here.

Norm-conserving is untested rather than refused --- no reference case is
committed for it.

\subsection*{Two spin channels for a noncollinear run without spin--orbit coupling}

With \code{noncolin} but \code{lspinorb = .false.} there is no $j$ to resolve
by: nothing couples $s_z$ to the orbital motion, so $s_z$ is still a good
quantum number and the orbitals are $|l\,m\rangle\chi_\sigma$ --- an up and a
down copy of every harmonic, $2\,\mathrm{natomwfc}$ columns where a collinear
run has $\mathrm{natomwfc}$. What comes back is nonetheless \emph{not} twice as
many channels: the columns are routed into an \textbf{up and a down density of
states}, which is \code{partialdos\_nc}'s \code{nspin0 = 2} and is the shape an
LSDA projection has. This is the regime of any magnetic run with
scalar-relativistic datasets, and it is what \code{h10-chain-noncolin.in} and
\code{alas-magnetoelectric-nosoc.in} are.

\begin{pycode}
from defumat import Calculator

calc = Calculator.from_file("h-atom-noncolin.in",   # noncolin, no lspinorb
                            pseudo_dir="../pseudo")
pdos = calc.get_pdos(symmetrize=False)

print(pdos.nspin)                # -> 2: an up channel and a down one
print(pdos.channels[0])          # -> #1 H1 s: no s_z, the spin is an axis now
print(pdos.charges.polarization) # -> [0.99857905]: the moment the s shell carries
\end{pycode}

The decomposition is along the \textbf{global} $z$ and not along the local
moment, exactly as QE's is: an in-plane texture reports two equal channels,
which is a statement about the frame the orbitals are built in rather than a
defect.

\textbf{No \code{projwfc.x} reference for this regime is committed here}, so
what validates it is an identity instead, and a sharper one: without spin--orbit
coupling a state polarized along $+z$ block-diagonalises the noncollinear
Hamiltonian into the two collinear ones, so the projection of such a run must
reproduce an LSDA run of the same cell channel for channel --- and the LSDA
route is itself validated against \code{projwfc.x}. On a hydrogen atom in a 12~bohr box the two
runs agree to \textbf{3e-12~Ry} in total energy, their majority curves to
\textbf{2e-5 of the peak} and their minority ones to \textbf{0.4 per cent}, and
their Löwdin charges to \textbf{5e-3} on a polarization of \textbf{0.9986}
electrons --- which is the number a code that binned both spins into one channel,
or split the weight evenly, would get wrong while passing every sum rule. The two
bounds differ for an arithmetic reason rather than a chosen tolerance: the
noncollinear branch reaches $\rho_\downarrow$ as $(n - |m|)/2$, a cancellation of
two numbers of order $0.1$ where the collinear branch carries $\rho_\downarrow$
itself, and that is worth $1.1\times10^{-4}$~Ry on the \emph{empty} minority
eigenvalue where every occupied one agrees to $10^{-6}$.

Unlike a spin--orbit run this case \emph{does} have an $m$ to report, so
\code{charges\_lm} is filled where a $j$-resolved projection leaves it
\code{None}. \code{print\_lowdin} allocates it only \code{IF (nspin /= 4)} and
so prints none for either noncollinear branch; the reason it gives --- a
spin-angle function has no $m$ --- is true of the spin--orbit branch alone, and
the table here is the richer one on purpose.

\textbf{A reduced $\mathbf k$-set is allowed here}, and what makes it so is that
the group average runs over the right operator. A point-group operation turns
two things at once in this regime --- the real harmonics among themselves and
the spin frame the spinor is written in --- so the matrix averaged over is the
tensor product $D^l \otimes \overline{U}$, with $U$ the $SU(2)$ rotation that
carries $\mathrm i \sigma_y D^{*}$ for an operation that is a symmetry only
together with time reversal. Averaging the $m$ indices alone leaves the spin
frame where it was.

\begin{pycode}
from defumat import Calculator

# fcc nickel, noncollinear, moment along (1,1,1): twelve operations, six of
# them with time reversal, and six of the twelve mix the two spin components.
calc = Calculator.from_file("ni-noncol-111.in", pseudo_dir="../pseudo")
pdos = calc.get_pdos()                   # symmetrize = True is the default
# (nspin, nat, l, m): the majority d weights of the one atom, per m.
print(pdos.charges.charges_lm[0, 0, 2])
# -> [0.94676054 0.9178314  0.9178314  0.94676054 0.91783139]
\end{pycode}

The number: against the same cell run at \code{nosym} on the closed
$4\times4\times4$ grid, the wedge's Löwdin charges agree to
\textbf{1.1e-5 electrons} per column with the symmetrisation on and to
\textbf{5.4e-2} without it. On the four-hydrogen helix
(\code{h4-cycloid-90.in}), where the shell is $s$ and the spin factor is
therefore the \emph{whole} symmetrisation, the same pair reads \textbf{4.1e-6}
against \textbf{3.8e-4}.

\emph{The conjugation is measured rather than assumed, and the obvious cell
cannot measure it.} Nickel with its moment along $z$ --- which is what
\code{ni-ldau-noncol.in} is --- has every operation turning the spin about $z$
alone, so its $2\times2$ matrix is diagonal, its phase drops out of a modulus,
and $U$, $U^{T}$, $U^{\dagger}$ and the identity all give the same answer to
the last digit. Turning the moment onto a three-fold axis separates them:
$U$, $U^{T}$ and $U^{\dagger}$ land 3.9e-2, 2.1e-2 and 2.1e-2 from the closed
grid, which is worse than not symmetrising the spin at all, and only
$\overline{U}$ reaches 1.1e-5.

\begin{refuses}
\textbf{Refuses.} A \textbf{symmetrised} projection of a \code{lspinorb}
run --- run with \code{nosym} and the whole $\mathbf k$-grid, where both codes
use every point and neither averages, or pass \code{symmetrize=False}. Its
columns are $|j\, m_j\rangle$ and the operator \code{sym\_proj\_so} averages
over is $D^j$ for $j = \tfrac12 \ldots \tfrac72$, which is a different matrix
from the one the branch above uses and is not reachable by relabelling it.
\emph{The noncollinear branch without spin--orbit coupling is no longer
refused}: its columns are $|l m\rangle \otimes |\sigma\rangle$, the operator is
the tensor product $D^l \otimes \overline{U}$ of the harmonic rotation and the
conjugated $SU(2)$ one, and it is built. A \textbf{fully-relativistic dataset with
\code{lspinorb = .false.}}: QE dispatches the orbitals on \code{has\_so} and
their labels on \code{lspinorb}, so it builds $j$-resolved columns and calls
them up/down ones. An $m$-resolved charge is \emph{absent} rather than refused
for a \textbf{spin--orbit} run: \code{charges\_lm} is \code{None} there, since a
spin-angle function has an $m_j$ and no $m$, and there is no $p_x$ weight to
report.
\end{refuses}

% =====================================================================
\section{Pseudopotentials and functionals}

\textbf{Kinds:} norm-conserving, \textbf{ultrasoft} and \textbf{PAW} --- the
two-grid split, the augmentation charge, the overlap operator, self-consistent
$D_{ij}$, and PAW's one-centre terms including its radial Poisson solve and
spherical quadrature. UPF v2 is read. \textbf{The kinds mix within one cell},
which is an ordinary \code{pw.x} run and is now an ordinary one here: a
norm-conserving species in a PAW cell has projectors and therefore rows of the
$(n_{kb}, n_{kb})$ coefficient matrix, and it has no augmentation charge and no
one-centre correction, so those rows stay zero and its bare $D_{ij}$ arrives
through the same route it always did. Displaced silicon with
\code{Si.pz-n-kjpaw\_psl.0.1.UPF} on one site and \code{Si.pz-vbc.UPF} on the
other matches \code{pw.x} to \textbf{4.7e-9}~Ry and \textbf{1.15e-5}~Ry/bohr on
a force of 0.0758 --- between the same cell run all-PAW (3.9e-7) and
all-norm-conserving (2.0e-5), so each species keeps its own accuracy.

\textbf{Functionals:} LDA (PZ), \textbf{PBE}, \textbf{revPBE}, \textbf{PBEsol},
on the plane-wave grid and on the PAW spheres. The functional comes from the
pseudopotential headers unless \code{input\_dft} overrides it. Only the
\emph{energy} is written down; $v_{xc}$, and a GGA's $v_1$ and $v_2$, come from
\code{jax.grad} of it.

\textbf{Meta-GGA, potential-only branch:} \code{input\_dft = 'tb09'}
(Tran--Blaha) and \code{'bj06'} (Becke--Johnson). These invert the rule above
--- they \emph{are} potentials, there is no energy --- so the consequences are
enforced rather than documented: \code{run\_scf} warns that its total is not the
value of any functional it minimised, and forces, stress, phonons and response
all refuse. Silicon's gap goes from LDA's 0.49~eV to \textbf{1.13~eV} against an
experimental 1.17.

\textbf{Van der Waals:} Grimme \textbf{D2} (\code{vdw\_corr = 'grimme-d2'}).
Bilayer graphene matches \code{pw.x} to \textbf{3.1e-9 Ry} and binds at
3.23~\AA{} where PBE alone has no minimum.

\begin{refuses}
\textbf{Refuses.} UPF~v1; an unimplemented functional (rather than silently
substituting LDA); energy-carrying meta-GGAs (TPSS, SCAN, M06L), whose potential
has a $\partial E/\partial\tau$ piece acting on the wavefunction that is not
written; plain \emph{ultrasoft} under a meta-GGA, which has no partial waves to
reconstruct $\tau$ from inside the sphere where PAW does; EXX; and
\textbf{D3, Tkatchenko--Scheffler, MBD and XDM} --- where \code{pw.x} warns and
silently runs with no correction at all.
\end{refuses}

% =====================================================================
\section{Spin, magnetism and spin-orbit}

\begin{center}
\small
\begin{tabular}{@{}lll@{}}
\toprule
\textbf{regime} & \textbf{how to ask} & \textbf{note} \\
\midrule
unpolarised & default & \\
collinear & \code{nspin = 2} & one Fermi level, or two under \code{tot\_magnetization} \\
noncollinear & \code{noncolin = .true.} & magnetization is a vector \\
spin-orbit & \code{+ lspinorb = .true.} & $j$-resolved projectors; NC, US and PAW \\
\bottomrule
\end{tabular}
\end{center}

\textbf{Three spin numbers are kept apart}, and \code{System} exposes all three:
\code{nspin} says which regime is in force, \code{npol} is the number of spinor
components of a \emph{wavefunction}, and \code{nspin\_mag} the number of
components of a \emph{density}. They coincide at 1 and 2 and come apart at 4,
where \code{npol = 2} but \code{nspin\_mag} is \textbf{one} for a nonmagnetic
spin-orbit run --- which is why such a run costs about what a doubled
unpolarized one costs.

\textbf{Fixed occupations under \code{nspin = 2}} fill each channel to its own
count rather than searching for a level: \code{tot\_magnetization} gives
\code{nelup} and \code{neldw}, and \code{iweights\_only} fills
$\mathrm{NINT}(n_\uparrow)$ bands in one channel and $\mathrm{NINT}(n_\downarrow)$
in the other. That is the \emph{only} shape the combination has ---
\code{pw.x} refuses fixed-occupation LSDA without a \code{tot\_magnetization},
and requires an integer one --- and this package makes the same two refusals at
input, before anything is allocated. Each channel's highest occupied level is
reported as \code{fermi\_energy\_up} / \code{fermi\_energy\_down} beside the
overall \code{homo} and \code{lumo}, which is what \code{iweights} returns.
Validated on an oxygen atom at \code{tot\_magnetization = 2}, whose four-up /
two-down filling is the Hund's-rule ground state: \textbf{4.9e-9 Ry} against
\code{pw.x}.

\textbf{Magnetic fields and constraints:} a uniform field over the cell, a field
inside one atom's sphere (through a \code{LOCAL\_MAGNETIC\_FIELDS} card that
\code{pw.x} has no counterpart for), Elk's \code{reducebf}, and all four of QE's
\code{constrained\_magnetization} schemes. \textbf{The field's energy is not part
of the reported total} --- QE prints \code{etcon} and never adds it, and Elk
excludes its external field's energy by the same convention.

\emph{A seed field has no natural size, so the symmetry filter has no scale in
it.} An infinitesimal field is the standard way to start a texture: small enough
not to bias the energy, large enough to pick the state. What the magnetic filter
asks of a field is a \emph{pattern} --- does the operation map it onto itself,
up to time reversal --- so each field is compared after being divided by its own
largest component, and it is ignored only below $10^{-12}$~Ry, which is the same
floor that decides whether the run is magnetic at all. One rule: a field big
enough to make a run magnetic is big enough to cut its group. (It used to be two,
seven orders of magnitude apart, and a field in between turned the run magnetic
and then contributed nothing to the filter --- which is the one thing worse than
no field, because the full crystal group then averages away the texture the field
was applied to create.)

\emph{\code{reducebf} decays the field after the convergence test}, so a run can
stop with the field still substantially on. The state reported is then the ground
state of a functional that includes a Zeeman term, and \textbf{the error in the
total is second order in the residual field}, not the field energy --- the state
is stationary. On a hydrogen atom with a 0.1~Ry local field, \code{reducebf =
0.5} stops with $6\times10^{-6}$~Ry still applied and a total right to
$5\times10^{-14}$~Ry, while \code{reducebf = 0.99} stops after \emph{six}
iterations with $9.5\times10^{-2}$~Ry applied and a total $4.7\times10^{-5}$~Ry
out. A run that stops above $10^{-4}$~Ry now says so, naming
\code{field\_scale}, the residual and \code{field\_energy}; a field held at full
strength (\code{reducebf = 1}) is a deliberate calculation and is left alone.
\code{reducebf} is also held to Elk's own $[0.5, 1]$
(\code{readinput.f90:1264}), which this used to accept any value outside:
above 1 the symmetry-breaking field \emph{grows} every iteration, and below 0.5
it is gone before the density has responded to it.

\begin{refuses}
\textbf{Refused.} \code{constrained\_magnetization = 'fsm'} together with
\code{lspinorb = .true.} The fixed-spin-moment search measures $\chi = \delta
m_a / \delta B_a$ one cartesian component at a time and inverts it the same
way, which models $\partial m_a/\partial B_b$ as \emph{diagonal} --- three
independent one-dimensional searches. Spin-orbit coupling ties the moment to the
lattice and is exactly what makes that untrue: pushing along $x$ moves the moment
along $z$ as well. The scheme would still converge sometimes and find the wrong
field the rest of the time, with nothing in the output to say which. Use a
penalty scheme (\code{'total'}, \code{'atomic'}) to hold a moment under
spin-orbit coupling.

\textbf{Also refused:} \code{noncolin = .true.} with \code{angle1}/\code{angle2}
set and no \code{starting\_magnetization}. The angles are the \emph{direction}
of a moment whose length is \code{starting\_magnetization}, so with no length
there is no moment: \code{domag} goes false, \code{nspin\_mag} collapses to 1,
and what runs is an unpolarized calculation with spinor wavefunctions ---
converged, successful, and not the calculation the angles asked for. \code{pw.x}
runs that input silently. A spinor run with nothing magnetic in it stays
accepted, because that is an ordinary calculation; the refusal is about a
\emph{stated} direction with nothing to point.

\textbf{The polar-angle constraint is guarded on the axis.}
\code{'total direction'} used to return a NaN potential in all three components
whenever the moment lay along $z$ --- which is where a run seeded from
\code{starting\_magnetization} with no angles starts. $\arccos'$ diverges at the
poles while the transverse moment's own derivative vanishes there, so the chain
is $0 \times \infty$; and at a \emph{round-off} transverse moment $m_z/|m|$
rounds to bit-exactly 1 and the clamp halves the tangent instead. The angle is
written $\operatorname{atan2}(|m_\perp|, m_z)$ with both arguments masked at the
value the derivative is taken at, which reproduces QE's \code{fact1} to 1e-12
off the axis and QE's zeroed transverse factors on it, plus its literal
$10^{-14}$ escape field along $x$ --- without which the moment sits at an
unstable stationary point of its own penalty and never turns
(\code{add\_bfield.f90:184-192}).
\end{refuses}

\textbf{Restarting an SCF from the middle} (\code{checkpoint\_dir},
\code{checkpoint\_every}, \code{max\_seconds}). A job killed at its wall clock
used to lose every iteration it had run: a state reached disk only after
\code{run\_scf} returned. It now writes on a cadence and \emph{resumes from the
same directory}, so a resubmitted batch script is the same command and
continues.

\begin{lstlisting}[language=Python]
from defumat import Calculator
calculator = Calculator.from_file("scf.in", checkpoint_dir="ckpt")
scf = calculator.get_scf()          # run again after a kill: it continues
\end{lstlisting}

\emph{A restart is three things}, and the third is the one with no obvious home:
the state (density, wavefunctions, \code{becsum}, \code{ns}, \code{tau}); the
\emph{mixer's history}, without which Anderson restarts empty and pays back the
iterations the checkpoint saved; and the \emph{loop state} --- the iteration
number, \code{dr2} and the diagonalisation threshold \code{ethr}. That last one
matters because QE's threshold schedule is indexed on the iteration number, so a
resume re-entering at iteration 1 with a fresh threshold converges on a different
schedule than the one it left. With all three carried the restart is exact: on
silicon at \code{conv\_thr = 1e-12}, stopping and resuming costs the \emph{same}
total iteration count as an uninterrupted run at three different mixing
parameters (17/17, 13/13, 11/11), with energies agreeing to $2\times10^{-15}$~Ry.
Both \qe{} and Elk restart this way, and carry the same three parts.

\code{max\_seconds} lets the loop stop \emph{itself} before the scheduler does,
writing a checkpoint on the way out, so a wall-clock kill lands after a
checkpoint rather than between two. It is checked between iterations, so set it
at least one iteration short of the real limit.

\emph{A derived calculator reads the checkpoint rather than the state it was
derived from.} \code{with\_spin}, \code{with\_positions} and their relatives hand
the parent's converged state to the first SCF as a starting point, and
\code{run\_scf} looks in its checkpoint directory only when nothing was handed to
it --- so a promoted run killed at its wall clock and resubmitted used to start
from the seed every time and never read the file it had been writing. The
checkpoint wins, for the reason it wins over a seed the caller passed: it is
strictly later state. Measured on fcc aluminium, a non-magnetic ground state
converged at \code{1e-8} and promoted into a two-channel run at the input's
\code{conv\_thr = 1e-12}: the first line takes two iterations, and the
resubmitted one re-enters at iteration 2 carrying the checkpoint's own
\code{ethr = 1e-13} and finishes in one, at the same $-4.1854697253$~Ry. Before,
every resubmission repeated both iterations from the seed, starting at
\code{ethr = 1e-2} and paying the re-diagonalisation the schedule then asks for.

\begin{lstlisting}[language=Python]
calculator = Calculator.from_file("al-metal.in")
calculator.get_scf()                          # the non-magnetic ground state
spinor = calculator.with_spin(2)              # carries it as a seed
scf = spinor.get_scf(checkpoint_dir="ckpt")   # resubmit after a kill: continues
\end{lstlisting}

The seed and the \code{magnetization} beside it are withheld together, which is
not bookkeeping: a resume refuses a \code{magnetization} argument by name,
there being nothing left for it to decide, and \code{with\_moments} sets one by
default --- so dropping only the state would turn a resubmitted texture run from
one that silently restarted into one that raises.

\emph{A field is carried, and only one kind of field is refused.} An applied
\code{LOCAL\_MAGNETIC\_FIELDS} card, and a penalty constraint of any flavour,
leave the field exactly as the input built it, so the resume rebuilds it from
\code{scf.in} with the rest of the calculator. Elk's \code{reducebf} fades a
\emph{scale} rather than the field, and that scale is saved and restored --- a
run whose field had decayed over forty iterations comes back where it left,
not at full field.

\begin{refuses}
\textbf{Refused.} Checkpointing is switched off, once and with a message, for a
run whose field is \emph{driven} rather than applied: the fixed-spin-moment
schemes (\code{constrained\_magnetization = 'fsm'}, \code{'atomic fsm'},
\code{'atomic fsm direction'}), where the field is the controller's own state
and is replaced after every iteration. No input file says what it reached, so a
resume would start again from the input field and converge somewhere else
without saying so. \code{checkpoint\_every} must be at least 1.

A DFT+U state used to be refused here as well, on the grounds that \code{ns}
without its manifold is an array of numbers about nothing. That is true of the
file and it is not what a resume reads it against: a state is loaded against a
system, and the manifold is rebuilt from the \code{HUBBARD} card before
\code{ns} is looked at. \code{load\_state} restores it from the
\code{calculation} when one is given, and leaves it \code{None} when only a
system is, where what is missing is the label \code{hubbard\_occupations}
reads rather than any part of the state.

\textbf{Know the cost.} The payload is the wavefunctions, which is tens of
gigabytes on a large cell, so this is a scratch-directory operation and the
cadence is a knob rather than every iteration. It has not been timed on a large
cell; the default of 10 is a guess until it is.
\end{refuses}

\textbf{A starting magnetic texture} (\code{STARTING\_MOMENTS} card): one
starting moment per \emph{atom}, cartesian, in Bohr magnetons, one line per atom
in the order of \code{ATOMIC\_POSITIONS}. \emph{Bohr magnetons} is meant
literally: a row of \code{0.0 0.0 2.0} seeds a $2\,\mu_B$ moment on that atom
whatever its valence charge, which is the same number the \code{atomic}
constraint aims at and the same one the per-site report gives back. A row asking
for more moment than the atom has valence electrons is clamped to it, with a
warning --- a superposition of atomic charges cannot express more, and the
alternative is a negative spin channel. \code{starting\_magnetization} in the
namelist follows \qe{}: below 1 it is a \emph{fraction} of the valence charge,
and at or above 1 it is Bohr magnetons and every species is divided by its own
valence charge (\code{input.f90:1448}). QE's \code{starting\_magnetization},
\code{angle1} and \code{angle2} are per \emph{species}, so a helix, a cycloid or
a skyrmion cannot be stated in a \code{pw.x} input at all.

\emph{It decides three things}, and the first is why it matters more than a
starting guess usually does. It decides the \textbf{symmetry group}; it decides
whether the run is \textbf{magnetic at all} (a noncollinear input whose texture
is given only through this card is a magnetic run, with a four-component density
--- this used to be read off the per-species \code{starting\_magnetization}
alone, so such an input converged, reported a total energy and never had a
magnetization); and it seeds the \textbf{starting density}, one moment per atom
rather than one per species. The magnetic group is filtered by the
per-atom moments $\mathbf{m}_i$ --- an operation survives only if it maps that
axial vector field onto itself, possibly followed by time reversal --- and it is
built from the input, once, never from the converged state. Built from
per-species values it is a \emph{ferromagnet's} group, which keeps operations a
texture does not have, and \code{sym\_rho} then averages the texture away a few
iterations in while the charge converges and \code{dr2} reports nothing.

A per-atom \emph{applied} field is filtered with the moments, together, under one
choice of time reversal, exactly as Elk's \code{findsym.f90} tests \code{bfcmt0}.
So a \code{LOCAL\_MAGNETIC\_FIELDS} card that asks for a texture now cuts the
group by itself and needs no card of its own.

\begin{lstlisting}
STARTING_MOMENTS
  0.00000000  1.50000000  0.00000000
 -0.75000000  1.29903811  0.00000000
 -1.29903811  0.75000000  0.00000000
\end{lstlisting}

\textbf{A texture from Python, without writing a file}
(\code{Calculator.with\_moments}, \code{System.with\_moments}). The card above is
the only way to \emph{state} a texture in an input, and a sweep over magnetic
configurations --- a 120-degree N\'eel state against a collinear one, a cone at
five angles --- wants to state several from one script. \code{with\_moments}
takes an \code{(nat, 3)} array of Bohr magnetons in \code{ATOMIC\_POSITIONS}
order, or \code{None} to drop the card, and returns a new calculator whose
k-points have been rebuilt for the new moments' magnetic group.

\begin{lstlisting}[language=Python]
import numpy as np
from defumat import Calculator

ferro = Calculator.from_file("h4-chain-ferro.in")  # tests/data/qe
angles = np.deg2rad([0.0, 90.0, 180.0, 270.0])
cycloid = 0.6 * np.stack([np.sin(angles), np.zeros(4), np.cos(angles)], axis=1)

print(ferro.system.symmetry_group().nsym, ferro.system.kpoints.nk)   # 16 9
textured = ferro.with_moments(cycloid)
print(textured.system.symmetry_group().nsym, textured.system.kpoints.nk)  # 4 12
\end{lstlisting}

\emph{The obvious workaround is silently wrong}, which is why this is a method
rather than an attribute. \code{dataclasses.replace(system,
starting\_moments=...)} swaps the field and leaves \code{kpoints} reduced with
the group the \emph{old} moments had, while \code{symmetry\_group()} is a
property recomputed from the new ones --- 9 k-points reduced with 16 operations,
symmetrised with 4, measured on the cell above. On a system built nonmagnetic it
also flips \code{domag}, so the k-set carries a \code{time\_reversal} a magnetic
run must not have. Nothing compares the two.

\textbf{Holding a texture} (\code{constrained\_magnetization = 'atomic
texture'}) constrains the \emph{full} unit vector of every atom's moment,
%
\[
  E_{\mathrm{con}} = \lambda \sum_i \left(
    1 - \frac{\mathbf{m}_i \cdot \hat{\mathbf{n}}_i}{|\mathbf{m}_i|} \right),
\]
%
with $\hat{\mathbf{n}}_i$ the normalised \code{STARTING\_MOMENTS} rows. It is not
one of QE's schemes and has no \code{i\_cons} number. QE's nearest, \code{'atomic
direction'} (\code{i\_cons = 2}), constrains $m_z/|m|$ --- the polar angle
\emph{alone} --- so it is exactly satisfied by every texture lying in a plane
that contains $z$, a helix and a collinear state alike, at zero penalty: it
cannot hold one, and reports success while failing to.

\textbf{Or hold it with a field instead of a penalty.} A penalty is a term
added to the energy, so at convergence it is still pushing --- the residual
angle \emph{is} that push, and it cannot be made zero by choosing $\lambda$
better. Elk's fixed-spin-moment scheme does the other thing: it drives an
external field until the moment sits where it was asked to, and the converged
state is then a stationary point of the \emph{unconstrained} functional under
that field. \code{'atomic fsm'} and \code{'atomic fsm direction'} are that
scheme resolved by atom (Elk's \code{fsmtype = 2} and $-2$),
%
\[
  \mathbf B_i \leftarrow \mathbf B_i - \tau\,(\mathbf m_i - \mathbf m_{\mathrm{fix},i}),
\]
%
with the targets read from \code{STARTING\_MOMENTS} in Bohr magnetons. The
direction-only variant finishes by projecting $\mathbf B_i$ perpendicular to
$\mathbf m_{\mathrm{fix},i}$ (Elk's \code{r3vo}), so the field can \emph{turn} a
moment and can never lengthen or shorten one --- and note what that buys: it has
no $1/|\mathbf m|$ anywhere in it, where the direction-only \emph{penalty} above
has one in its gradient and runs away.

\textbf{It is transcribed and it is not yet shown to win.} The argument above is
sound and the implementation is checked against Elk's own routines, but on the
two-hydrogen cell measured below it does not converge at all: at half Elk's
default gain the site residual rings over 2000 iterations with an envelope that
\emph{grows}, from $0.15$ to between $0.3$ and $0.66\,\mu_B$, while the inner
calculation's own accuracy swings between $10^{-6}$ and $6\times 10^{-2}$. That
is an unstable loop rather than a slow one. The cause is the cell rather than
the scheme: left alone this pair carries $0.000235\,\mu_B$, so its $m(B)$ is
nearly a step and no fixed-gain controller is stable against a nearly vertical
response --- and the cell was chosen for a \emph{penalty}, which does not care.

The \code{secant} update fails in a different and more interesting way. It is
perfectly well behaved, reaching $6.5\times 10^{-6}$ in the density, and it
converges to \emph{the wrong state}: the moment lengths come out right to eight
per cent and the angles wrong by 145 degrees per site, with the residual sitting
on a plateau for 900 iterations. That is its model rather than its stability. It
measures $\mathrm{d}\mathbf m/\mathrm{d}\mathbf B$ one atom and one cartesian
component at a time, and what sets a texture's angles is the exchange
\emph{between} atoms, which is exactly the off-diagonal part a diagonal model
discards. A diagonal controller can size a moment and cannot orient a pair of
them.

So use a robust magnet, and read the vector penalty's table below as the answer
until one has been measured.

\textbf{\code{lambda} means two different things three orders of magnitude
apart, and this is the trap of the family.} For a \emph{penalty} it is a
stiffness, and 1 to 10 is the useful range. For a \emph{feedback} scheme it is
Elk's \code{taufsm}, a gain that multiplies a moment error in Bohr magnetons to
give a field in Rydberg, and Elk's own default is \textbf{0.01}. Leave it unset
and you get that; write a penalty-sized value and the run warns by name.
Measured: at $\tau = 0.2$ the vector scheme runs 200 iterations without
converging and drives the moments to $0.82\,\mu_B$ against a target of $0.26$.

\textbf{Which penalty holds a texture, measured.} \emph{Not this one.}
\code{'atomic texture'} constrains a direction and not a length, so its
potential carries a $1/|\mathbf m|$ and \emph{grows} as a site's moment
shrinks --- positive feedback, and on the two-sublattice cell below no
$\lambda$ converged. \code{'atomic'} constrains the moment as a vector, so its
potential is $2\lambda(\mathbf m - \mathbf m_{\mathrm{target}})$ and bounded,
and it is what to reach for. Two hydrogen atoms of one species asked to sit 120
degrees apart (\code{tests/data/qe/h2-texture-120.in}), \code{conv\_thr = 1e-8},
targets at the converged sphere moment of $0.26\,\mu_B$:

\begin{center}
\begin{tabular}{llccrr}
\toprule
scheme & $\lambda$ & converged & iterations & angle & error/site \\
\midrule
(none)           & ---   & \textbf{yes} &  10 & $180.00^\circ$ & $30.00^\circ$ \\
\code{atomic}    & 0.05  & yes &  10 & $170.95^\circ$ & $25.47^\circ$ \\
\code{atomic}    & 0.2   & yes &   7 & $153.52^\circ$ & $16.76^\circ$ \\
\code{atomic}    & 1.0   & yes &   9 & $130.30^\circ$ & $5.15^\circ$ \\
\code{atomic}    & 3.0   & yes &  17 & $123.65^\circ$ & $1.83^\circ$ \\
\code{atomic}    & 10.0  & yes &  38 & $\mathbf{121.13^\circ}$ & $\mathbf{0.55^\circ}$ \\
\code{atomic}    & 30.0  & \textbf{no} & 200 & $11.91^\circ$ & $64.41^\circ$ \\
\code{atomic texture} & 0.1  & \textbf{no} & 200 & $118.39^\circ$ & $4.30^\circ$ \\
\code{atomic texture} & 0.02 & \textbf{no} & 400 & $140.03^\circ$ & $10.02^\circ$ \\
\code{atomic texture} & $\geq 2$ & \textbf{no} & 200 & collapsed & --- \\
\bottomrule
\end{tabular}
\end{center}

Read the first row and the sixth together: left alone the 120-degree state
relaxes to the collinear antiferromagnet in ten iterations and \emph{reports
success}, and at the largest $\lambda$ the SCF tolerates the vector penalty
holds it to \textbf{0.55 degrees per site} in 38. So a magnetic structure that
is not the ground state can be converged here, which is what a user mapping out
$E(\theta)$ needs and what the cone and spiral collapses recorded elsewhere in
this document have no constraint applied to them. The recipe is
\code{'atomic'}, a \code{STARTING\_MOMENTS} card scaled to the moment you
\emph{expect} rather than the moment you seed with, and $\lambda$ raised until
the angle stops moving or the SCF stops converging --- which on this cell are
within a factor of three of each other. \code{'atomic texture'} warns and names
this table.

\emph{How far each atom is from its target} is on
\code{SCFResult.site\_residuals}, one number per atom --- degrees for the two
direction schemes, Bohr magnetons for \code{'atomic'}. \code{constraint\_energy}
is a single scalar over every site, so under \code{'atomic texture'} one flipped
site out of fifteen reads as a small number indistinguishable from partial
convergence everywhere; the per-site array is the only form of the comparison
that can tell those apart.

\begin{refuses}
\textbf{Refused.} \code{'atomic texture'} needs \code{noncolin = .true.} (a
collinear moment has no direction beyond its sign) and a
\code{STARTING\_MOMENTS} card with no zero row (a zero vector carries no
direction). \code{STARTING\_MOMENTS} itself needs \code{noncolin} for any $x$ or
$y$ component, the same rule \code{B\_field} and \code{LOCAL\_MAGNETIC\_FIELDS}
are held to, and one row per atom.

\textbf{Checked, not remembered:} the card above parses through
\code{build\_system}, reaches \code{system.starting\_moments} unchanged, becomes
\code{m\_loc} verbatim, and \code{constraint\_targets} returns it normalised ---
three unit rows. The starting density built from it carries opposite moment
lobes on two antiparallel sites and integrates to zero over the cell.

\textbf{It reaches all three consumers}, which it did not until recently: the
charge density, the Hubbard occupation matrix, and a PAW or ultrasoft dataset's
atomic \code{becsum}. The last two used to be split by the per-\emph{species}
\code{starting\_magnetization} and so started every site of a species pointing
the same way, while the charge carried the texture --- so iteration~1
contradicted itself on exactly the systems the card exists for. For PAW that is
not merely a poorer guess: the one-centre terms are a \emph{function} of
\code{becsum}, so the two sites got different Hamiltonians than the density
asked for. And nothing in the SCF turns a moment, so a spinor DFT+U run --- which
is forced to \code{nosym}, leaving the starting matrix as the only steering
there is --- could converge with the shells polarised along the wrong axes and
report success.

\textbf{Warned about, with a number:} a textured starting \emph{density} handed
to \code{run\_scf} by the caller. The symmetry group is decided from the input
and never from the density, so a texture that reaches the run only that way is
invisible to the filter and \emph{will} be symmetrised away. \code{run\_scf}
measures $\lVert \rho - \mathrm{sym}(\rho)\rVert / \lVert \rho\rVert$ on the seed
before iteration~1 and warns above $10^{-3}$, reporting the charge and the
magnetization apart. It is not a small effect and it is not a rounding: a
four-atom 90-degree cycloid handed to the same cell's \emph{ferromagnetic} group
has its magnetization residual come out at \textbf{1.0} --- averaging four
directions 90 degrees apart over a group that permutes the four sites gives
exactly zero, and every site moment goes from $0.27\,\mu_B$ to $10^{-19}$ ---
while its charge residual is $5.7\times10^{-16}$, so the charge notices nothing.
From a \code{STARTING\_MOMENTS} card the same figure is $6\times10^{-16}$. Give
the texture as a card, or set \code{nosym = .true.}
\end{refuses}

\textbf{The moment on each atom} (\code{SCFResult.site\_moments},
\code{.site\_charges}). The two magnetic numbers a run has always printed are
integrals over the \emph{cell}, and both are zero for every compensated state
--- for an antiferromagnet, for a spiral, for a cycloid, and equally for the
nonmagnetic state any of those can collapse into. What separates them is the
charge and the moment integrated in a sphere around each atom, which is what
\qe{}'s \code{report\_mag} prints and what Elk computes as \code{mommt}:
%
\[
  q_a = \int w_a(\mathbf{r})\, n(\mathbf{r})\, d^3r,
  \qquad
  \mathbf{m}_a = \int w_a(\mathbf{r})\, \mathbf{m}(\mathbf{r})\, d^3r,
\]
%
with $w_a$ one inside a radius $r_m$ around atom $a$, falling linearly to zero
at $1.2\,r_m$, and every grid point belonging to at most one atom.

\begin{lstlisting}[language=Python]
from defumat import Calculator

calculator = Calculator.from_file("h2-mirror-afm.in")
scf = calculator.get_scf()

print(scf.magnetization)          # -2.3e-06  -- the cell total, zero
print(scf.absolute_magnetization) #  0.806933
for a, (q, m) in enumerate(zip(scf.site_charges, scf.site_moments)):
    print(f"atom {a}: charge {q:.4f}  moment {m[0]:+.4f} mu_B")
# atom 0: charge 0.6240  moment +0.2909 mu_B
# atom 1: charge 0.6240  moment -0.2909 mu_B
\end{lstlisting}

\emph{They are recorded at every iteration}, in \code{scf.history}, and not only
at convergence. That is the point rather than a convenience: a texture that is
going to be lost goes early and the run then converges cleanly, so the converged
value is exactly the one that cannot show it. The per-iteration console line
carries \code{|m|\_site = min..max}, the shortest pair that tells a live
compensated state from a dead one, and the full block is printed at the end with
$\theta$ and $\phi$ per atom --- angles \qe{} does not print and that nobody
should be reconstructing by hand from three components.

\textbf{A collinear run has a magnetic symmetry group too.} Until recently it did
not: the filter was built only for \code{noncolin}, so a compensated magnet
stated on \emph{one} species kept the operation carrying one sublattice onto the
other, and the density symmetriser then averaged $n_\uparrow$ over both sites ---
which is the statement that the staggered moment is zero. This is not an exotic
case: an \emph{altermagnet} is by definition compensated order whose sublattices
are related by a rotation, so the swap is a point-group operation of the chemical
structure. On two hydrogen atoms related by a mirror with $\pm0.6$ from the card,
the moments were gone at iteration~1 and the run converged to a perfectly good
total energy \textbf{6.6~meV} above the antiferromagnet, reporting an absolute
magnetization of $7\times10^{-6}$. With the filter the group is cut from 16
operations to the 8 that preserve the moment pattern, and the run agrees with the
\code{nosym = .true.} spelling --- the only one that used to work --- to
$10^{-8}$~Ry and $10^{-6}\,\mu_B$ in the same iteration count.

The collinear filter compares a \emph{number} per site, not a vector. Without
spin--orbit coupling a spatial operation does not turn the spin, so the moment is
permuted and never rotated; filtering it as an axial vector instead would cut a
$C_{2x}$ or a mirror on a $z$-ferromagnet and leave a group that is too
\emph{small}.

\begin{refuses}
\textbf{Refused.} An operation that reverses \emph{every} moment is a symmetry of
the crystal followed by time reversal. Those are \textbf{dropped} rather than
used, which is \qe{}'s own \code{colin\_mag = 1} behaviour: using one means
swapping the two spin channels while symmetrising, and the density symmetriser
here does not. The result is a subgroup of the true magnetic group --- correct,
and paid for in k-points rather than in physics.

\textbf{Where the site moments are not reported:} the response stack. All three
relaxation drivers carry them now --- \code{RelaxStep}, \code{VCRelaxStep} and
\code{SpiralRelaxStep} each hold \code{site\_charges} and \code{site\_moments} at
their own step, which is the record a \emph{relaxation} needs for the same reason
the SCF needs one per iteration: everything else a step reports is zero-sum over
the sites, so a compensated magnet that lost its moments between two geometries
reads exactly like one that kept them.

\textbf{The spheres are a choice, not a definition.} A plane-wave basis has no
muffin tins, so ``the moment on this atom'' is whatever region you integrate
over; $r_m$ defaults to a little under half the nearest-neighbour distance so the
tapered spheres are disjoint. The site charges therefore sum to \emph{less} than
the cell's charge --- the interstitial belongs to nobody --- and the site moments
sum to less than the cell moment. Compare them against another code only when
both used the same radius.
\end{refuses}

\textbf{DFT+U:} both rotationally-invariant functionals, on \code{atomic},
\code{ortho-atomic} and \code{norm-atomic} projectors, read from the
\code{HUBBARD} card. Forces come from differentiating through projectors that
move with the atoms.

\emph{The card selects which functional runs}, exactly as \code{pw.x} does it
--- the namelist's \code{lda\_plus\_u\_kind} is obsolete there and warned about.
A card carrying only $U$, $J_0$, \code{ALPHA} and \code{BETA} is Dudarev's
simplified functional; any $J$, $B$, $E_2$ or $E_3$ selects Liechtenstein's full
one, whose interaction is the whole Coulomb matrix
%
\[
  \mathrm{vee}(m_1 m_2 m_3 m_4)
  = \sum_{k} F^k \frac{4\pi}{2k+1}
    \sum_{q=-k}^{k} a_{kq}(m_1 m_3)\, a_{kq}(m_2 m_4)
\]
%
rather than one number, with the Slater integrals $F^k$ fixed by $U$ and $J$.
Against \code{pw.x} on antiferromagnetic FeO: \textbf{4.2e-9 Ry} at $J = 1$~eV.
\emph{Forces work for both}, and for the full functional \code{pw.x} does not
have them at all (\code{force\_hub.f90} stops on ``forces in the DFT+U+J scheme
are not implemented''); a central difference of the SCF energy agrees to
1.3e-4~Ry/bohr. \textbf{Seed that difference from the undisplaced run} --- with
a full interaction matrix the DFT+U landscape has more than one solution, and
cold-started displacements land on a different one, which does not look like an
error, it looks like a force.

\begin{lstlisting}[language=Python]
from defumat import Calculator
scf = Calculator.from_file("feo.in").get_scf()
print(scf.energy_terms["hubbard"])
\end{lstlisting}

\begin{lstlisting}
HUBBARD {atomic}
U Fe1-3d 4.3
J Fe1-3d 1.0     # selects the full functional
\end{lstlisting}

\textbf{Around-mean-field double counting} (\code{hubbard\_double\_counting =
'amf'}) is the alternative to the fully-localised limit that both QE's
functionals assume, and \code{pw.x} has neither the option nor the code. FLL
treats the occupations as integers, which is right for a localised magnetic
insulator and wrong for a metal; AMF subtracts the shell's mean occupation
before the interaction, so a uniformly filled shell is corrected by
\emph{exactly} nothing --- $0$ to $10^{-14}$ at any filling, where FLL is not.
On FeO its Hubbard energy is $-0.005$~Ry against FLL's $+0.232$. It needs the
full functional, there being no interaction matrix in the simplified one to
apply the shift to.

\textbf{Slater integrals from the orbital} (\code{hubbard\_slater = 'yukawa'})
computes the interaction instead of parameterising it: $F^k$ is the double
radial integral of the manifold's own all-electron partial wave against a
screened Coulomb kernel,
%
\[
  F^k = \iint \mathrm{d}r\,\mathrm{d}r'\;
        [r\phi(r)]^2 [r'\phi(r')]^2 \,
        (2k+1)\,\lambda\, i_k(\lambda r_<)\, \tilde k_k(\lambda r_>) ,
\]
%
which becomes the bare $r_<^k/r_>^{k+1}$ as $\lambda \to 0$. Give $\lambda$ on
the card (\code{LAMBDA Ni-3d 2.1}) or give a $U$ and let the screening length be
solved for --- either way $F^0$, $F^2$, $F^4$ and $J$ all come out of the
orbital. Nickel's unscreened 3d gives $F^4/F^2 = 0.626$ against the 0.625 that
QE and Elk both hardcode, and $J = 1.28$~eV.

\begin{lstlisting}
&system
   hubbard_slater = 'yukawa'
/
HUBBARD {ortho-atomic}
U Ni-3d 5.0      # lambda is solved for; F2, F4 and J follow
\end{lstlisting}

\textbf{On a spinor} (\code{noncolin}, with or without \code{lspinorb}) the
occupation matrix is one $2\times2$ matrix in spin space rather than two real
matrices, measured on projector columns that are themselves spinors, and its
off-diagonal blocks are what let the shell's moment point anywhere. Both
functionals work there. Against \code{pw.x}: \textbf{7.7e-9 Ry} on relativistic
BN and \textbf{1.4e-9 Ry} on fully-relativistic iron carrying $U = 2.20$,
$J = 1.75$~eV with its moment turned into the $xy$ plane, with
$\mathrm{Tr}[n^{\sigma\sigma}]$ matching every digit \code{write\_ns} prints.

\textbf{Tensor moments} put that matrix in a basis where each element is a
multipole rather than a number: an orthonormal, complete set $\Gamma^{kpr}_t$
with $k$ the rank in the orbital index, $p$ the rank in spin and $r$ their
coupling, so that
%
\[
  n = \sum_{kprt} w^{kpr}_t \, \Gamma^{kpr}_t ,
\]
%
with real $w$. The charge is $w^{000}$, the spin moment is $w^{011}$, and
$\mathbf{L}\cdot\mathbf{S}$ is $w^{110}_0$ alone --- to $10^{-11}$ over all one
hundred moments of a $d$ shell.

\textbf{Fixing one} (a \code{TENSOR\_MOMENTS} card) converges the field to a
state carrying it, which selects an orbital or multipolar ordering the field
would not find on its own. The penalty is written down and its potential is
\code{jax.grad} of it, and, like a constrained moment's, its energy is not part
of the reported total --- so the number that comes back is the unconstrained
functional at the constrained state.

\begin{lstlisting}
&system
   tensor_moment_penalty = 20.0
/
TENSOR_MOMENTS
 N-2p 1 1 0 0 -0.40      # hold L.S at -0.40
\end{lstlisting}

Being quadratic, the penalty settles where its gradient balances the material's
own restoring force, so the moment \emph{approaches} its target as the penalty
stiffens rather than reaching it: on nitrogen's $2p$, whose free
$\mathbf{L}\cdot\mathbf{S}$ is $-0.0006$, holding it at $-0.40$ gives $-0.109$,
$-0.263$, $-0.355$ and $-0.388$ at $\lambda = 1, 5, 20, 80$, with the energy
rising monotonically by $2.6$, $15.3$, $27.4$ and $32.6$~mRy.

\begin{refuses}
\textbf{Refuses.} A \emph{collinear} DFT+U run needs \code{nosym} when its
magnetic group carries time-reversed operations: there the occupation matrix is
two separate channels and such an operation exchanges them, which is not
written. \textbf{A spinor run does not} --- it keeps both spin indices, so the
operation acts on the matrix itself, and both halves of that action are in: the
$SU(2)$ representation of each operation beside the rotation of the $m$ indices,
and, for a time-reversed one, the transpose that goes with an antiunitary
operation. A fixed tensor moment needs a spinor: for a collinear
run the spin-rank moments keep only their $z$ component and the constraint
degenerates into the fixed spin moment \code{constrained\_magnetization}
already is. \code{hubbard\_slater = 'yukawa'} needs a \textbf{PAW}
dataset, and this is a measurement rather than a preference: only PAW carries
the all-electron partial wave. A norm-conserving $\chi$ is normalised and gives
a pseudo-orbital $F^0$ far too small (silicon's 3p: 8.8~eV), and an
\emph{ultrasoft} $\chi$ is not even normalised --- iron's 3d has
$\int (r\chi)^2 = 0.41$, the rest of the norm living in the augmentation charge,
and its $F^0$ comes out at 2.1~eV where the answer is above twenty. The
\code{norm} that comes back with each manifold is what says whether it is bound
inside the cutoff, and the cutoff defaults to the dataset's own augmentation
radius because a partial wave \emph{diverges} outside it. Elk's
\code{inpdftu = 2} and \code{3} (Slater and Racah \emph{input}) have no card
syntax: they are a change of coordinates on the three numbers $(U, J, B)$
already spans, and QE's card has spent the names $E_2$ and $E_3$ already. The
FLL/AMF interpolation (Elk's \code{dftu = 3}) is not here because it is not in
Elk either any more --- documented in the 11.0.2 manual, removed from the 11.0.2
source.
\end{refuses}

\begin{refuses}
\textbf{Refuses.} The intersite $V$, background channels, the
orbital-resolved variant, the \code{wf}/\code{pseudo} projector sets, and
noncollinear \code{ns\_nc}; and PAW $+$ GGA $+$ a noncollinear magnetization
together.

A \textbf{magnetic field in a noncollinear run with no starting moment} is
refused rather than run. \code{domag} comes from \code{starting\_magnetization}
and from nothing else (\code{setup.f90}), so such a run has
\code{nspin\_mag = 1} and a density with no magnetization channels for the field
to act on. Set \code{starting\_magnetization} for one species --- any nonzero
value will do, the field decides where the moment goes. Worth knowing:
\code{pw.x} \emph{accepts} this input and silently drops the field, because
\code{add\_bfield} writes it into the potential's magnetization channels and
\code{vloc\_psi\_nc} applies those only \code{IF (domag)}.

The \textbf{residual solver} cannot take a fixed occupation whose filled/empty
boundary cuts a \textbf{degenerate multiplet} --- which is exactly what an
atom's Hund's-rule filling does. Which member of the multiplet the eigensolver
returns is arbitrary, so the density map is not a function of the density and a
Newton step has nothing to converge on; it is diagnosed by name rather than left
as a \code{NaN}. The mixer damps its way past it, and a \emph{gapped}
fixed-occupation LSDA cell agrees between the two solvers to 5e-12~Ry.
\end{refuses}

\subsection{The source-free exchange-correlation field, and the torque}

Let us ask what turns a moment. The exchange-correlation field of a local
functional is built pointwise: at each point the density is resolved onto the
direction of $\mathbf m(\mathbf r)$, the collinear expression is evaluated
there, and the splitting is attached back to $\hat{\mathbf m}$. So
$\mathbf B_{xc}(\mathbf r) \parallel \mathbf m(\mathbf r)$ everywhere, and the
torque the field exerts,

\begin{equation*}
\boldsymbol\tau = \int \mathbf m(\mathbf r)\times\mathbf B_{xc}(\mathbf r)\,
\mathrm d^3 r,
\end{equation*}

is \textbf{identically zero at every density}, converged or not, constrained or
not. A functional of $|\mathbf m|$ alone cannot turn a moment, which is a
statement about the approximation rather than about this implementation, and it
means a test asserting that this quantity vanishes at self-consistency passes on
a broken code exactly as it passes on a working one.

A magnetic field in nature has no sources. \code{nosource} makes the
exchange-correlation one obey that, by projecting out its longitudinal part ---
Elk's \code{src/projsbf.f90}, which takes the divergence, solves a Poisson
equation for it and adds the gradient of the solution back, and which on a
plane-wave grid is one line in reciprocal space,
$\mathbf B(\mathbf G) \to \mathbf B(\mathbf G) - \hat{\mathbf G}\,
(\hat{\mathbf G}\cdot\mathbf B(\mathbf G))$, with $\mathbf G = 0$ left alone
since a Poisson equation cannot move a constant. The projected field is no
longer parallel to $\mathbf m$, so the torque stops being zero.

\begin{pycode}
from defumat import Calculator

# fcc nickel, noncollinear, LDA, with nosource = .true. in &system
calc = Calculator.from_file("ni-noncol-111-nosource.in", pseudo_dir="../pseudo")
torque = calc.get_exchange_torque()

print(torque.parallel_fraction)   # -> 0.982: B_xc is no longer along m
print(torque.sites.shape)         # -> (1, 3): one atom, one torque vector
\end{pycode}

\code{parallel\_fraction} is
$\int|\mathbf m\cdot\mathbf B_{xc}| / \int|\mathbf m||\mathbf B_{xc}|$, and it
is reported beside the torque so that a zero can be told from a silence: one
means the field is everywhere along the magnetization and the zero is the
approximation speaking. Without \code{nosource} it is $1$ to round-off.

\textbf{The numbers.} On the cell above, $\int|\nabla\cdot\mathbf B_{xc}|^2$
falls from $1.5\times10^{-2}$ to $1.1\times10^{-32}$, and the moment goes from
\textbf{0.6576 to 0.7045}~$\mu_B$ --- the size of change the source-free
literature reports for a transition metal. On the four-hydrogen helix the same
flag takes the site moments from 0.4675 to \textbf{0.0003}~$\mu_B$, and that
collapse is the cell rather than the method: the projection removes roughly the
radial part of each atom's field, and a Stoner-marginal magnet crosses its own
criterion when it does.

\textbf{And the torque fires.} Turning one site's moment inside its own sphere,
away from where it sat, on the same nickel cell:

\begin{center}
\begin{tabular}{lrr}
\hline
turned by & local functional & source-free field \\
\hline
$0^\circ$  & 7e-22 Ry & 2e-18 Ry \\
$15^\circ$ & 4e-20 Ry & \textbf{1.64e-4 Ry} \\
$45^\circ$ & 5e-20 Ry & \textbf{4.47e-4 Ry} \\
$90^\circ$ & 1e-21 Ry & \textbf{6.32e-4 Ry} \\
\hline
\end{tabular}
\end{center}

The left column is the point of the table: the local functional reports zero at
\emph{every} angle, not only at the one it converged to.

Note that the \emph{total} torque over the cell is zero for a second reason, and
that one is a symmetry rather than an approximation: without spin--orbit
coupling the energy does not change under a global rotation of every spin, so
the net torque on the cell cannot either. What to read is
\code{torque.sites}; the total is a check on the assembly.

\begin{refuses}
\textbf{Refuses.} Everything that differentiates the total energy ---
\textbf{forces, stress, phonons, the elastic constants and the whole linear
response} --- because the projection changes the potential and leaves the energy
expression alone, exactly as a potential-only meta-GGA does, so the converged
state is stationary for a functional nothing writes down. \code{run\_scf} warns
for the same reason: the total energy a \code{nosource} run reports is not the
value of anything it minimised and is not comparable with a run without it.
\textbf{Magnons}, and this is the sharpest one: the transverse kernel
$f_{xc}^{+-} = B_{xc}/m$ is a scalar only because the local functional makes
$\mathbf B_{xc}$ parallel to $\mathbf m$, which is precisely what the projection
breaks --- the quotient is still finite and the magnon it gives is meaningless.
A \textbf{collinear} run, since a transverse projection gives a field along $z$
components along $x$ and $y$ that an $(n, m_z)$ density has nowhere to carry
(Elk's \code{init0.f90} requires \code{ncmag} for the same reason). \textbf{PAW},
whose one-centre $\mathbf B_{xc}$ on the spheres is a second copy of the field
that this projection does not reach, where Elk projects its muffin-tin part
together with the interstitial one. A \textbf{spin spiral}, since
$\nabla\cdot\mathbf B$ is a laboratory-frame statement and the magnetization is
stored in the rotating frame. A \textbf{potential-only meta-GGA}, because
stacking two potentials with no energy behind them leaves nothing at all that
says what the run is a stationary point of. And the \textbf{ultracell}, whose
per-cell potential stops being a function of that cell's density once a
division by $|\mathbf G|^2$ couples the whole modulation.

Elk's other half is \textbf{not built}: \code{sxcscf} scales the spin part of
the field by a constant, which is the same paper's second knob and a different
approximation. Elk labels \code{nosource} an experimental feature and that
label carries here.
\end{refuses}

% =====================================================================
\section{Forces, stress and relaxation}

\begin{equation*}
\mathbf F_I = -\frac{\partial E}{\partial\boldsymbol\tau_I}\bigg|_{\psi\ \mathrm{fixed}},
\qquad
\sigma_{ij} = -\frac{1}{\Omega}\frac{\partial E}{\partial\varepsilon_{ij}}\bigg|_{\psi\ \mathrm{fixed}}
\end{equation*}

Both are one gradient. With ultrasoft the orthonormality constraint
$\langle\psi_n|\hat S|\psi_m\rangle=\delta_{nm}$ is carried explicitly, so its
multiplier term --- QE's Pulay contribution --- falls out of the same derivative
rather than being added afterwards. A variable-cell relaxation minimises the
\textbf{enthalpy}, which is why a relaxed crystal carries the applied pressure
rather than having zero stress:
\begin{equation*}
H = E + P\,\Omega,
\qquad \frac{\partial H}{\partial h} = \Omega\,(P\,\mathbf 1 - \sigma)\,h^{-T}
\end{equation*}

\begin{pycode}
from defumat.forces import compute_forces
from defumat.stress import compute_stress
from defumat.workflows.relax import run_relax
from defumat.workflows.vc_relax import run_vc_relax

forces = compute_forces(calculation, result)      # method=... for QE's own terms
print(forces.forces)                              # (nat, 3) in Ry/bohr

stress = compute_stress(calculation, result, terms=True)
print(stress.pressure)                            # kbar

relaxed = run_relax(system, pseudos, forc_conv_thr=1.0e-4)
cell = run_vc_relax(system, pseudos, press=500.0) # kbar
print(cell.pulay_error)   # the frozen-basis error, reported not hidden
\end{pycode}

QE's \code{force\_lc}, \code{force\_cc}, \code{force\_ew}, \code{force\_us},
\code{addusforce} and \code{force\_corr} are transcribed too, behind the same
registry --- \textbf{as a cross-check, not as the implementation}. The two share
no machinery, which is what makes disagreement informative; it is how the
augmentation force's sign and a missing gradient correction were found.

For a noncollinear or spin-orbit run the same expression is evaluated on
two-component spinors: the nonlocal term takes the pseudopotential's bare
\code{dvan\_so} (a complex $2\times2$ matrix in spin space) and the
orthonormality constraint takes \code{qq\_so}, and the state is one coefficient
vector of length $2\,n_{\mathrm{pw}}$ rather than two. Nothing else in the
functional changes.

\begin{pycode}
from defumat import Calculator

calculator = Calculator.from_file("pt2-soc-force.in")   # noncolin, lspinorb
forces = calculator.get_forces()          # (nat, 3) Ry/bohr
stress = calculator.get_stress()          # (3, 3)   Ry/bohr^3
print(forces.max_force, stress.pressure_kbar)
\end{pycode}

Validated against \code{pw.x} with \code{tprnfor}/\code{tstress} on a four-atom
noncollinear hydrogen chain (norm-conserving) to $8.9\times10^{-7}$ Ry/bohr,
doubled fcc platinum with spin-orbit coupling to $7.5\times10^{-6}$ (ultrasoft)
and $7.3\times10^{-7}$ (PAW), and the stress on those three plus QE's own three
\code{pw\_spinorbit} cases to $\le 1.2\times10^{-6}$ Ry/bohr$^3$; and, without
any Fortran at all, against a central finite difference of the frozen energy to
$6.2\times10^{-9}$ Ry/bohr.

\textbf{With a gradient-corrected functional} the derivative is checked
separately, because the gradient correction of a \emph{vector} magnetization is a
branch of its own: bcc iron with its moment in the plane, ultrasoft, PBE
(\code{fe-noncolin-pbe-stress.in}) agrees with \code{pw.x} to
$6.7\times10^{-9}$~Ry in energy and $1.6\times10^{-7}$~Ry/bohr$^3$ in
\textbf{stress} --- 0.019 kbar out of 152.71, and the same level the collinear
ultrasoft cases reach. With one atom in a bcc cell the force is zero by symmetry,
so the stress is the only derivative that geometry has, which is why it is the
quantity compared. \qe{} resolves the density onto a \emph{fixed} quantization
axis whenever the starting moments are all parallel to one direction
(\code{compute\_ux}), using $(n \pm \mathrm{sign}(\mathbf m\cdot\hat u)\,|\mathbf
m|)/2$ so that ``up'' stays up across a node where $\mathbf m$ passes through
zero; that is the branch iron takes, and both codes pick the same axis. A
genuinely \textbf{canted} cell takes plain $|\mathbf m|$ instead, with a real kink
at every node, and that branch has no external number yet --- not for want of
trying here: \code{pw.x} itself limit-cycles at $5\times10^{-6}$~Ry on a canted
PBE hydrogen chain that converges to $10^{-11}$ under LDA in 62 iterations.

\begin{refuses}
A noncollinear or spin-orbit run has forces, a stress and a relaxation, but only
through the \emph{autodiff} method: \code{method='analytic'} is refused by name,
because \code{force\_us}/\code{stres\_knl} are transcriptions of QE's
hand-derived expressions and have no spinor form. The force on an atom of a
\textbf{spin spiral} is refused separately --- its two spinor components live on
different plane-wave spheres, so the nonlocal term needs the projectors of both;
\code{dE/dq} is what a spiral has instead. \textbf{Elastic constants and
electrostriction} stay refused for \code{noncolin}: they differentiate the same
functional a third time, and while the Sternheimer solve their first-order
wavefunctions come from now \emph{does} have a spinor form, no third derivative
has been run with one. The refusal is about the assembly rather than the solve,
which is the distinction that cost this project a whole quantity once before.

\code{method='analytic'} is refused for one more reason, and this one is about
\emph{storage} rather than about physics: \code{addusforce} needs the whole
$Q_{ij}(\mathbf G)$ table, and a cell whose stored table would exceed
\code{AUG\_MAX\_BYTES} (2~GiB; \code{DEFUMAT\_AUG\_MAX\_BYTES} moves it)
keeps a radial table and rebuilds the array a chunk at a time instead, so there
is no whole table to read. Materialising one to keep the term would undo the
thing that scheme exists for, which is why this is a refusal. It costs the
cross-check and nothing else: the default autodiff force runs on that route, and
on a displaced \code{si2-us.in} it returns $3.92220321\times10^{-2}$~Ry/bohr on
\emph{both} storage schemes, identical to every printed digit, agreeing with the
resident cell's transcribed force to $1.4\times10^{-6}$~Ry/bohr. Before this the
term read \code{qgm[0]} off an empty tuple and raised a bare \code{IndexError}
from inside a \code{jax.jit} trace --- in a relaxation with
\code{force\_method='analytic'}, at the first ionic step with the SCF already
paid for.
\end{refuses}

% =====================================================================
\section{Response: dielectric, phonons, Raman}

The \textbf{velocity operator} comes first and everything else uses it: it is one
\code{jvp} of $\hat H(\mathbf k)$ at a frozen sphere, because
$\partial\hat H/\partial k_\alpha = i[\hat H, r_\alpha]$ in the periodic gauge.
\code{band\_velocities} exposes it --- \code{Calculator.get\_band\_velocities}
is the short way --- and because the overlap carries a velocity too, a band
velocity is
$\langle\psi|\partial_k\hat H - \epsilon\,\partial_k\hat S|\psi\rangle$.

\begin{pycode}
calc = Calculator.from_file("scf.in", pseudo_dir="pseudo")

velocities = calc.get_band_velocities()          # .velocities is (nk, nbnd, 3)
elsewhere = calc.get_band_velocities(kpoints=path)
\end{pycode}

Every first-order quantity then comes from the Sternheimer equation --- the response
of an occupied orbital to a perturbation, projected outside the occupied
manifold so that no empty states and no energy denominators appear:
\begin{equation*}
\big(\hat H - \epsilon_n\hat S + \alpha\hat Q\big)\,|\Delta\psi_n\rangle
  = -\,\hat P_c^{+}\,\Delta V\,|\psi_n\rangle
\end{equation*}
It is solved self-consistently, because $\Delta V$ contains the response of the
potential to the response of the density,
\begin{equation*}
\Delta V_{\mathrm{scf}} = \Delta V_{\mathrm{ext}}
  + \int K(\mathbf r,\mathbf r')\,\Delta n(\mathbf r')\,d\mathbf r'
\end{equation*}
and the kernel $K$ is one \code{jvp} of \code{v\_of\_rho} rather than a second
implementation. From its solutions:
\begin{gather*}
\epsilon^\infty_{\alpha\beta} = \delta_{\alpha\beta}
  + \frac{8\pi}{\Omega}\sum_{n\mathbf k} w_{\mathbf k}
    \langle\Delta_\alpha\psi_n|\partial_\beta\psi_n\rangle,
\qquad
Z^{\ast}_{I,\alpha\beta} = \frac{\partial F_{I\alpha}}{\partial\mathcal E_\beta} \\[3pt]
C_{I\alpha,J\beta} = \frac{\partial^2 E}{\partial\tau_{I\alpha}\partial\tau_{J\beta}},
\qquad
\omega^2\mathbf e = \frac{1}{\sqrt{M_I M_J}}\,C\,\mathbf e,
\qquad
\frac{\partial\epsilon_{ij}}{\partial\tau_{I\gamma}}\ \ \text{(Raman)}
\end{gather*}

The Raman tensor is the second-order energy differentiated once more along a
displacement, at \emph{frozen first-order} wavefunctions --- the $2n{+}1$
theorem, which is why a third derivative costs one more \code{jvp} and not a new
solve.

\begin{pycode}
from defumat.response import (dielectric_tensor, dynamical_matrix,
                               raman_tensors, vibrational_spectrum)

eps = dielectric_tensor(calculation, scf.wavefunctions, scf.eigenvalues,
                        scf.density)
print(eps.epsilon)                     # 3x3, cartesian

ph = dynamical_matrix(calculation, scf.wavefunctions, scf.eigenvalues,
                      scf.density, scf.becsum)
print(ph.frequencies)                  # cm^-1

raman = raman_tensors(calculation, scf, born_charges=True, keep_internals=True)
spectrum = vibrational_spectrum(calculation, scf)   # per-mode Raman and IR
print(spectrum.raman_activity, spectrum.depolarisation)
\end{pycode}

\code{born\_effective\_charges} gives $Z^{\ast}$ on its own, and
\code{mode\_activities} is the bare assembly --- a transcription of
\code{dynmat.x}'s \code{RamanIR}, kept a pure function of arrays so it is
testable without a self-consistent run.

\textbf{At $\Gamma$ this code carries a term \code{ph.x} does not, and
\code{origin\_tangent} is the knob that puts \code{ph.x}'s convention back.}
The plane wave with $\mathbf k+\mathbf G=0$ exists only at $\Gamma$ exactly, and
there the $l=1$ projector's $\mathbf k$-derivative is finite: the product
$f_1(q)\,Y_{1m}(\hat{\mathbf q})\to c\sqrt{3/4\pi}\,q_m$ is linear in the
vector, so its gradient is $c\sqrt{3/4\pi}\,\delta_{m\alpha}$ rather than zero.
Quantum ESPRESSO zeroes that row, which makes its $\mathrm{d}H/\mathrm{d}\mathbf k$
discontinuous at $\Gamma$ --- a $\mathbf k$ of $10^{-4}$ off it computes the
term. The default here is the derivative the operator has;
\code{Calculation(..., origin\_tangent=False)} reproduces \code{ph.x}. What it
is worth is the weight of $\Gamma$ in the k-sum, so a \emph{shifted} mesh cannot
see it at all, an unshifted one sees a fraction, and a $\Gamma$-only cell sees
all of it: on the O$_2$ molecule $\epsilon_{xx}$ reads $1.116446$ against
$1.110915$, and $Z^{\ast}_{xx}$ $0.10110$ against $0.13372$, where \code{ph.x}
gives $1.110915996$ and $0.13367$.

\textbf{A Born charge is validated on a polar crystal, and the reason is worth
one sentence.} In silicon $Z^{\ast}$ is zero by symmetry, so the $-0.0757$ both
codes print is a basis-set residue and agreeing about it says nothing about the
antisymmetric part, which is what a Born charge physically is. On zincblende
AlAs the two sites carry different species and $Z^{\ast}$ is a real charge:
$+2.101065$ on the Al and $-2.165827$ on the As, against \code{ph.x}'s
$+2.10106$ and $-2.16581$. The three silicon datasets agree with \code{ph.x} to
every printed digit, $1.0\times10^{-5}$ and $9\times10^{-6}$ for
norm-conserving, ultrasoft and PAW.

Ultrasoft and PAW work for the dielectric constant ($\le$1.2e-4 against
\code{ph.x} on four cases) \emph{and for the dynamical matrix} --- silicon's
optical mode comes out at \textbf{513.295} and \textbf{513.378}~cm$^{-1}$
against \code{ph.x}'s 513.275 and 513.404, which is tighter than the
norm-conserving case's 0.05. Four things make the difference and all four
switch themselves off when $S$ is the identity: the source term is
$(\mathrm{d}H/\mathrm{d}u - \varepsilon\,\mathrm{d}S/\mathrm{d}u)|\psi\rangle$;
the first-order state has an occupied block the Sternheimer solve does not
produce; the mixed state changes at \emph{frozen} states, because the
augmentation charge and the projectors travel with their atom; and the
orthonormality multipliers are variables of the functional that move --- as a
\textbf{matrix}, since a diagonal one is not invariant under the
occupied-manifold rotation the state tangent is free in. Metals work for the
response and for a norm-conserving dynamical matrix; and a symmetry-reduced
wedge works as of the rank-$n$ symmetriser.

\textbf{The Raman tensor is ultrasoft and PAW as well}, at \textbf{1.2e-4} on
both against a central difference of $\epsilon$ over re-converged displaced
cells, where the norm-conserving control on the same script is 6.8e-4. It took
\textbf{two tangents that are only right together}, and either alone is worse
than neither: the state tangent is $P_c\,\mathrm{d}\psi + \psi^{\mathrm{ort}}$,
the occupied block again; and $|b\rangle = P_c\,\mathbf r|\psi\rangle$ is
\emph{not} the solution of its own linear equation once \code{adddvepsi\_us}
has applied $S$ to it and added the augmentation dipole, so $\mathrm{d}b$ is
the tangent of a \emph{composition} rather than of a solve. What located them
is that $\partial\epsilon/\partial\tau$ is a sum of five partial derivatives
and each one can be finite-differenced on its own against a re-converged run:
three agreed to 7e-4 and two did not.

\textbf{A spinor run has a dielectric constant and Born charges}, on all three
datasets --- what an augmented one adds is in the amber box below, with its own
numbers. A nonmagnetic one is the same physics on a doubled
space, four bands of two electrons against eight of one, and it comes out the
same to $5\times10^{-14}$ in $\epsilon$ and $3\times10^{-15}$ in $Z^{\ast}$;
the symmetrised 8-point wedge of an unshifted grid gives the closed 64-point
grid's tensor to $7\times10^{-13}$, and \code{ph.x}'s own number to
$4.3\times10^{-5}$, which is the scalar route's figure to the digit.

\begin{lstlisting}[language=Python]
silicon = Calculator.from_file("si-epsilon.in")
spinor  = silicon.with_spin(4)          # 8 bands of one electron, not 4 of two
spinor.get_dielectric_tensor().isotropic    # 13.806641, against 13.806646
\end{lstlisting}

\noindent\code{with\_spin} promotes the converged state rather than starting
again, so the second ground state is not paid for twice.

\textbf{A magnetic spinor is refused, and the reason is a measurement rather
than a missing term.} An iodine atom with spin--orbit coupling and its moment
along $z$ (\code{i-atom-soc.in}, \code{nosym}, a 0.164~eV gap) gives
$\epsilon = \mathrm{diag}(1.3565721,\,1.3565721,\,1.5744824)$: uniaxial along
the moment, off-diagonals at $10^{-15}$, and \emph{nothing imposes it} --- the
run carries no symmetry, so the two equal entries are a measurement. Turn the
moment to $x$ and the distinct axis moves with it, returning the same two
numbers to nine digits, which is what makes the anisotropy a statement about
spin--orbit coupling rather than about the box. And yet \code{ph.x}, on a ground
state the two codes agree on to the printed digit, gives 1.494594 along the
moment against this 1.574482 --- \textbf{5.3\%}, where the two components across
it agree to $5\times10^{-4}$. Every threshold QE's \code{dmxc\_nc} has that a
derivative of the potential does not (a clamped $\zeta$, a zeroing above
$|\zeta| = 1$, an $|\mathbf m| \le 10^{-10}$ rule) fires at \textbf{zero} of
that cell's 157464 grid points, so it is not a convention at an edge.

What is \emph{not} in question is the bare response underneath. On the same
cell, with no kernel and no field in the calculation at all, $\chi_0$ under a
potential probe reproduces a central difference of the density to
$6.8\times10^{-7}$ for a charge probe and $1.1\times10^{-6}$ for a field along
the moment, both unchanged over a factor of eight in the step, and this is the
first fully relativistic dataset the solve has run on, so \code{dvan\_so} is
live in it. The disagreement therefore sits above $\chi_0$: in the
exchange--correlation kernel, in the field's own source term, or in \code{ph.x}.
It is unlocated between those three, and a number that is 5\% out while looking
like a working calculation is what a refusal is for.

\textbf{A subset of the atoms can be displaced on its own.} The cost of a
dynamical matrix is that $3\,n_{\mathrm{at}}$ bare perturbations and
$3\,n_{\mathrm{at}}$ first-order wavefunctions are held at once, which is what
stops it running on a large cell; \code{atoms=} solves only the perturbations
asked for and returns that subset's own block, diagonalised with its own
masses. That is the frozen-substrate approximation --- a molecule on a fixed
slab --- and the modes it leaves are the frustrated translations and rotations,
which are the physics of such a calculation rather than an artefact. A
\code{on\_row} callback receives each row as it is assembled, so a run that
outlives its job keeps what it paid for.

\begin{lstlisting}[language=Python]
phonons = calc.get_phonons(atoms=range(57))   # the molecule, slab held fixed
phonons.frequencies                            # (3 * 57,) in cm^-1
\end{lstlisting}

\noindent It works on norm-conserving, ultrasoft and PAW datasets. The
validation needs no reference of any kind, which is unusual here: a partial
dynamical matrix is not an approximation --- it solves fewer perturbations of
the \emph{same} equations --- so the whole-cell run is exact for what the two
share, and the subset's rows reproduce it to \textbf{9.3e-15} (norm-conserving),
\textbf{2.8e-14} (ultrasoft) and \textbf{4.9e-15} (PAW) Ry/bohr$^2$. The
physical check is that freezing one of silicon's two atoms turns the optical
mode's reduced mass $m/2$ into $m$, so the frequency falls by $\sqrt2$:
$519.205 \to 367.146$, a ratio of 1.414173.

\textbf{A soft mode converges in \code{ecutrho}, and the finite-difference
route hides it.} The analytic Hessian and a finite difference of the forces are
the same quantity in exact arithmetic and are not on a grid: the FFT box breaks
continuous translation invariance, so the energy carries a small ripple as an
atom moves across it (the ``egg-box'' effect). The analytic second derivative
takes the \emph{point} curvature of that ripple; a finite difference averages it
over the displacement. The two therefore differ by the amount the dense grid is
not translationally invariant --- and because that error is roughly a fixed
curvature, it matters in proportion to how \emph{small} the mode's own curvature
is. A stiff mode barely notices; a soft one can be several per cent out.

It is \code{ecutrho} and not \code{ecutwfc} that fixes it, the ripple living on
the dense grid where the density and the local potential do. Measured on
N$_2$ in a 12~bohr cube (SG15 ONCV PBE, \code{nosym}, \code{K\_POINTS gamma},
one atom frozen, the other at 1.15~\AA), with the finite difference taken as a
central difference of the forces at $\pm0.03$~\AA:

\begin{center}
\begin{tabular}{lrrr}
\toprule
 & transverse (cm$^{-1}$) & gap & stretch (cm$^{-1}$) \\
\midrule
40/160~Ry, analytic          & 314.12 &        & 1483.49 \\
40/160~Ry, finite difference & 296.59 & 17.5   & 1490.77 \\
40/320~Ry, analytic          & 306.35 &        & 1488.1 \\
40/320~Ry, finite difference & 301.9  & \phantom{0}4.4 & 1495.2 \\
\bottomrule
\end{tabular}
\end{center}

\noindent Raising \emph{only} \code{ecutrho} moves the frustrated-rotation pair
by 8~cm$^{-1}$ analytically and 5~cm$^{-1}$ by finite difference, closing their
gap from \textbf{17.5 to 4.4}~cm$^{-1}$, while the stretch is 0.5\% off either
way at both cutoffs --- which is the signature. The analytic asymmetry is 3e-8
and 1.4e-8~Ry/bohr$^2$, so neither number is a convergence failure of the solve.

\textbf{What makes this worth stating rather than discovering:} both routes
converge, each agrees with itself, and they drift \emph{together}, so nothing in
either calculation announces the error. The practical rule is that a soft mode
wants its own \code{ecutrho} convergence pair, and the shift between them is the
error bar to quote --- not the agreement between the analytic and
finite-difference routes at one cutoff, which is a weaker statement than it
looks. (Measurements contributed from a production study of a molecule on a
surface, where the modes of interest are 10--15~meV.)

\textbf{A trap worth knowing:} a response on a reduced $k$-set is a \emph{polar
vector field} and must be symmetrised as one. Running the whole grid instead is
sound only if the grid is closed under the point group --- a \textbf{shifted}
Monkhorst--Pack grid is not.

\subsection{Phonons away from $\Gamma$}

A displacement pattern $\mathbf u_s(\mathbf R) = \mathbf u_s e^{i\mathbf q\cdot\mathbf R}$
is not a displacement of the unit cell --- it is periodic only on a supercell
commensurate with $\mathbf q$, which is exactly what a frozen-phonon calculation
pays for. Linear response does not pay it, because the \emph{first-order}
quantities are still lattice periodic once their $e^{i\mathbf q\cdot\mathbf r}$
is taken out. The perturbation is $e^{i\mathbf q\cdot\mathbf r}$ times a
periodic function, so $\Delta V_{\mathbf q}|\psi_{\mathbf k}\rangle$ is a Bloch
state at $\mathbf k + \mathbf q$: the calculation stays in the unit cell, with
one extra plane-wave sphere.

That asymmetry is the whole of it. The operator, the projector and the
preconditioner belong to $\mathbf k+\mathbf q$; the eigenvalue subtracted is the
one of the band being perturbed, at $\mathbf k$:
\begin{equation*}
\big(\hat H_{\mathbf k+\mathbf q} - \epsilon_{n\mathbf k}\hat S_{\mathbf k+\mathbf q}
   + \alpha\hat Q\big)\,|\Delta\psi_{n\mathbf k}\rangle
 = -\,\hat P_c^{\,\mathbf k+\mathbf q}\,\Delta V_{\mathbf q}\,|\psi_{n\mathbf k}\rangle
\end{equation*}
and the response density is $\Delta n_{\mathbf q}(\mathbf r) = (2/\Omega)\sum_{n\mathbf k}
w_{\mathbf k}\,\psi^{*}_{n\mathbf k}\Delta\psi_{n\mathbf k}$, which is complex ---
at $\Gamma$ it collapses to twice a real part and nowhere else.

\begin{pycode}
# nosym = .true., noinv = .true., on an unshifted grid -- see the box below
calc = Calculator.from_file("si-nosym.in", pseudo_dir="pseudo")

ph = calc.get_phonons_at_q(q=(0.5, 0.5, 0.5))          # crystal coordinates, L
print(ph.frequencies)                                   # cm^-1
print(ph.matrix.shape)                                  # (nat, 3, nat, 3), complex

xpoint = calc.get_phonons_at_q(q=(0.0, -1.0, 0.0), q_cartesian=True)
\end{pycode}

\noindent \code{q} is in crystal coordinates of the reciprocal lattice unless
\code{q\_cartesian} is set, which reads it in the $2\pi/a$ units \code{ph.x}
prints. The matrix is \textbf{complex}: $D(\mathbf q)$ is hermitian rather than
symmetric, and only at $\mathbf q = 0$ --- or where the crystal has inversion
--- does it collapse to a real one.

Silicon on a 64-point unshifted grid, against the vendored \code{ph.x}: at $L$
the six modes are 101.789, 101.789, 382.235, 405.099, 488.028 and
488.028~cm$^{-1}$ against 101.843, 101.843, 382.241, 405.121, 488.019 and
488.019, the worst of them \textbf{0.054}~cm$^{-1}$; at $X$ the worst is
\textbf{0.038}. That is the same floor the zone centre has and for the same
reason (QE interpolates its radial form factors from a $\mathrm{d}q = 0.01$
table where this integrates them directly). \code{ph.x} runs the eight-point
wedge with its response symmetrised where this runs the whole grid with nothing
symmetrised, so the comparison is also the only check the two routes get
against each other.

\textbf{Only the Ewald sum knows about $\mathbf q$}, among the terms that do not
involve the response. The second derivative of the external potential at frozen
density is diagonal in the atom index --- both derivatives are taken at the same
site, so the two modulation phases cancel --- and it is therefore the same
matrix at every wavevector. The ion--ion sum is not diagonal, and it is the one
piece rebuilt.

\begin{refuses}
\textbf{Refuses.} A symmetry-reduced $k$-set: a response has to be averaged over
the operations that leave the \emph{perturbation} invariant, and for a phonon at
$\mathbf q$ those are the ones with $S\mathbf q = \mathbf q + \mathbf G$ --- the
small group of $\mathbf q$, which is not the crystal's group and is not written.
Run with \code{nosym} and \code{noinv} on an \textbf{unshifted} grid, which is
the grid that is closed under the point group. There is no \textbf{dispersion}
either: that needs the star of $\mathbf q$ and \code{q2r}/\code{matdyn}'s
Fourier interpolation on top of this. Ultrasoft and PAW are refused ($S$ moves
with the atoms, and the orthonormality multipliers carry a term whose bra is at
$\mathbf k$ and whose ket is at $\mathbf k+\mathbf q$, where $q_{ij}$ pairs
projectors on one sphere); so are metals, spin polarization, spinors, spin
spirals, meta-GGA and DFT+U. A \textbf{nonlinear core correction} is refused for
a reason that belongs to $\mathbf q$ rather than to the dataset: two atoms' core
charges overlap inside a nonlinear $E_{\mathrm{xc}}$, so unlike every other
frozen term that one \emph{is} a function of $\mathbf q$.
\end{refuses}

\begin{refuses}
\textbf{Refuses.} $\chi^{(2)}$ and the electro-optic tensor (the second-order
source term has nothing to build the $\langle u_i|r_k|u_j\rangle$ piece from ---
and it is \textbf{42\%} of the answer, measured);
a partial dynamical matrix (\code{atoms=}) of a
\emph{symmetrised} run, since \code{symdvscf} acts on the $3\,n_{\mathrm{at}}$
stack as one object --- an operation rotates the cartesian direction \emph{and}
permutes the atom, so it mixes a mode inside the subset with one outside it,
and such a run wants \code{nosym}; the dynamical matrix of an ultrasoft or PAW \emph{metal}
(the \code{wg}/\code{wk} weight split was derived for a response whose
\code{becsum} dependence is entirely inside \code{dpsi}, and an insulator cannot
say which weight the three further tangents belong with); any response under a
potential-only meta-GGA; and a shifted grid with a reduced $k$-set.

A \textbf{spinor} run (\code{noncolin}, with or without \code{lspinorb}) has a
dielectric constant and Born charges on all three datasets, provided it carries
\emph{no net moment}. A \textbf{textured} one
(\code{nspin\_mag = 4}) is refused by name and the refusal carries its own
measurement: it runs, its tensor is uniaxial along the moment and follows the
moment when that is turned, and it is 5.3\% from \code{ph.x} in the component
along the moment with the two across it agreeing to $5\times10^{-4}$. An
\emph{ultrasoft or PAW} one runs, and the object it was refused for is not
written here either: $\mathrm dD_{ij}$ is a $2\times2$ matrix in spin space
(\code{set\_int3\_nc}), and one \code{jvp} of \code{newd} already passes
through the noncollinear recombination, whose \code{fcoef} sandwich is linear,
so the tangent comes out dressed. What had to be written was the \emph{position}
operator's spin blocks, the augmentation charge having a dipole that couples the
two components (\code{compute\_qdipol\_so}). Fully-relativistic ultrasoft AlAs
reaches \code{ph.x} to $3.5\times10^{-5}$ on a spin-orbit shift of
$8.6\times10^{-3}$, and its Born charges to $3.2\times10^{-6}$ and
$3.9\times10^{-5}$ on the two atoms ($2.101143$ against $2.10114$, $-2.165831$
against $-2.16587$); a relativistic \emph{PAW} insulator has no reference here,
so that combination is validated only through the identity against the scalar
run. A spinor \emph{metal} is refused, and for the reason every metal is: an
\code{epsilon\_infinity} is not defined there.

The \textbf{phonons, the Raman tensor, the strain response} and everything
built on them stay refused for a spinor, and for their own reason rather than
this one: \code{symmetrize\_displacement} carries the axial landmine the field
path had --- three of the four \code{nspin\_mag} channels are a magnetization
and rotate with $\det(R)$ and the time-reversal sign --- with the atom
permutation on top, and a symmetriser that has never been checked against the
whole zone is worse than a refusal. That is not rhetoric: on the field path the
wrong rotation was \textbf{66\%}, worse than no symmetrisation at all.

Also refused: a ground state converged under a \textbf{magnetic field or a
constrained moment}. The whole stack rebuilds its potential from the field the
\emph{input} asked for, and that is not the field the density belongs to ---
\code{reducebf} scales a field away as the run proceeds (7\% of its input
value after 25 iterations at 0.9) and the fixed-spin-moment scheme replaces it
outright, which is why \code{SCFResult} carries \code{magnetic\_field} and
\code{field\_scale} separately. Threading that pair through is the missing
piece for a field put in by hand; a \code{constrained\_magnetization = 'fsm'}
moment needs the \emph{induced} field as well, and its field comes from a
feedback update rather than from a penalty, so it is not a derivative of
anything. Re-run the ground state without the field.
\end{refuses}

\subsection*{LO--TO splitting and the static dielectric constant}

At $\Gamma$ a polar crystal's dynamical matrix is not analytic. An optical mode
that moves the ions against each other builds a macroscopic electric field, and
what the field costs depends on the \emph{direction} the zone centre is
approached from, so the matrix acquires a rank-one term:
\begin{equation*}
D_{I\alpha,J\beta} \;\to\; D_{I\alpha,J\beta}
  + \frac{4\pi e^2}{\Omega}\,
    \frac{(\mathbf q\!\cdot\!Z^{\ast}_I)_\alpha\,
           (\mathbf q\!\cdot\!Z^{\ast}_J)_\beta}
          {\mathbf q\cdot\epsilon^\infty\cdot\mathbf q}
\end{equation*}
It raises the longitudinal mode and leaves the transverse ones alone --- the
LO--TO splitting --- and it vanishes identically where $Z^{\ast}$ does, so a
non-polar crystal is untouched. Contracting the same two ingredients the other
way gives the ionic screening of a \emph{static} field, one oscillator per polar
mode:
\begin{equation*}
\epsilon^0_{ij} = \epsilon^\infty_{ij}
  + \frac{4\pi e^2}{\Omega}\sum_\nu
    \frac{p^\nu_i\,p^\nu_j}{\omega_\nu^2},
\qquad
p^\nu_i = \sum_{I\gamma} Z^{\ast}_{I,i\gamma}\,z^\nu_{I\gamma}
\end{equation*}
where $p$ is the same mode dipole the infrared activity is the square of. So
$\epsilon^0$ and $\epsilon^\infty$ are two ends of one contraction, and
together with the frequencies they satisfy the Lyddane--Sachs--Teller relation
$\epsilon^0/\epsilon^\infty = (\omega_{\mathrm{LO}}/\omega_{\mathrm{TO}})^2$.

\begin{pycode}
spectrum = calc.get_vibrational_spectrum(neutralize=True)
print(spectrum.static_permittivity)          # eps_0, 3x3
print(spectrum.ionic_permittivity)           # the polar modes' share of it

lo = calc.get_vibrational_spectrum(loto_direction=(1, 0, 0), neutralize=True)
print(lo.frequencies)                        # the LO modes along x
\end{pycode}

\code{nonanal} adds the term to a dynamical matrix on its own and
\code{loto\_modes} diagonalises the result; \code{polar\_mode\_permittivity}
is the oscillator sum, and \code{neutral\_born\_charges} the sum-rule
correction below. All four are pure functions of arrays, so they can be checked
without a self-consistent run behind them.

Both halves agree with \code{dynmat.x} --- which applies the same term and the
same oscillator sum to the tensors this code writes to its \code{fildyn} --- to
every digit it prints. \textbf{The Lyddane--Sachs--Teller relation is the check
neither half can pass alone}, since one side is a ratio of two diagonalisations
and the other a sum of dipoles over $\omega^2$; on AlAs it holds to
\textbf{5.0e-11}.

\textbf{It needs a charge-neutral $Z^{\ast}$ and that is worth knowing before
reading an LO frequency.} The sum rule $\sum_I Z^{\ast}_I = 0$ says a rigid
translation builds no field; a computed $Z^{\ast}$ misses it by the basis-set
error, which \emph{charges the crystal} --- AlAs at \code{ecutwfc = 10} misses
it by $-1.257$, and the non-analytic term then lifts a longitudinal
\emph{acoustic} mode from 1.8 to 33.8~cm$^{-1}$ and puts the LO frequency
7.7~cm$^{-1}$ low. \code{neutralize=True} shares the violation out over the
atoms, which is what \code{dynmat.x} does under any \code{asr} but
\code{'no'}. Comparing against another code does not see this: it is handed the
same charges and reproduces the same wrong number. LST does.

\begin{refuses}
\textbf{Refuses.} Everything the Raman tensors and the Born charges refuse,
since it is built from both: PAW $Z^{\ast}$, metals, and the rest of the list
above. The splitting is a $\Gamma$ quantity by construction --- it is what
$q\to0$ leaves behind --- so a dispersion through it needs phonons at
$q\neq0$, which are not implemented.
\end{refuses}

\subsection*{Elastic constants, electrostriction and the elasto-optic tensor}

The same machinery in the \emph{strain} coordinate rather than the atomic one.
\textbf{The strain \emph{response} itself works on all three pseudopotential kinds}
(\code{strain\_response}): the augmentation charge $Q_{ij}(\mathbf r)$ is a
function of the \emph{cell}, so a homogeneous strain deforms the table it is
tabulated on where a displacement only translates it, and on top of that come
the four terms the dynamical matrix needed. Against a central difference of
\emph{re-converged} strained runs --- which needs no \code{ph.x}, since
\code{ph.x} has no strain perturbation --- it agrees to \textbf{4.6e-4}
(ultrasoft) and \textbf{4.7e-4} (PAW) against a norm-conserving 1.9e-4 on the
same cell. The \emph{derivatives} below are norm-conserving still --- but the variational
second-order energy they all stand on is not: it reproduces \code{dielec.f90}'s
dielectric constant to \textbf{3.4e-10} (ultrasoft) and \textbf{6.9e-11} (PAW)
against a norm-conserving 8.4e-10, which took four terms that each vanish when
$S$ is the identity. One of them is worth stating on its own, because it is a
distinction rather than a term: a \emph{state} is projected with
$1 - \sum|\psi\rangle\langle\psi|S$ and a \emph{right-hand side} with
$1 - \sum S|\psi\rangle\langle\psi|$, and collapsing the two --- which is exact
at $S = 1$ --- costs a factor of fourteen. \textbf{The third derivative in this \emph{strain} coordinate runs on all three
pseudopotential kinds too}, and what had it refused was a reference rather than
a term. Three tangents are what it takes: the two that closed the Raman tensor,
and the commutator that sources $db$ written with the multiplier \emph{matrix}
rather than the frozen scalar eigenvalue, which is the right form for a reason
that needs no measurement: the equation defining $b$ carries that eigenvalue on
both of its sides, so differentiating it produces
$\mathrm{d}\varepsilon_n$ twice, in the operator and in the source. Against a
central difference of $\epsilon$ over re-converged strained cells,
extrapolated in the step: \textbf{7.2e-6} (ultrasoft) and \textbf{2.2e-5}
(PAW), on a floor of \textbf{2.5e-5} measured by running the norm-conserving
control through the identical script, where the analytic answer is exact.
The frozen scalar reads 1.3e-2 there and the source written with the traced
diagonal alone 1.4e-3.

That term had been excluded by measurement, and the measurement was a step
size: it read the \emph{displacement} coordinate's Raman tensor going from
1.2e-4 to 1.14e-3 against a central difference at $h = 0.02$, and that
difference is itself 1.1e-3 from its own limit, reading $-69.1942$, $-69.1350$
and $-69.1205$ at $0.02$, $0.01$ and $0.005$ on ultrasoft silicon, so its
truncation error had been cancelling the missing term. Extrapolated, the
two forms swap places. \textbf{The elastic constants are the one thing here
that stays refused on an augmented dataset}, and for terms of their own:
$C_{ijkl}$ is a second derivative with no electric-field solution in it, so
none of the above reaches it, and measured against a five-point second
difference of the total energy it is \textbf{22 per cent} out on $C_{1111}$ and
2.9 per cent on $C_{1212}$, the same on ultrasoft and PAW where the
norm-conserving control is exact to 6.2e-08. Ask
\code{electrostriction(elastic=False)} for everything else; $M$ and $Q$ need
the compliance and go with them.
\code{strain\_response} solves the first-order problem for a strain --- Abinit's
metric-tensor formulation, obtained for nothing here because the cell was
already written in reduced coordinates --- and two things follow from it:
\begin{equation*}
C_{ijkl} = \frac{\partial\sigma_{ij}}{\partial\varepsilon_{kl}},
\qquad
\frac{\partial\chi_{ij}}{\partial\varepsilon_{kl}}\ \ \text{(elasto-optic)}
\end{equation*}
The second is a \textbf{third} derivative of the energy, and it comes from one
\code{jvp} of the second-order energy at frozen first-order wavefunctions ---
the $2n{+}1$ theorem again. Through the thermodynamic identity of Tanner,
Bousquet and Janolin it gives the four electrostriction tensors $m$, $q$, $M$
and $Q$.

\begin{pycode}
from defumat.response import elastic_constants, electrostriction

es = electrostriction(calculation, scf)     # solves the strain response itself
print(es.m, es.q, es.M, es.Q)              # the four electrostriction tensors
print(es.photoelastic)                      # the elasto-optic tensor
print(es.elastic.voigt)                     # C_ij in Voigt notation, GPa

# or the elastic constants alone, from a strain response you already have
C = elastic_constants(calculation, scf.wavefunctions, scf.eigenvalues,
                      scf.density, es.strain)
print(C.tensor.shape, C.voigt)

# on an ultrasoft or PAW dataset, everything but the elastic constants and
# the two families built on them: es.M and es.Q come back as None
es = electrostriction(calculation, scf, elastic=False)
print(es.depsilon_dstrain, es.photoelastic, es.m, es.q)
\end{pycode}

Converged silicon gives $C_{11}$ = 198.5 GPa and $C_{12}$ = 68.9, against a
measured 165.7 and 63.9; the three independent components of
$\partial\epsilon/\partial\varepsilon$ match a central difference over
re-converged strained cells to \textbf{2e-4}, which is that difference's own
floor, and the rank-4 tensor is cubic to 3e-14 with nothing imposing it.

\begin{refuses}
\textbf{Refuses.} Both are \code{nspin = 1}, insulators and
\textbf{clamped-ion}. \code{d(chi)/d(strain)} and the elasto-optic tensor run
on norm-conserving, ultrasoft and PAW datasets; \textbf{the elastic constants
are norm-conserving only} (\code{require\_a\_measured\_elastic\_regime}), with
22 per cent on $C_{1111}$ measured on the augmented pair rather than assumed,
and $M$ and $Q$ go with them because they need the compliance, so ask for
\code{elastic=False} to get the rest. A diverged first-order solution is
refused rather than consumed in silence
(\code{require\_converged\_responses}): at QE's \code{alpha\_mix = 0.7} the
strain response of a \emph{slab} diverges, and consuming it gave a $C_{ijkl}$
that was not even symmetric under $C_{ijkl} = C_{klij}$. The elastic constants
additionally need the whole $k$-grid rather than a wedge, because their
functional builds its own density and symmetrises it as a scalar inside the
chain rule.
\end{refuses}

\subsection{The piezoelectric tensor}

The polarization a strain induces is the stress an electric field induces ---
one mixed second derivative read either way round:
\begin{equation*}
e_{k,ij} \;=\; \frac{\partial P_k}{\partial\varepsilon_{ij}}
        \;=\; \frac{\partial\sigma_{ij}}{\partial\mathcal E_k}
        \;=\; -\frac{1}{\Omega}\,
              \frac{\partial^2 E}{\partial\varepsilon_{ij}\,\partial\mathcal E_k}
\end{equation*}
It is the Born charge's construction with one coordinate changed. $Z^\ast$ is a
\code{jvp} of the \emph{force} along the field's response; the force is
\code{jax.grad} of the frozen-state energy in the atomic positions and the
stress is \code{jax.grad} of the same functional in a strain, so this is a
\code{jvp} of the \emph{stress} along the same response --- three of them, one
per field direction, on top of a dielectric constant that was going to be
solved anyway. The orthonormality constraint contributes nothing here, because
$\langle\psi|\hat S|\psi\rangle$ is a sum over a sphere of integers and carries
no cell.

\begin{pycode}
calc = Calculator.from_file("alas.in", pseudo_dir="pseudo")

piezo = calc.get_piezoelectric_tensor()
print(piezo.voigt)              # (3, 6) in C/m^2, columns xx yy zz yz xz xy
print(piezo.e14)                # -0.7638: zincblende's one component
print(piezo.e.shape)            # (3, 3, 3) in e/bohr^2, e[k, i, j]
print(piezo.dielectric.epsilon) # the field response is shared, not re-solved

# The same number contracted rather than taped: no forward-over-reverse tape,
# so it is the route for a dense mesh or a large cell.
cheap = calc.get_piezoelectric_tensor(method="zstar_eu")
print(abs(cheap.e14 - piezo.e14))   # 3.4e-15
\end{pycode}

\code{pw.x} computes no piezoelectric tensor --- the word occurs once in the
whole distribution, in a citation in a comment --- and Elk's task 380 reaches
it by a finite difference of the Berry-phase polarization over one full ground
state per strain tensor. So the validation is internal, and there are four
independent statements. Silicon is centrosymmetric and its whole tensor
vanishes (2.4e-5 C/m$^2$ against AlAs's 0.764 from the same code, with nothing
imposing it on a \code{nosym} run); AlAs is $\bar43m$ and only
$e_{14} = e_{25} = e_{36}$ survive, to 1.7e-14; an eight-point wedge reproduces
the sixty-four-point closed grid to \textbf{4.5e-9}, and on an ultrasoft cell,
where one term of this derivative is quadratic in a per-$k$ tangent and no
average of the finished tensor can repair it, a three-point wedge reproduces its
own closed grid to \textbf{4.8e-6} --- which is the shared field response's
floor on that cell and not the tensor's, since $\epsilon_\infty$ itself splits
between the same two runs by five times as much; and the same mixed
derivative contracted the other two ways --- \code{zstar\_eu.f90}'s expression
with a strain label, and the strain response against the field's bare
perturbation --- agree to \textbf{6.2e-15} and \textbf{1.3e-7}. What anchors
the sign and the normalisation is that the \emph{same} assembly run in the
position coordinate is the Born charge, and that is \code{ph.x}'s number.

\textbf{All four of those statements are blind to the $k$-mesh, and the
$k$-mesh is worth thirteen per cent.} The three routes contract the same field
response, and the $Z^\ast$ anchor is the same assembly in the position
coordinate, so none of them can see the Brillouin-zone sum this quantity is.
Measured against Elk's route on the same AlAs cell, one converged ground state
per strain and a finite difference of the Berry-phase polarization, which
shares no machinery with any of the above:

\begin{center}
\begin{tabular}{@{}lrrrr@{}}
\toprule
$e_{14}$, C/m$^2$ & \code{4 4 4} & \code{6 6 6} & \code{8 8 8} & $10^3$ \\
\midrule
this code, a \code{jvp} of the stress & $-0.763786$ & $-0.687475$ & $-0.672897$ & $-0.669907$ \\
Berry-phase finite difference        & $-0.661386$ & $-0.661964$ &             & \\
\bottomrule
\end{tabular}
\end{center}

\code{method} is what that second call is for, and on a big run it is not a
detail: both routes consume the *same* field response, so neither checks the
other's solve, but \code{'autodiff'} holds a forward-over-reverse tape of every
radial and reciprocal-space intermediate a cell derivative rebuilds, where
\code{'zstar\_eu'} contracts the response against the bare strain perturbation
and holds nothing. Measured on ultrasoft AlAs, the first peaks at 139.6~GiB at
64 $k$-points and grows by about 1.6~GiB a point, which puts a $6^3$ mesh on a
whole cluster node. The default is the taped one because it is the route that
extends to the next derivative; reach for the other when the mesh is what you
are converging, which on an ultrasoft cell is the difference between a
workstation and a cluster node.

\code{zstar\_eu.f90} line~90 hands an augmented dataset to
\code{zstar\_eu\_us.f90}, and this route was refused on one until what those
three hundred further lines are worth in the strain coordinate had been found:
it is two contractions, because only one leg of this derivative moves $S$. They
are the multipliers' own first-order change together with the
\code{add\_for\_charges} sandwich beside it, and the screening of the
frozen-state density response --- the second of which is not about the
constraint at all, but about what ``bare'' means when the coordinate is a
strain, since freezing the density as an \emph{array} misses that the density
is a function of the cell as well as of the states. With all of them in, the
two routes read $1.473304405$ and $1.473304408$ on a small ultrasoft AlAs cell
where they had been 1.6 per cent apart, and $+0.8213303507$ against
$+0.821330452$ on the cell above at 64 $k$-points, which took 3.6~GiB and
nineteen minutes on the contracted route where the taped one took 139.6~GiB.
\textbf{PAW is still refused by name}, because its one-centre energy is a
function of \code{becsum} directly and the cross term that reaches is on no
grid.

The $-0.7638$ printed above is the \code{4 4 4} of the committed input, and it
is thirteen per cent from where the calculation is going; the Berry-phase value
stays inside 0.7 per cent over four string meshes and two ground states, so it
is the response route that is unconverged rather than the reference. The
successive differences shrink by about a factor of five a rung, so on that
pattern roughly $0.0008$ would be left after $10^3$ and the two routes
would sit about one per cent apart, which is comparable with what the reference
itself is known to. That last step is an extrapolation from four points rather
than a measurement, so what is established is the thirteen per cent, not where
the last per cent goes. The reason
is that $e_{14}$ is a mixed derivative of an energy that is already stationary
in the density, so what is left is a zone sum with no variational protection,
and it needs a denser mesh than the ground state does: the \code{6 6 6} and
\code{4 4 4} rows above share one SCF, to $-16.89293132223534$~Ry in every
printed digit, and that same change of mesh moves the Berry value by $0.0006$
and this code's value by $0.076$. \textbf{Converge the mesh on this quantity
rather than on the total energy.}

\textbf{The code says so now, and it can measure it.} Nothing inside this
quantity can see its own zone sum --- the three routes share one field response
and the $Z^\ast$ anchor is the same assembly in the position coordinate --- so
\code{get\_piezoelectric\_tensor} warns with the curve above, and the result
carries \code{nk} and \code{grid}, the sample it actually integrated. What
replaces the warning with a number for \emph{your} crystal is a ladder: one
ground state and one response per mesh, and the relative change across the last
step.

\begin{pycode}
ladder = calc.get_piezoelectric_kmesh_ladder(meshes=(4, 6))

print(ladder.nk)       # (64, 216): what each rung integrated
print(ladder.e14)      # (-0.763786, -0.687475) C/m^2
print(ladder.drift)    # 0.111: the last step moved the tensor 11 per cent
print(ladder.settled)  # False, and the warning above says so with that number
print(ladder.tensor)   # the densest rung, carrying kmesh_drift
\end{pycode}

The rungs are the irreducible wedge of each unshifted grid by default, which is
the same sample at a fraction of the cost --- on an ultrasoft cell the wedge and
the closed grid gave the same tensor to $1.6\times10^{-6}$ --- and
\code{wedge = False} runs the whole grid. A rung is a whole calculation, so a
ladder costs what its densest rung costs: those two rungs took 73~s and 191~s
and peaked at 14.4~GB, all of it the taped route's tape at 216 $k$-points, which
is the case for \code{method = 'zstar\_eu'} on anything larger.
\textbf{The last step is smaller than
the error}: across this crystal's four rungs the steps are 10, 2.1 and 0.44 per
cent while the densest is still 1.2 per cent from the Berry value, so the
remaining distance runs about three times the last step and \code{settled}
means the mesh has stopped being the largest thing wrong rather than that it is
converged.


\begin{refuses}
\textbf{Refuses.} \textbf{Clamped-ion} --- the atoms are carried by the strain
and do not relax, which is also what Elk's task computes; the internal-strain
term that makes the constant comparable with experiment is not here, and for
zincblende it very nearly cancels this one. And \textbf{polar crystals by
name}: what the derivative above gives is the \emph{improper} tensor, and the
proper piezoelectric response differs from it by
$\delta_{ki}P_j - \delta_{ij}P_k$, which needs the spontaneous polarization
itself. Those terms vanish identically whenever the two Cartesian labels they
pair are different --- so $e_{14}$ of a zincblende crystal never carries the
ambiguity --- and they vanish for every component of a class that admits no
invariant vector, which is what is checked. \textbf{PAW}, and \textbf{an ultrasoft dataset below eight $k$-divisions in
each direction} --- the ultrasoft half of what used to be one refusal was lifted
on 2026-09-20 by the ladder, which puts $e_{14}$ 1.26 per cent from a
Berry-phase finite difference at \code{8 8 8} and 0.57 at \code{10 10 10} on
zincblende AlAs, against the norm-conserving calibration's 1.62 and 1.19 on the
same meshes; below eight the dataset and the $k$-mesh are not separable, the
same cell reading 15.6 per cent out at \code{4 4 4}, so a coarse ultrasoft run
is refused where a coarse norm-conserving one is only warned about. A measured
ladder overrides the threshold: \code{kmesh\_drift} under one per cent is
evidence where a division count is a guess taken from one cubic crystal. What
follows is the history of that refusal, and it is a gap rather
than a missing term: nothing in the assembly is norm-conserving and the
displacement leg of it is the Born charge, which is now validated on a
\emph{polar} ultrasoft crystal as well (\code{alas-epsilon-us.in}, $2.101065$
and $-2.165827$ against \code{ph.x}'s $2.10106$ and $-2.16581$), so what is
left refused is the \emph{strain} leg alone. The one run taken on an ultrasoft
cell so far does not lift it: \code{alas-piezo.in} gave $+0.815929$ against a
Berry-phase finite difference's $+0.687757$, but that is the same \code{4 4 4}
mesh at which the norm-conserving calibration cell is itself thirteen per cent
out, so it measures the $k$-convergence above rather than the dataset. What has
since been settled is the \emph{internal} half of that question: the two
assemblies agree on an augmented dataset to $2.6\times10^{-9}$~C/m$^2$ where
they had been 1.6 per cent apart, so what the refusal is waiting on is a
comparison against something outside this code and not a missing term. What the
refusal is no longer standing in front of is \emph{untested} code: the PAW half
of the wedge completion, which reaches the one-centre terms that no grid
integral follows, is now exercised on a committed PAW zincblende cell, where it
takes a three-point wedge from $6.96\times10^{-4}$ away from its own closed grid
to $1.58\times10^{-6}$, a factor of 440. Beyond that: insulators, \code{nspin = 1}, and
everything the Sternheimer regime and the differentiable cell refuse (a metal, a
spin spiral, a magnetic field, a shifted grid run with \code{nosym}).
\end{refuses}

\subsection{Spin-polarized response}

The Sternheimer equation is solved per spin channel,
\begin{equation*}
  \bigl(H_\sigma - \varepsilon_{n\sigma} S + \alpha_{pv} Q_\sigma\bigr)
  \,|\Delta\psi_{n\sigma}\rangle
  = -P^{+}_{c,\sigma}\,\Delta V_\sigma\,|\psi_{n\sigma}\rangle ,
\end{equation*}
with the occupied manifold $P_{c,\sigma}$ built from \emph{that channel's}
filling: $N^{\rm occ}_\uparrow = \mathrm{NINT}(n_\uparrow)$ and
$N^{\rm occ}_\downarrow = \mathrm{NINT}(n_\downarrow)$, which for a magnetic
insulator are different numbers. Quantum ESPRESSO never meets this because LSDA
doubles \code{nks} there and each spin channel is a separate $k$-point.

\begin{pycode}
# chi_0 for a spin-polarized run -- the part that is validated unconditionally.
from defumat.response.sternheimer import make_sternheimer

solver = make_sternheimer(calculation, result, spin_polarized=True)
drho = solver.chi0(dv)          # dv is (2, n1, n2, n3) on the dense grid
\end{pycode}

Validated against a central difference of the density under a spin-dependent
probe potential --- $1.8\times10^{-6}$ on an antiferromagnetic hydrogen chain (a
smeared metal) and $1.1\times10^{-6}$ on triplet O$_2$ (the sliced branch,
ultrasoft) --- and against the identity that a cell with no magnetization gives
the same $\varepsilon_\infty$ at \code{nspin = 1} and \code{nspin = 2}:
$6.2\times10^{-14}$ on 13.806646105. The probe must differ between the channels,
because $\chi_0$ is block diagonal in spin and one equal in both would not tell
the blocks apart.

The \emph{screened} response and the Born effective charges are validated for
\code{nspin = 2} as well, on triplet O$_2$ against the vendored \code{ph.x}:
$\varepsilon_\infty = \mathrm{diag}(1.11091441, 1.11091441, 1.19799977)$ against
$(1.110915996, 1.110915996, 1.198004867)$, and $Z^*$ against the \emph{raw}
(no acoustic sum rule) block, $0.13372$ and $0.20042$ against $0.13367$ and
$0.20023$. Two things had to be right for those and neither was a new
derivation. The exchange--correlation kernel is \emph{defined} to be zero where
the magnetization reaches the charge --- \code{dmxc\_lsda} does not regularise
$|\zeta| = 1$, it skips it --- so what was missing was that convention, applied
to the \emph{argument} the derivative is taken at rather than to its value.
And the orthonormality multipliers carry a weight on the column of
$\Lambda_{mn}$ and need a \emph{mask} on its row: the solver keeps
$\max(N^{\rm occ}_\sigma)$ bands in every channel because the shape is static,
so a magnetic insulator's minority channel carries empty bands whose rows are
not removed by the zero weight. That leak is worth a factor of three on
$Z^*_{xx}$, and nothing already validated can see it --- the term it feeds is
$dS/du$, which vanishes for a norm-conserving dataset, and the mask is all ones
whenever the two channels are filled to the same depth.

\begin{refuses}
\textbf{Refuses.} \code{tot\_magnetization} together with a smearing: the two
channels have separate Fermi levels and \code{Smearing} carries one scalar
\code{ef}. A filling whose boundary lands inside a \emph{degenerate multiplet}
--- a Hund's-rule atom is exactly this case --- where the solve converges and
returns an answer 100\,\% away from a finite difference, because which member
falls below the cut is arbitrary; that is the physics, not a missing term. The
dynamical matrix, the strain response, the elastic
constants, the Raman tensor and the electrostriction tensors for
\code{nspin = 2}: the solve is spin-polarized, their second-derivative assembly
is not --- \emph{except} the Born effective charges, which are validated above. The \emph{screened} response of a magnetic system with vacuum: for
\code{nspin = 2} the kernel \code{dv\_of\_drho} is the second derivative of the
LSDA energy in the two channel densities and diverges wherever a channel density
reaches zero --- measured on triplet O$_2$, 1504 of 91125 grid points and exactly
1504 NaN --- \emph{lifted for the LDA}, which is the paragraph above, and
standing for a \textbf{GGA}, whose \code{dgcxc\_spin} carries its own density
and $\zeta$ gates that nothing here has transcribed. And any potential-only
meta-GGA (\code{tb09}, \code{bj06}), which has
no total energy to differentiate --- refused by name now rather than surfacing as
\code{v\_of\_rho} asking for \code{tau}.
\end{refuses}

% =====================================================================
\subsection{Magnons: the transverse spin susceptibility}

A magnon is the collective precession of a magnet's own magnetization: a spin
wave. It lives at the pole of the transverse susceptibility, which obeys a
Dyson equation with \emph{no} Coulomb term --- a transverse spin fluctuation
moves no charge, and the Hartree interaction belongs to the charge channel:
\begin{equation*}
  \chi^{+-} = \bigl[1 - \chi_0^{+-} f_{xc}^{+-}\bigr]^{-1}\chi_0^{+-},
  \qquad
  f_{xc}^{+-}(\mathbf r) = \frac{B_{xc}(\mathbf r)}{m(\mathbf r)} .
\end{equation*}
The Kohn--Sham part $\chi_0^{+-}$ is a sum over pairs with one member in each
spin channel, the majority state at $\mathbf k$ and the minority at
$\mathbf k+\mathbf q$; its own spectral weight is the \textbf{Stoner
continuum} of independent spin flips, which costs at least the exchange
splitting. The kernel pulls a mode out from underneath that continuum, and that
mode is the magnon.

\textbf{The kernel is not a second derivative and this is why it works for a
GGA.} Rotate the whole ground state rigidly by an angle $\theta$: then
$\delta\mathbf m_\perp = \theta\,\mathbf m$ and
$\delta\mathbf B_{xc,\perp} = \theta\,\mathbf B_{xc}$, so the transverse
response of the exchange--correlation field to the magnetization is
$B_{xc}/m$ --- an exact consequence of rotational invariance of
$E_{xc}[n,|m|]$, whatever the functional. It also means the Goldstone theorem
holds \emph{by construction}, so the residual below tests the assembly and the
truncations rather than the physics of $f_{xc}$.

\begin{pycode}
import numpy as np
from defumat import Calculator

calc = Calculator.from_file("ni-fcc-magnon.in")

# one wavevector: the spectral function, and the Stoner continuum under it
chi = calc.get_spin_susceptibility((0.0, 0.0, 0.0), np.linspace(0.0, 0.06, 9),
                                   nbnd=30, ecut_response=60.0)
chi.goldstone              # how far X_0 B_xc = m is from holding: the error bar
chi.enhancement            # lambda_max(X_0 F) at omega = 0; one at q = 0
chi.magnon                 # the pole, in Ry, or None where there is none
chi.magnon_mev
chi.spectral_function      # -Im chi_00 / pi
chi.kohn_sham_spectral_function
chi.method                 # "crossing", "peak" or "unstable"

# a path: one fixed-density run serves every q
disp = calc.get_magnon_dispersion([(0, x, x) for x in (0.0, 0.25, 0.5)],
                                  np.linspace(0.0, 0.06, 9),
                                  nbnd=30, ecut_response=60.0,
                                  goldstone_correction=True)
disp.energies_mev          # (nq,), nan where the grid holds no pole
disp.enhancements          # > 1 is an instability, not a magnon
disp.kernel_scale
\end{pycode}

The functional entry points are
\code{defumat.workflows.run\_spin\_susceptibility} and
\code{run\_magnon\_dispersion}.

\textbf{$\mathbf q$ must be a difference of two k-points of the run's own
grid.} That is where the states at $\mathbf k+\mathbf q$ come from --- there is
no second diagonalisation --- and it is Elk's requirement that \code{vecql} be
commensurate with the mesh, reached by construction. What survives of the
$\mathbf k+\mathbf q$ problem is an umklapp shift of one gather index, and
$\chi_0(\mathbf q + \mathbf G) = \chi_0(\mathbf q)$ under a relabelling of the
matrix's own $\mathbf G$ index holds to $1.3\times10^{-17}$, which is what says
that shift is right.

\textbf{The pole is found as a root, not as a peak.} Below the Stoner onset
$\lambda_{\max}(\chi_0(\omega) f_{xc})$ is real and rises monotonically, so the
magnon is the root of $\lambda_{\max} = 1$ --- no broadening enters it, a
handful of frequencies suffice, and it reports \emph{below zero} where the
state is unstable instead of returning the frequency grid's own edge.

\textbf{The Goldstone residual is the number to read first.} A global spin
rotation costs no energy, so $\chi_0(\mathbf q{=}0,\omega{=}0)\,B_{xc} = m$
exactly --- but only with a complete band set \emph{and} a complete
$\mathbf G$-set, so what \code{goldstone} reports is a double truncation.
\textbf{Which axis binds depends on the element.} On hydrogen the band count
does it: $8.0 \to 0.50$ per cent over $n_{\text{bnd}} = 12 \to 140$. On fcc
nickel the band count does nothing at all (10.02, 10.13, 10.22 per cent at 30,
60, 100 bands) and only \code{ecut\_response} moves it: $10.0 \to 6.1 \to 2.0$
per cent at 12, 30, 60~Ry, because $B_{xc}$ of a 3d shell has structure a small
sphere cannot represent.

\textbf{Against Elk's committed reference, nickel's magnon agrees to 0.8\%.}
Elk ships an \code{Ni-magnetic-response} example --- fcc nickel at
$\mathbf q = (0.1,0.1,0)$ on a $10^3$ shifted grid --- whose pole sits at
117.92~meV. Its input applies a small magnetic field, which is worth 1.99~meV
of Zeeman gap, so its field-free magnon is 115.93~meV against \textbf{115.04}
here on the same cell, grid and wavevector --- and 123.08 with the Goldstone
correction applied, which moves it \emph{away}, so neither is chosen for
agreeing. That is across an all-electron LAPW basis versus a norm-conserving
plane-wave one, and across a 25~Ry response sphere versus 60. The moment agrees
on the same run (0.6178 against 0.6152 $\mu_B$), and so does the sign of the
static $\chi$. \textbf{That 0.7\% of kernel is worth 7\% of magnon} is the
number to take away: it is how steeply $\lambda(\omega)$ crosses one, and so how
far the Goldstone residual has to fall before a magnon energy is worth a second
digit.

\textbf{An enhancement above one is an instability, and that is a result.} It
says the collinear state is not a minimum at that wavevector --- the transverse
susceptibility has changed sign --- so the ferromagnet is not the ground state
there. Nickel's enhancement \emph{falls} away from $\mathbf q = 0$, which is
what a stable ferromagnet looks like. Most simple models do not: a half-filled
band on any lattice prefers antiferromagnetic order, so an fcc hydrogen crystal
comes out soft at finite $\mathbf q$ at every spacing tried, and the calculation
says so rather than producing a plausible dispersion anyway.

\textbf{Do not compare that against a spiral $E(q)$ without checking what the
spiral converged to.} It is the obvious cross-check and it does not work: on fcc
hydrogen the 90-degree spiral at $q_3 = \tfrac12$ converges \emph{off} the
magnetic branch, so its energy gain is demagnetization and not a spin wave; a
15-degree cone collapses the same way, because nothing holds it. (The collapsed
$|m|$ is not a number to quote --- two runs of the same input give 0.0001 and
0.0435, which is where a wandering SCF stopped rather than a property of the
state.)

\textbf{A constraint does not rescue it on this cell, and the reason is the cell.}
A spiral SCF accepts \code{constrained\_magnetization} and a magnetic field ---
neither is refused --- and the quantity a constraint acts on is the
\textbf{rotated-frame} magnetization, which is the right object for a helix. But
\code{'fsm'} does not converge here, and it does not converge at $q = 0$ either,
where there is no rotated frame at all: fcc hydrogen at this spacing is a
\emph{marginal} magnet, so $m(B)$ is nearly a step --- 0.057~Ry of field takes the
moment from 0.027 to 0.719 --- with very little in between for a constraint to
hold. The moment lands saturated, the inner SCF then stops converging under the
field that put it there, and since the field steps only on self-consistent pairs it
freezes. Two numbers say that cell was the wrong choice before the run: its bare
moment is 0.0435 and the target was 0.0273, so the constraint asked for a moment
\emph{below} the cell's own while the only other stable point was saturated in the
opposite direction. Use a robust magnet, and give \code{electron\_maxstep} room:
it is \textbf{shared} between the inner SCF and the outer field loop.

\textbf{What to read when one does not converge.} \code{'fsm'} drives a
\emph{feedback field} rather than a penalty, so its \code{constraint\_energy} is
0 by construction and tells you nothing. \code{SCFResult.constraint\_residual} is
$\mathbf m - \mathbf m_{\text{target}}$, signed and per component, with the same
per iteration in \code{history}; a residual that has changed \emph{sign} is an
overshoot rather than a slow approach. And \code{converged = False} on such a run
is often not a density failure --- \code{accuracy} can sit below
\code{conv\_thr} while the constraint is unmet, which is a separate warning with
its own advice. Elk's own magnon-spiral example runs \code{fsmtype = -1} with a
large field along \code{momfix} instead: the \emph{direction} fixed hard rather
than a small magnitude pinned. That scheme is not implemented here.

\code{goldstone\_correction=True} scales the kernel by
$1/\lambda_{\max}(\mathbf q{=}0)$, the same factor at every $\mathbf q$, so the
uniform mode sits at exactly zero. It is what the magnon literature does, it is
off by default, and it is honest only if the residual error is
$\mathbf q$-independent --- check that by running the dispersion at two
$(n_{\text{bnd}},\ \code{ecut\_response})$ settings.

\textbf{An ultrasoft magnet needs one more term in every matrix element}, and
without it there is no magnon at all rather than a slightly wrong one.
$\langle\psi_{n\mathbf k\uparrow}|e^{-i(\mathbf q+\mathbf G)\cdot\mathbf r}
|\psi_{m\,\mathbf k+\mathbf q\downarrow}\rangle$ is the plane-wave overlap of
the pseudo states only when all of the charge is in $|\psi|^2$; with an
ultrasoft dataset what is missing is

\begin{equation*}
\sum_a\sum_{ij} Q^a_{ij}(\mathbf q+\mathbf G)\,
\langle\psi_{n\mathbf k}|\beta^{\mathbf k}_i\rangle
\langle\beta^{\mathbf k+\mathbf q}_j|\psi_{m\,\mathbf k+\mathbf q}\rangle,
\end{equation*}

the augmentation charge at the same wavevector the exponent carries. Nothing is
derived for it: $q^a_{ij}(\mathbf b)$ at an arbitrary wavevector is the
primitive the ultrasoft Berry phase already uses, called once per $\mathbf G$ of
the response sphere (13~ms each, so a sphere of 59 vectors is under a second).

The number is the Goldstone identity on fcc nickel with an ultrasoft dataset:
\textbf{0.984} without the term against \textbf{0.071} with it, at the same
sphere and band count. The bare figure is not an independent error bar and the
guide will not pretend it is --- the identity's $\mathbf m$ is read off the
dense-grid density, which already carries \code{addusdens}'s augmentation, so
0.984 is the augmentation's \emph{share of the moment}. What measures the
assembly is the augmented number, and it is truncation rather than a missing
term: over \code{ecut\_response} of 6, 8, 12, 16, 24 and 48~Ry it reads 0.1063,
0.0706, 0.0992, 0.0987, 0.0854 and \textbf{0.0128}. \emph{Not} monotone --- the 8~Ry
point is the default and sits below the next two, reproducibly --- so read the
identity at the setting a dispersion is to be run at rather than trusting a
trend.

\begin{refuses}
\textbf{Refuses.} \code{nspin != 2}: without a collinear magnetization the
transverse block does not decouple and $B_{xc}/m$ does not exist, and a
\textbf{noncollinear} ground state needs the full $4\times4$ spin-density
response, which is a different object. \textbf{PAW}, and no longer
ultrasoft. $\langle u_i|e^{-i(\mathbf q+\mathbf G)\cdot\mathbf r}|u_j\rangle$
is not the plane-wave overlap once the charge is not all in $|\psi|^2$, and the
missing augmentation charge $Q_{ij}(\mathbf q+\mathbf G)$ is now carried --- it
is the primitive the ultrasoft Berry phase already uses, called once per
$\mathbf G$ of the response sphere. What PAW is refused for is a \emph{second}
term that is not this one: its exchange-correlation field has a one-centre part
on the spheres, where the kernel $B_{xc}/m$ is built from the grid field alone,
so a PAW magnon would carry the right matrix element and the wrong
enhancement. A \textbf{spin
spiral}, whose rotating frame is not the one the kernel was derived in. A
\textbf{Hubbard $U$}, which responds to the magnetization too. An
\textbf{applied magnetic field or a constrained moment}, because the rotation
argument assumes the energy is invariant under a global spin rotation and a
field is exactly what breaks it --- a magnon in a field is gapped by the Zeeman
energy, which this construction does not carry. A \textbf{potential-only
meta-GGA}, there being no $E_{xc}$ for that argument to be about. And a
\textbf{symmetry-reduced k-set}: $\chi_0(\mathbf G,\mathbf G')$ is a matrix in
two $\mathbf G$ indices that nothing here symmetrises, and the
$\mathbf k+\mathbf q$ fold needs the grid closed under translation by
$\mathbf q$ in any case.
\end{refuses}

\section{Optical spectra and excitons: TDDFT}

The response above is a \emph{static} one, solved orbital by orbital. An
absorption spectrum needs the frequency axis and needs the susceptibility as a
\textbf{matrix} in reciprocal lattice vectors, so this is the one place where a
sum over states earns its keep. \code{run\_absorption} builds
\begin{equation*}
\chi^0_{\mathbf{GG}'}(\omega) = \frac{1}{\Omega}\sum_{\mathbf k}\sum_{ij}
  \frac{(f_i - f_j)\,\rho^{ij}_{\mathbf G}\,\rho^{ij\,*}_{\mathbf G'}}
       {\epsilon_i - \epsilon_j + \omega + i\eta},
\qquad
\rho^{ij}_{\mathbf G} = \langle u_i|e^{-i\mathbf G\cdot\mathbf r}|u_j\rangle
\end{equation*}
and solves the Dyson equation with an exchange-correlation kernel:
\begin{equation*}
\epsilon^{-1} = 1 + X\big(1 - X - \tilde f_{\mathrm{xc}} X\big)^{-1},
\qquad X = v^{1/2}\chi^0 v^{1/2},
\qquad \tilde f_{\mathrm{xc}} = v^{-1/2} f_{\mathrm{xc}} v^{-1/2}
\end{equation*}
The \textbf{bootstrap} kernel of Sharma, Dewhurst, Sanna and Gross
(\emph{PRL} \textbf{107}, 186401 (2011); Elk's \code{fxctype = 210}) is that
equation's other half, so the two are solved together to self-consistency:
\begin{equation*}
f^{\mathrm{BS}}_{\mathrm{xc}}
  = -\,\frac{\epsilon^{-1}(\mathbf q, 0)\,v(\mathbf q)}{\epsilon_0(\mathbf q,0)-1},
\qquad \epsilon_0 = 1 - v\chi^0
\end{equation*}
It is parameter-free and diverges as $1/q^2$, which is what binds an
electron--hole pair; ALDA's kernel is finite at $\mathbf q = 0$ where $v$
diverges, so its head and wings vanish \emph{identically} and no amount of it
can bind one.

\textbf{Quantum ESPRESSO has no counterpart.} \code{TDDFPT/} is a
Liouville--Lanczos solver with RPA and ALDA; it has no bootstrap kernel and no
Dyson equation in $\mathbf G$ space. The reference is Elk.

\begin{pycode}
from defumat.workflows import run_absorption
import numpy as np

omega = np.arange(0.0, 0.6, 0.005)          # Ry
spectrum = run_absorption(system, pseudos, scf.density, omega,
                          kernel="bootstrap",   # or rpa / alda / lrc / bootstrap-1
                          nbnd=60, ecut_response=8.0,
                          broadening=0.012, scissor=0.05)

print(spectrum.absorption)                  # Im eps_M, isotropic average
print(spectrum.frequencies_ev)              # the same grid in eV
print(spectrum.alpha, spectrum.iterations)  # the LRC-equivalent head, and the
                                            # bootstrap's fixed-point count
print(spectrum.static_residual)             # how far the band truncation reaches
\end{pycode}

\code{chi\_0} is the whole cost and the kernel is nearly free, so comparing
kernels means building it once: \code{independent\_response} returns it and
\code{solve\_dyson} screens it, which is what \code{run\_absorption} does
internally and what the notebook shows.

Three knobs here are this code's own. \code{ecut\_response} (Elk's
\code{gmaxrf}) selects the $\mathbf G$-set the matrix is built over and has its
own convergence: too small and the local-field effect is missing, and the cost
is quadratic in it. \code{scissor} is \emph{part of the method} rather than a
convenience --- the paper's Eq.~(3) replaces $\chi^0$ by a gap-corrected model
response, and the velocity matrix elements are renormalised by
$e_{ij}/(e_{ij}\mp\Delta)$ along with it. And \code{static\_residual} is the
diagnostic for the one error this phase cannot refuse: how many empty states
the sum ran over. It is $\epsilon_M(0)$ from this run, in RPA, minus the same
quantity from the Sternheimer solve, which is band-\emph{complete}; it is
kernel-matched on purpose, since differencing a bootstrap number against a
Sternheimer one that contains $f_{\mathrm{xc}}$ would measure the kernel
instead.

\textbf{A trap that is invisible in the answer:} the macroscopic dielectric
function is the inverse of the $3\times3$ \emph{head} of $\epsilon^{-1}$, not
the head of the inverse of the whole matrix. The second is the microscopic
$\epsilon$, whose head in RPA is exactly the no-local-field result; it is
smooth, positive, has the right peaks and is 9\% too large on silicon. Both are
returned (\code{epsilon} and \code{epsilon\_no\_local\_fields}), because the gap
between them \emph{is} the local-field effect.

Validated by an identity rather than against another code's spectrum: the same
$\epsilon_M(0)$ is reached by this sum over states plus a Dyson inversion and by
the projected conjugate-gradient solve of \code{dielectric\_tensor}, which
shares no machinery with it and never sees an empty state. On silicon the two
agree to \textbf{1.3e-2} on a constant of 22 --- and that residue is the band
truncation, which is why it is reported rather than tuned away. The
\code{screening = "hartree"} switch on \code{dielectric\_tensor} exists for
this comparison: the identity holds only when both sides use the same kernel.

\begin{refuses}
\textbf{Refuses.} Finite $\mathbf q$ (the perturbed states would live at
$\mathbf k + \mathbf q$); ultrasoft and PAW (every matrix element gains the
augmentation charge $Q_{ij}(\mathbf G)$); metals (the $f_i - f_j$ weight kills
the $i=j$ term, so the intraband Drude response is absent entirely --- Elk's
\code{tddftlr} does not add it either); \code{nspin} $\neq$ 1, noncollinear
magnetism, spin--orbit coupling, spin spirals and a Hubbard $U$; a
symmetry-reduced $k$-set, because symmetrising $\chi^0_{\mathbf{GG}'}$ rotates
\emph{two} $\mathbf G$ indices at once and the rank-$n$ symmetriser is
Cartesian; the \code{alda} kernel with a gradient-corrected functional, since
ALDA's kernel is pointwise in the density and a GGA's is not; and
\textbf{non-convergence of the bootstrap fixed point is an error}, as it is in
\code{tddftlr.f90}, because a spectrum from a kernel still in motion is the
spectrum of nothing.
\end{refuses}

\subsection{The conductivity tensor, the Kerr angle and the anomalous Hall effect}

The whole complex tensor, interband plus a Drude term:
\begin{equation*}
\sigma_{ij}(\omega) = \frac{i}{\Omega}\sum_{\mathbf k}\sum_{n,m}
  \frac{W_n\,(1-f_m)}{\epsilon_{mn}}
  \left[\frac{v^i_{nm}v^j_{mn}}{\omega-\epsilon_{mn}+i\eta}
      + \frac{\overline{v^i_{nm}v^j_{mn}}}{\omega+\epsilon_{mn}+i\eta}\right]
  \;+\; \frac{\omega_{p,ij}^2}{4\pi(\gamma - i\omega)}
\end{equation*}
with $v = \partial H/\partial\mathbf k$ from one \code{jvp} of $H(\mathbf k)$ at
a frozen sphere. \emph{That} is the one thing here not transcribed from
\code{dielectric.f90}, and it is not a refinement: \code{epsilon.x} accumulates
$\langle\psi_1|\mathbf G|\psi_2\rangle$, and $[H,\mathbf r]\neq\mathbf p$ when
the pseudopotential is nonlocal.

What is new beside \code{run\_absorption} is the \textbf{antisymmetric} part,
$\sigma_{xy}$ --- the response that rotates light reflected off a magnet, and
whose static limit is the intrinsic anomalous Hall conductivity. Time reversal
forces it to zero, so a nonmagnetic crystal has none however strong its
spin--orbit coupling; a \emph{collinear or coplanar} magnet without spin--orbit
coupling has none either, because a spin rotation composed with conjugation
survives as an antiunitary symmetry. (A \emph{noncoplanar} magnet does have
one with no spin--orbit coupling at all --- the topological Hall effect --- so
the usual pairing is a route and not a theorem.)

\begin{pycode}
calc = Calculator.from_file("ni-soc.in", pseudo_dir="pseudo")

sigma = calc.get_optical_conductivity(nbnd=36, window=0.6, nw=200)
print(sigma.sigma_s_per_cm.shape)   # (nw, 3, 3) in S/cm
print(sigma.kerr)                   # complex Kerr angle, degrees, per frequency
print(sigma.hall_conductivity)      # the static antisymmetric part, S/cm
print(sigma.plasma_ev)              # hbar omega_p, eV, a (3, 3) tensor
print(sigma.dielectric)             # 1 + 4 pi i sigma / omega
print(sigma.fsum, sigma.band_cut_gap)   # two of the three diagnostics
print(sigma.degenerate_pairs)       # the third: pairs the guard dropped, usually 0
\end{pycode}

There is no reference output for any of it, so the validation is internal and
analytic. The \emph{symmetric} part reproduces the independent-particle
dielectric function of the TDDFT stack to \textbf{3.6e-14}, which is a
different assembly (a Dyson solve over a response sphere) reaching \code{ph.x}
through \code{dielectric\_tensor}. The \textbf{f-sum rule} is the absolute
anchor: $\int_0^\infty\mathrm{Re}\,\sigma_{aa}\,d\omega$ must equal
$(\pi/2\Omega)\sum W_n\langle n|\partial^2H/\partial k_a^2|n\rangle$, and the
familiar $\pi n_e/2$ is only the \emph{local} Hamiltonian's special case ---
silicon's diamagnetic weight is measurably \textbf{0.943} and not one. Measured
against a central difference of the velocity operator, the ratio is 1.258 on a
$4^3$ grid, 1.045 on $6^3$ and \textbf{0.990} on $8^3$: it converges in the
\emph{k-grid}, because what separates the two is $\sum_k w_k\,
\partial^2\epsilon_n/\partial k^2$, which vanishes over the zone by periodicity
and not on a coarse mesh. Aluminium's plasma frequency comes out at
\textbf{11.89 eV} against a free-electron 16.27, which is the zone-boundary gaps
removing Fermi surface, and its tensor is isotropic to \textbf{2.5e-4} eV with
nothing imposing that.

\begin{refuses}
\textbf{Refuses.} \textbf{Ultrasoft and PAW run}, and did not until P99: the
current operator of a generalised eigenproblem is $\langle n|\partial_a H -
\varepsilon_m \partial_a S|m\rangle + (\varepsilon_m-\varepsilon_n)K^a_{nm}$
and the last piece --- the augmentation dipole --- is neither tangent of the
\code{jvp} that gives the first two. It is the Kubo curvature's own term and
its numbers are there; what this module adds is that the two assemblies are the
\emph{same} object, agreeing to \textbf{5.1e-15} on norm-conserving and on
ultrasoft AlAs, which is what says an index order transposed in either would
have shown. A \textbf{spin spiral}, whose two spinor
components live on different spheres. A \textbf{symmetry-reduced $k$-set},
because the antisymmetric part is an axial vector --- the escape is the whole
unshifted grid, which is closed under the point group. And
\textbf{\code{nspin} = 2}, whose two channels are two band structures whose
conductivities add; a spinor run is the regime a magneto-optical spectrum wants
anyway, since $\sigma_{xy}$ needs spin--orbit coupling and collinear spin has
none.

\textbf{Three things are reported rather than refused, and all three must be
read.} \code{band\_cut\_gap} is the gap between the last band summed and the
first dropped: stopping \emph{inside} a degenerate multiplet keeps some members
and drops others, and silicon's antisymmetric $\sigma$ --- zero by time
reversal --- comes out at 4.0e-13 when the cut falls at a real gap and at
1.0e-5 when it does not. \code{degenerate\_pairs} counts the occupied/empty
pairs that came back closer together than \code{degeneracy\_tol} (1e-5~Ry,
what an SCF converges an \emph{empty} band to) and were therefore dropped.
Zero is the ordinary case. When it is not zero it usually still does not
matter --- on the smeared frequency route the two orderings of such a pair
cancel analytically, because what survives the $1/\epsilon_{mn}$ is
$w_k[f(\epsilon_n)-f(\epsilon_m)]$ and that is itself linear in the gap ---
and the code warns exactly where that cancellation fails: \code{method =
"curvature"}, whose $1/\epsilon_{mn}^2$ leaves a $1/g$ divergence, and a
fixed-occupation run, where $f$ is not a function of energy at all. And the
\textbf{anomalous Hall conductivity of a metal is mesh-hungry}: the static
Berry-curvature integral is concentrated on near-degeneracies at the Fermi
surface, so it is quoted with its $k$-convergence and not on its own.
\end{refuses}

\subsection{The shift current: the bulk photovoltaic effect}

Illuminate a crystal with no inversion centre and it carries a direct current,
with no junction and no built-in field. The intrinsic part of that is the
\emph{shift current}: the photoexcited electron is born displaced, and what the
current counts is the real-space shift between the valence and conduction
Wannier centres. It is a second-order optical response,
$J^a = 2\sigma^{abc}(0;\omega,-\omega)E_b(\omega)E_c(-\omega)$, with
\begin{equation*}
\sigma^{abc}(\omega) = -\frac{i\pi e^3}{4\hbar^2}\int\![d\mathbf k]
  \sum_{n,m} f_{nm}\left(I^{abc}_{mn}+I^{acb}_{mn}\right)
  \left[\delta(\omega_{mn}-\omega)+\delta(\omega_{nm}-\omega)\right],
  \qquad I^{abc}_{mn} = r^b_{mn}\,r^{c;a}_{nm} .
\end{equation*}
The whole difficulty is $r^{c;a}_{nm}$, the \emph{generalised derivative}
$\partial_a r^c_{nm} - i(A^a_{nn}-A^a_{mm})r^c_{nm}$ --- neither piece of which
is separately computable, since the diagonal Berry connection is gauge dependent
and is not a function of $H$ at one $\mathbf k$. The escape is the Aversa--Sipe
sum rule, which writes the covariant object entirely in gauge-free quantities
living at a single $\mathbf k$:
\begin{equation*}
r^{c;a}_{nm} = \frac{i}{\omega_{nm}}\left[
  \frac{v^c_{nm}\Delta^a_{nm} + v^a_{nm}\Delta^c_{nm}}{\omega_{nm}}
  - w^{ca}_{nm}
  + \sum_{p\neq n,m}\left(\frac{v^c_{np}v^a_{pm}}{\omega_{pm}}
                        - \frac{v^a_{np}v^c_{pm}}{\omega_{np}}\right)\right],
\end{equation*}
with $v^a_{nm} = \langle n|\partial_a H|m\rangle$ and
$w^{ab}_{nm} = \langle n|\partial_a\partial_b H|m\rangle$. Both are derivatives
of code that already exists: $v$ is one \code{jvp} of $H(\mathbf k)$ at a frozen
sphere and $w$ is that \code{jvp} differentiated by a second one. Nothing is
transcribed. The interband dipole $r^a_{nm} = -i\,v^a_{nm}/\omega_{nm}$ is
\emph{exact} for an eigenstate of the full $H(\mathbf k)$, which is why a
plane-wave code needs no counterpart to Wannier90's position-operator matrices.

\begin{pycode}
calc = Calculator.from_file("alas-raman.in", pseudo_dir="pseudo")

sc = calc.get_shift_current(nbnd=14, window=0.9, nw=180)
print(sc.sigma.shape)            # (nw, 3, 3, 3) in A/V^2, real
print(sc.component(0, 1, 2))     # sigma^xyz(omega), -43m's only free component
print(sc.frequencies_ev)         # the photon energy axis in eV
print(sc.truncation)             # the band-truncation diagnostic: read it
print(sc.band_cut_gap)           # where the band set was cut, in Ry
\end{pycode}

There is no reference binary --- \code{pw.x} computes no photocurrent of any
kind, and Elk's \code{nonlinopt.f90} is second-harmonic generation, a different
response --- so the validation is internal and it is three statements. The
\textbf{sum rule against a parallel-transport finite difference} of the dipole,
which shares nothing with it: no second derivative, no intermediate sum, only
the phases of neighbouring $\mathbf k$-points fixed to one another. That
comparison is decisive in exactly one form, because the sum rule is an
\emph{identity over a complete basis}: held on a frozen sphere of 158 plane
waves, AlAs agrees to \textbf{1.8e-4} at 158 bands against 4.3e-2 at 120 and
6.0e-2 at 20. A residue that falls off a cliff at completeness is a correct
assembly with a truncated sum; one that plateaus is a bug. Then the
\textbf{crystal class}: AlAs comes out exactly $\bar 43m$ on a \code{nosym}
grid, only the components with all three labels different surviving, while
centrosymmetric silicon gives zero from the same machinery. And the
\textbf{absorption edge}: $\sigma^{xyz}$ is 6.9e-25 A/V$^2$ at 2 eV against a
peak of 1.1e-4 near 4.4 eV, which is the one check no index bookkeeping can pass
by accident --- no absorption, no photocurrent.

\textbf{All three are blind to an overall constant}, and this document says so
rather than leaving it to be discovered: a $\sigma$ wrong by a factor of two is
still exactly $\bar 43m$, still zero on silicon and still zero below the gap.
Two further checks close that gap. The \textbf{spin sum} --- the same cell run
\code{nspin} = 1 and as a two-component spinor gives the same tensor to
\textbf{4.6e-9}, peak ratio 1.000000. And the \textbf{overall scale against the
published literature}, where the \emph{ordering} is what makes the comparison
sharp rather than merely consistent. AlAs's first peak converges to
\textbf{35.0} $\mu$A/V$^2$ at 4.17 eV (33.9 at \code{ecutwfc} = 16, then 33.7,
34.8 and 35.0 at 22 with 14, 22 and 30 bands), where calculations across the
fourteen III--V and II--VI zincblende semiconductors span 14 $\mu$A/V$^2$ (CdSe,
smallest) to 83 (AlSb, largest) and find the aluminium compounds the strongest of
the family. A factor of two either way breaks that ordering rather than moving
the number: 17 would put AlAs at the bottom of the family, below the cadmium
chalcogenides, and 70 would put it beside AlSb, the heaviest and narrowest-gap
member. The number to compare is the peak below about 6 eV, which is the range
published spectra cover. That sweep also showed \code{truncation} to be
\emph{predictive} rather than merely indicative: it falls 3.5\% to 1.0\% over the
band sweep while the value moves 3.9\%.

\textbf{The convention is declared, because the literature has two.} What is
implemented is $j^a = 2\sigma^{abc}(0;\omega,-\omega)\,
\mathrm{Re}[\mathcal E_b(\omega)\mathcal E_c(-\omega)]$, the definition of the
paper the formulae above come from. Some authors write
$J = \sigma\mathcal E\mathcal E$ and report numbers twice as large for the same
physics, so a comparison against a paper that does not state its normalisation
is worth nothing to better than a factor of two.

\begin{refuses}
\textbf{Refuses.} \textbf{Ultrasoft and PAW}, one order further out than the
Kubo curvature, the optical conductivity and second-harmonic generation, none
of which refuse them any more. The
dipole carries $\partial S/\partial k$, its derivative carries
$\partial^2 S/\partial k_a\partial k_b$, and the augmentation dipole
$\langle\psi_n|T^\dagger\partial_k T|\psi_m\rangle$ that belongs with them is
here too. What is not here is \emph{that term's own $k$-derivative}, which is
built from $\langle\beta|\psi\rangle$ and $\partial\beta/\partial k_a$ and so
wants $\partial^2\beta/\partial k_a\partial k_b$ --- a second derivative of a
bare projector, which nothing in this package forms. A
\textbf{spin spiral}. A \textbf{symmetry-reduced $k$-set}, because
$\sigma^{abc}$ is a polar rank-3 tensor and a wedge sum is not the cell's until
it is averaged over the point group; the escape is the whole unshifted grid. A
\textbf{metal}, whose partially filled bands contribute pairs with vanishing
$\omega_{nm}$ where every denominator above diverges --- the Fermi-surface terms
that replace them are a different assembly. \textbf{DFT+$U$}, whose Hubbard term
is built from $\phi_U(\mathbf k + \mathbf G)$ and so carries a velocity and a
second velocity of its own. And \textbf{\code{nspin} = 2}. A \emph{spinor} run
is supported and tested.

\textbf{The band count is reported rather than refused, and it is the largest
approximation here.} The intermediate sum over $p$ converges only as it
approaches completeness, so a production band count carries a few per cent that
no symmetry check sees --- \code{truncation} is the in-run estimate and is taken
over the top \emph{quarter} of the bands, because dropping one band reports
5.3e-5 where the error is 4 per cent. The route that would remove it is written
down in \code{NONLINEAR.md}: the two sums are a Sternheimer solve rather than a
spectral sum, and what stops it being taken is that the operator to invert is
indefinite. \code{band\_cut\_gap} carries the same warning it does for the
conductivity.
\end{refuses}

\subsection{Second-harmonic generation}

Shine light of frequency $\omega$ on a crystal with no inversion centre and some
of it comes back out at $2\omega$. The tensor that says how much is the
second-order optical susceptibility
$P^a(2\omega) = \chi^{abc}(-2\omega;\omega,\omega)\,E_b(\omega)E_c(\omega)$, a
polar rank-3 tensor, symmetric in its last two labels and identically zero in
every centrosymmetric crystal. It is a sum of three physically distinct pieces
--- the pure interband term $\chi_{II}$, the intraband modulation of it
$\eta_{II}$, and the modulation of the transition energy
$\tfrac{i}{2\omega}\sigma_{II}$ --- each of which is a contraction of the
interband dipoles $r^a_{nm} = -i\,v^a_{nm}/\omega_{nm}$ over two resonant
denominators, one at $\omega$ and one at $2\omega$:
\begin{equation*}
\chi^{abc}(-2\omega;\omega,\omega) = \frac{1}{\Omega}\sum_{\mathbf k, m, n}
  \left[\frac{c^{(1)}_{mn}}{\omega_{mn}-\omega+i\eta}
      + \frac{c^{(2)}_{mn}}{\omega_{mn}-2(\omega-i\eta)}\right],
\end{equation*}
where the coefficients $c^{(1,2)}_{mn}$ carry a further sum over an intermediate
state and are built once, before the frequency axis is touched. This is the same
sum over states the shift current uses, contracted differently; it needs no
second derivative of $H$ at all, because that triple sum \emph{is} the sum-rule
expansion of the generalised derivative written out rather than collapsed.

\begin{pycode}
calc = Calculator.from_file("alas-shg.in", pseudo_dir="pseudo")

shg = calc.get_shg(nbnd=22, window=0.6, nw=240, broadening=0.010)
print(shg.chi.shape)                  # (240, 3, 3, 3), complex, in pm/V
print(shg.component(0, 1, 2)[0])       # chi^xyz at the lowest frequency
print(shg.d_coefficient(0, 1, 2)[0])   # d = chi/2, what the literature quotes
print(shg.frequencies_ev[:3])          # the FUNDAMENTAL photon energy, in eV
print(shg.truncation)                  # the band-truncation diagnostic
print(shg.chi_ii, shg.eta_ii, shg.sigma_ii)   # the three parts, kept apart
\end{pycode}

\textbf{There is a real reference here}, which is unusual in this corner of the
package: Elk's \code{nonlinopt.f90} (task 125) computes the same tensor, and
\code{tests/data/elk/} carries its output for the same AlAs crystal, the three
parts separately as well as their sum. \code{pw.x} has none --- its
\code{el\_opt.f90} is the \emph{static} electro-optic response, and the branch
that reaches even that is the one this document's Raman section records as
broken in the vendored 7.5.

What the comparison is \emph{not} is digit-for-digit, and saying so is part of
quoting it: an all-electron LAPW basis is not a norm-conserving pseudopotential
one, and a second-order susceptibility carries two energy denominators, so a gap
difference is not a scale factor. On the same crystal and the same $6^3$ mesh:
the \textbf{resonance position} agrees to \textbf{0.5\%} (2.152 eV against
2.163), the \textbf{peak height} to \textbf{7\%} (27.52 a.u.\ against 25.70),
and the \textbf{static value} to \textbf{11\%} ($-3.10$ a.u.\ against $-3.50$).
The last is the one that moves with the basis, and it is shown converged:
\code{ecutwfc} 10, 30 and 45 give $-2.71$, $-3.10$ and $-3.13$, so the residue is
the pseudopotential rather than the cutoff.

Four internal checks close what the reference cannot. \textbf{$\bar 43m$}: AlAs
comes out exactly zincblende on a \code{nosym} grid with nothing imposing it.
\textbf{Silicon}: an inversion centre forbids every component, and the
centrosymmetric control runs through the same machinery. \textbf{The two-photon
edge}: $\mathrm{Im}\,\chi$ turns on at \emph{half} the smallest direct gap, not
at the gap, because the second channel's denominator is $\omega_{mn}-2\omega$ ---
the only check here blind to both the scale and the symmetry, and a lost factor
of two there moves the absorption edge by an octave. And \textbf{the spinor},
which is the scale's own anchor: the same cell run \code{nspin} = 1 and as a
two-component spinor must give the same tensor, which is what catches a factor of
two in the spin sum.

\textbf{The convention is declared, because the literature has two.} What is
computed is $\chi^{(2)}$, defined by $P^a = \chi^{abc}E_bE_c$; nonlinear-optics
papers more often quote $d^{abc} = \chi^{abc}/2$, so a number compared against a
published $d$ without halving it is wrong by exactly two and nothing about its
shape, symmetry or zeros would say so. \code{d\_coefficient} returns the halved
form.

\textbf{The scissors shift is the one knob worth explaining.} An LDA or GGA gap
error does not scale a second-order susceptibility, it moves its resonances, and
\code{scissor} applies the usual rigid shift to the empty states. The dipoles are
built from the \emph{unshifted} gaps --- a scissors correction is a statement
about energies, not about wavefunctions --- which is what Elk reaches by scaling
its momentum matrix elements instead. Validated against Elk's own scissored run:
at $\Delta = 0.05$~Ha the $2\omega$ peak moves \textbf{0.0502}~Ry against a
half-scissor of 0.0500, and its height falls to \textbf{0.60} of the unscissored
value against Elk's 0.58.

\begin{refuses}
\textbf{Refuses.} Four, and they are what a \emph{sum over velocity matrix
elements} refuses rather than what the shift current does: a \textbf{spin
spiral}, a \textbf{symmetry-reduced $k$-set} (a polar rank-3 tensor is not the
cell's until averaged over the point group; the escape is the whole unshifted
grid), a \textbf{metal}, \textbf{DFT+$U$}, and \textbf{\code{nspin} = 2}. A
\emph{spinor} run is supported and tested.

\textbf{Ultrasoft and PAW are not among them, and were until P99 because this
module inherited the shift current's guard whole.} That refusal belongs to the
generalised derivative $r^{c;a}$, which the shift current differentiates and
this assembly does not: the triple sum over $l$ above \emph{is} that
derivative's sum-rule expansion written out, so what an augmented dataset needs
here is the velocity matrix element and nothing beyond it. That element is
$\langle n|\partial_a H - \varepsilon_m\partial_a S|m\rangle + (\varepsilon_m -
\varepsilon_n)K^a_{nm}$, whose last piece is the augmentation dipole of the
Kubo section, and deleting it moves $|\chi^{(2)}|$ on ultrasoft AlAs by \textbf{40.8}~pm/V on a
peak of \textbf{1577.3} --- 2.6\,\%, where the same deletion is worth 0.5\,\% in
the linear conductivity, because $\chi^{(2)}$ carries the velocity matrix
element three times rather than twice.

\textbf{This does not lift the \code{chi\^{}(2)} refusal in the Raman section},
and the distinction matters. That one is about the $2n+1$ route through the
Sternheimer solver, where a field enters only through a source term and the
$\langle u_i|r_k|u_j\rangle$ piece has nothing to build it from; a sum over
states never needed that object. The electro-optic tensor and a truncation-free
static $\chi^{(2)}$ are still refused for the same missing term.

\textbf{The band count is reported rather than refused}, exactly as for the shift
current and for the same reason: the sum over the intermediate state is an
identity only over a complete basis. \code{truncation} is the in-run estimate,
taken over the top quarter of the bands.
\end{refuses}

% =====================================================================
\section{Topological invariants}

Everything is built from one primitive --- the overlap of occupied manifolds at
neighbouring $k$-points,
\begin{equation*}
M^{(\mathbf b)}_{mn} = \langle u_{m\mathbf k}|\hat S|u_{n,\mathbf k+\mathbf b}\rangle
\end{equation*}
because a determinant of overlaps is blind to the unitary mixing a degenerate
eigensolver leaves. The Berry curvature is the phase of a plaquette of them and
the Chern number their lattice sum, which is an \textbf{exact integer on any
mesh}:
\begin{equation*}
F_{12}(\mathbf k) = \ln\big[U_1(\mathbf k)U_2(\mathbf k{+}\hat 1)
  U_1(\mathbf k{+}\hat 2)^{-1}U_2(\mathbf k)^{-1}\big],
\qquad C = \frac{1}{2\pi i}\sum_{\mathbf k} F_{12}(\mathbf k)
\end{equation*}
$Z_2$ has two independent routes --- the flow of Wannier charge centres, and,
when there is an inversion centre, the parity product over the eight TRIM
points:
\begin{equation*}
(-1)^{\nu} = \prod_{i=1}^{8}\ \prod_{n\ \mathrm{occ}}^{\mathrm{Kramers}}
  \xi_{2n}(\Gamma_i)
\end{equation*}

\begin{pycode}
from defumat.workflows.topology import (run_z2, run_z2_3d,
                                         run_berry_curvature)

z2 = run_z2(system, pseudos, scf.density, axis=2)   # Wilson loop by default
print(z2.z2, z2.gap_step)   # read gap_step before believing the integer

chern = run_berry_curvature(system, pseudos, scf.density, shape=(12, 12))
print(chern)                # an exact integer on any mesh

nu = run_z2_3d(system, pseudos, scf.density)   # the four 3D indices
print(nu)                   # (nu0; nu1 nu2 nu3)
\end{pycode}

\begin{refuses}
\textbf{Refuses.} A collinear spin-polarized run (two Hamiltonians, so two
independent sets of invariants and no statement of which is meant); a $Z_2$
without spin--orbit coupling, where every invariant is zero for a reason that
has nothing to do with the band structure; the parity route without an
inversion centre; and a manifold that is not separated from the band above it
(\code{gap\_tol}). A PAW run needs the converged \code{becsum} passed beside
the density --- \code{run\_z2(..., becsum = scf.becsum)}, which
\code{Calculator} does for you --- because \code{ddd\_paw} is a property of the
wavefunctions and cannot be rebuilt from a density; leaving it out is worth
1.03~eV in the eigenvalues, and an invariant computed from those is still an
integer.
\end{refuses}

\textbf{Running both $Z_2$ routes is the check.} Where they disagree the parity
one is the answer, because it has no mesh and the Wilson route does.
\textbf{Two things bite in a plane-wave code and both are silent:} neighbouring
$k$-points do not share a $G$-sphere, so coefficients are aligned by Miller
index; and the wrap at the zone edge is a \emph{shift} of that index, without
which the Chern number comes out smooth and non-integer.

\subsection{The curvature as a map: the Kubo route}

Two constructions of the same quantity, chosen by name:
\begin{equation*}
  \Omega_n^{12}(\mathbf k) = -2\,\mathrm{Im}
  \sum_{m \neq n}
  \frac{A^{1}_{nm} A^{2}_{mn}}{(\varepsilon_n - \varepsilon_m)^2},
  \qquad
  A^{a}_{nm} = \langle \psi_n |\, \partial_{k_a} H - \varepsilon_m\,
  \partial_{k_a} S \,| \psi_m \rangle
  + (\varepsilon_m - \varepsilon_n) K^{a}_{nm}
\end{equation*}
with \code{method="kubo"}, against the lattice field strength of Fukui,
Hatsugai and Suzuki with \code{method="fhs"} (the default). The velocity
operator is one \code{jvp} of $H(\mathbf k)$ at a frozen plane-wave sphere ---
no dense $H(\mathbf k)$ is ever formed --- and the tangent for a crystal
direction $d$ is the reciprocal lattice vector $\mathbf b_d$.

The last term is the one a dataset with an augmentation charge needs, and it is
zero for every other. With $\hat S = T^\dagger T$ the states a Berry phase is
about are $T|\psi\rangle$ rather than $|\psi\rangle$, so the connection splits
as $\langle\Psi_n|\partial_a\Psi_m\rangle = \langle\psi_n|\hat S|\partial_a
\psi_m\rangle + K^a_{nm}$ with
\begin{equation*}
  K^{a}_{nm} = \sum_{ij}\langle\psi_n|\beta_i\rangle\,q_{ij}\,
  \langle\partial_{k_a}\beta_j|\psi_m\rangle
  - i\sum_{ij}\langle\psi_n|\beta_i\rangle\,
  d^{a}_{ij}\,\langle\beta_j|\psi_m\rangle,
  \qquad
  d^{a}_{ij} = \int r_a\,Q_{ij}(\mathbf r)\,\mathrm d\mathbf r,
\end{equation*}
the augmentation dipole \code{adddvepsi\_us} adds to an electric field's
position operator.

\textbf{Written that way $A^a$ is Hermitian and is the generalised velocity
matrix}, which is what \code{defumat.response.conductivity} contracts; the code
here writes the equivalent form with the \emph{outer} band's energy
$\varepsilon_n$ in both factors and the two blocks $K^\dagger$ and $K$ split
between them, because that is what falls out of resolving the identity inside
$-2\,\mathrm{Im}\langle\partial_1 n|\hat S|\partial_2 n\rangle$. The two agree
identically, and the reason they do is the identity below. Only
$\partial_{k_a}H$ and $\partial_{k_a}S$ can be got from
the two tangents of one \code{jvp}, and $K$ is neither: $K^\dagger + K$ is
\emph{exactly} $\langle\psi_n|\partial_{k_a}S|\psi_m\rangle$, the dipole
cancelling between the two and the rest adding, so the two factors of the sum
take different blocks and no rearrangement of $\partial S/\partial k$ supplies
it. That identity holds to $1.2\times10^{-16}$ on ultrasoft silicon,
$2.1\times10^{-16}$ on PAW silicon and $2.9\times10^{-16}$ on
fully-relativistic ultrasoft AlAs --- where the blocks are $q^{SO}_{ij}$ and
$d^{SO}$ and carry a spin index each --- which pins the index order and says
nothing about the dipole; the dipole is pinned against a central difference of the FHS
overlap itself, where at a step small enough that all three $k$-points hold the
same plane-wave sphere the connection comes out at
$1.67\times10^{-4}$ (ultrasoft) and $1.67\times10^{-4}$ (PAW) of a
norm-conserving control's own $1.70\times10^{-4}$, while dropping $K$ leaves
$1.57\times10^{-2}$ --- 94 times larger, and \emph{not moving} as the difference
shrinks.

\begin{pycode}
from defumat.workflows.topology import run_berry_curvature

curvature = run_berry_curvature(
    system, pseudos, density, shape=(4, 4), nocc=4, nbnd=12, method="kubo"
)
curvature.curvature          # (4, 4), Omega(k)
curvature.curvature_by_band  # (4, 4, nocc) -- see the refusal box
curvature.truncation         # what the highest empty band was worth
\end{pycode}

Validated against the FHS flux, which shares no machinery with it: a centred
plaquette shrunk around one k-point of zincblende AlAs converges onto the
pointwise Kubo value to $3.2\times10^{-3}$, and the whole $24\times24$ mesh
agrees plaquette by plaquette to $1.45\times10^{-4}$, improving 65-fold over a
sixfold refinement. The velocity matrix elements themselves reproduce a central
finite difference of $H(\mathbf k)$ to $1.8\times10^{-9}$, and on centrosymmetric
silicon the curvature vanishes \emph{pointwise} to $3.5\times10^{-5}$.

\begin{refuses}
\textbf{Refuses.} \code{method="kubo"} is for the map, never for the integer:
its denominator $(\varepsilon_n-\varepsilon_m)^{-2}$ is singular at a degeneracy
and its Brillouin-zone sum is an ordinary Riemann sum, which converges to an
integer without ever being one --- \code{method="fhs"} is the default and is
exact on any mesh. \textbf{Ultrasoft and PAW run}, and what the
\code{method="kubo"} comparison against \code{method="fhs"} cannot tell you is
whether the augmentation dipole is there: on ultrasoft AlAs it is worth
\textbf{0.34\,\%} of $\Omega$ where the two methods sit \textbf{50\,\%} apart
pointwise at $18\times18$, so that comparison would read the same with the term
deleted, and the numbers that carry it are the overlap's own finite difference
above. The sum over empty states is
truncated at \code{nbnd} and the truncation is \emph{reported} rather than
hidden --- read \code{BerryCurvature.truncation}, or \code{truncation\_abs}
wherever symmetry forces the curvature to vanish, since the ratio is then noise
over noise. \code{curvature\_by\_band} is gauge invariant only for a
non-degenerate band; inside a degenerate multiplet only the sum over the members
is defined, and \code{curvature} (the manifold total) always is.

A mesh point where an occupied and an empty band actually \emph{touch} --- a
Dirac cone, a Weyl node, the closing a topological transition runs through ---
is a genuine singularity and is \textbf{dropped, with a warning naming how
many}. \code{BerryCurvature.singular\_points} carries the count, and a nonzero
one means the mesh sum is not a Chern number of that state set: the curvature
diverges there and a zero was put in its place, which is smooth and plausible
and wrong. Move the mesh off the point, or use \code{method="fhs"}, which
needs no gap in a denominator. A band pair counts as touching below
$10^{-8}$ of the band width \emph{over the whole mesh} --- a fraction rather
than an energy, since a model Hamiltonian carries no unit, and over the whole
mesh because for a two-band model the spectrum at one $\mathbf k$ \emph{is}
the gap.
\end{refuses}

\subsection{The Berry-phase polarization}

The polarization of a crystal is not the dipole of its unit cell --- that
integral depends on where the cell is cut. What is a property of the crystal is
a \emph{phase}: the Berry phase the occupied manifold accumulates as $\mathbf k$
crosses the zone once along a reciprocal lattice vector $\mathbf b_i$
(King-Smith and Vanderbilt, PRB \textbf{47}, 1651 (1993)),
\begin{equation*}
  \Omega\,\mathbf P = \sum_i \frac{\varphi_i}{2\pi}\,\mathbf a_i,
  \qquad
  \varphi_i = \underbrace{\operatorname{Im}\ln \prod_{j}
    \det M^{(\mathbf b_i/N)}(\mathbf k_j)}_{\text{electronic}}
  \;+\; \underbrace{2\pi\sum_a Z^{\mathrm v}_a\, s_{a i}}_{\text{ionic}}
\end{equation*}
with $M$ the same overlap the invariants above are built from and $s_{ai}$ the
atom's crystal coordinate. It is the \emph{determinant} of the overlap that is
kept here where a Wilson loop keeps the eigenvalues, which is the whole
difference between the two calculations.

\textbf{The value is defined modulo a quantum} and the result object carries it.
A single number is not a physical statement: what is meaningful is a
\emph{difference} between two geometries taken on the same branch, which is what
a Born effective charge and a piezoelectric constant are made of.

\begin{pycode}
from defumat.workflows.polarization import run_polarization

p = run_polarization(system, pseudos, scf.density,
                     gdir=2,            # 0-based; pw.x's gdir = 3
                     nppstr=6,          # as pw.x counts it
                     transverse=(2, 2))
print(p.total_phase, p.quantum)         # -0.2412383  (mod 1)
print(p.polarization_si)                # -0.34909 C/m^2, mod 1.44708
print(p.strings.phases)                 # one phase per string, four here
\end{pycode}

\code{Calculator.get\_polarization()} reads \code{lberry}, \code{gdir} and
\code{nppstr} straight off a \code{pw.x} input, so a polarization run transfers
unchanged. \code{nppstr} counts the repeated endpoint the way \code{kp\_strings}
does; the mesh built here carries \code{nppstr - 1} distinct points and closes
the string with the reciprocal lattice vector, which is the same links and one
diagonalisation fewer per string.

\begin{refuses}
\textbf{Refuses.} A \textbf{metal} --- a Berry phase is a property of a gapped
manifold, and a smeared occupation does not say which bands the string carries.
\textbf{\code{nspin = 2}}: two channels are two independent string sets and one
phase each, which is a layout this does not have rather than a missing term.
And a \textbf{spin spiral}, whose two spinor components live on spheres centred
at $\mathbf k \pm \mathbf q/2$, so a string overlap is not a single gather. A
\emph{spinor} run (\code{noncolin}) is supported and has its own \code{pw.x}
reference; its quantum is half the scalar one, because a spinor band holds one
electron.
Ultrasoft and PAW are \emph{not} refused, and both are validated: zincblende
SiC matches \code{pw.x} string by string with an electronic phase of $-0.009696$
against $-0.00970$. It had to be SiC rather than silicon, because a
centrosymmetric crystal's electronic phase is zero whatever $q_{ij}(\mathbf b)$
does and so agrees with zero however wrong the augmentation is.
\end{refuses}

\textbf{Validated three ways, two of which need no other code.} Against
\code{pw.x}'s own \code{lberry} run on AlAs it agrees \emph{string by string} and
on every digit printed (electronic phase 0.00876, total $-0.24124$). The Zak
phase of the SSH model is pinned to $0$ or $\pi$ by symmetry and comes out
exactly that on a mesh of \emph{any} size ($10^{-10}$), and silicon is pinned by
inversion to $0$ or half a quantum, landing on the half ($10^{-8}$). Silicon is
also the second \code{pw.x} reference, and the one that exercises the quantum of
2: both its species carry an even valence, where AlAs's arsenic does not. The strongest is
$\mathrm dP/\mathrm du$ against the Born charges, which come from a Sternheimer
solve and share no machinery with a string of overlaps: AlAs on an
$8\times8\times8$ ground state gives $Z^*(\mathrm{Al}) = +2.148$ against
\code{ph.x}'s $+2.142$, charge-neutral to $5\times10^{-5}$.

\subsection{The magnetoelectric tensor}

How much electric polarization a magnetic field induces,
\begin{equation*}
  \alpha_{ij} = \frac{\partial P_i}{\partial B_j},
\end{equation*}
computed Elk's way (task 390): a full ground state at $B \mp \delta/2$ for each
Cartesian direction, the Berry-phase polarization of each, and a central
difference. The cheap route --- one derivative of the magnetization along an
electric field's response --- needs a noncollinear Sternheimer solve, and an
electric field as a ground-state perturbation, which this package does not have.

\textbf{Two symmetries must both be broken or the answer is zero for a reason
that has nothing to do with the calculation}: time reversal and inversion each
map $\alpha$ to $-\alpha$. Elk's own example is built around the first, and so is
the case here --- a non-magnetic semiconductor has no linear tensor at all, so a
\emph{large} external field is applied to break time reversal by hand and the
derivative taken against a small further change in it. The magnetism is induced
rather than spontaneous, which is what makes a plain gapped semiconductor usable.

\begin{pycode}
from defumat import Calculator

alas = Calculator.from_file("alas-magnetoelectric.in", pseudo_dir)
me = alas.get_magnetoelectric_tensor(delta=0.02, directions=(2,))
print(me.magnitude)        # 4.4376e-06
print(me.largest_step)     # how close a field step came to a branch
\end{pycode}

\begin{refuses}
\textbf{Refuses.} An unconverged ground state, which is what makes the next
limitation visible rather than silent. \textbf{Only the column parallel to the
seeded magnetization is reachable} on the committed cell: a field along the
direction \code{starting\_magnetization} points converges in 12 iterations, and
one with a transverse component does not converge in 150, because rotating the
moment is the slow mode. It is not the tilt --- a pure transverse field fails on
its own. Seed the magnetization along the column you want. A run carrying no
magnetization vector ($n_{\mathrm{spin,mag}} \neq 4$): the Zeeman term then multiplies nothing and every polarization comes out
identical, which looks like a symmetry result. And a field step that moves the
Berry phase by more than a quarter of its quantum, since a difference taken across
a branch is the quantum divided by $\delta$ --- large, smooth and meaningless.
Everything the polarization refuses is refused here too.
\end{refuses}

\textbf{The validation is the spin-orbit null.} Without the coupling a global spin
rotation maps $B$ to $-B$ and leaves every charge observable fixed, so $\alpha$
vanishes identically. On the same cell, same field, same ultrasoft family at the
same valence, with only \code{lspinorb} and the scalar/fully relativistic datasets
changed: $3.43\times10^{-9}$ against the coupled $4.44\times10^{-6}$, a ratio of
\textbf{1293}. Two things had to be right together to get there, and neither helps
alone: \code{chain} (seeding one field run from the other, which breaks the very
$B \to -B$ symmetry the null rests on --- off by default here, unlike Elk) and the
\emph{polarization's} threshold rather than the self-consistent field's.
\textbf{The absolute value is not calibrated against another code}: Elk's
\code{bfieldc} and this package's \code{B\_field} differ in normalisation by about
two orders of magnitude, and Elk's own tensor on this cell is not converged.

% =====================================================================
\section{Effective masses, site angular momenta, nesting, structure factors and
STM images}

Quantities Elk computes and \code{pw.x} mostly does not. The first two cost one
NSCF at a single k-point; the nesting function costs one NSCF on a dense grid,
and nothing after it; the structure factors cost one transform of a density that
already exists, and an STM image one band sum with different weights.

\subsection{The effective mass tensor}

\begin{equation*}
  \left(\frac{1}{m^*}\right)_{ab}
  = \frac{1}{2}\,\frac{\partial^2 \varepsilon_n(\mathbf k)}
                      {\partial k_a \partial k_b},
\end{equation*}
in units of $1/m_e$. The factor of a half is Rydberg atomic units and not a
normalisation: $\hbar^2/2m_e$ is exactly $1$~Ry~bohr$^2$, so a free electron has
$\varepsilon = |\mathbf k|^2$ and the tensor is the identity.

\textbf{The first derivative is analytic and only the second is a difference.}
$\partial \varepsilon_n/\partial k$ is the generalised Hellmann--Feynman
expression built from one \code{jvp} of $H(\mathbf k)$; its second derivative is
not reachable the same way, because an individual band's $|\partial_k \psi_n
\rangle$ is not what the Sternheimer projector produces (it removes the whole
occupied manifold, and the band whose mass is wanted is usually empty). So the
construction is one central difference of an \emph{exact} first derivative ---
six stencil points against Elk's twenty-seven, and no polynomial fit. Elk's
route, differencing eigenvalues, is kept beside it as \code{method="eigenvalue"}
because it shares nothing with the velocity operator.

\textbf{The whole stencil is built on one plane-wave sphere}, the centre's, and
its points are then moved onto it. Which plane waves satisfy
$|\mathbf k + \mathbf G|^2 \le \code{ecutwfc}$ is a step function of
$\mathbf k$, so two points a stencil apart can hold different numbers of them:
at $L$ on two-atom silicon at \code{ecutwfc = 30} the centre holds 754, the
$\pm h$ pair 752 and the $\pm 2h$ pair 744. A difference taken across that step
carries it divided by $h^2$, which is not $O(h^2)$ and which the Richardson step
therefore returns rather than removes; on one sphere what is differenced is one
band of one basis. It is worth \textbf{0.5\,\%} on the eigenvalue route at $L$
at the default stencil.

\begin{pycode}
from defumat import Calculator

calc = Calculator.from_file("si.scf.in")
mass = calc.get_effective_mass((0.0, 0.0, 0.0), nbnd=10)  # crystal coordinates

mass.inverse_mass          # (nbnd, 3, 3) in 1/m_e, cartesian
mass.principal(band=7)     # the three principal masses and their axes
mass.density_of_states_mass(band=7)   # (m_x m_y m_z)^(1/3)
mass.truncation            # the O(h^2) the Richardson step removed
mass.multiplets            # per multiplet, including the degenerate ones
\end{pycode}

On two-atom silicon at $\Gamma$ against the vendored \textbf{Elk} binary
(all-electron LAPW, PBE, $a = 10.26$~bohr) with a PAW dataset here: the
non-degenerate $\Gamma_{1v}$ agrees to \textbf{0.02\%} ($0.8601272$ against
$0.8603044$, i.e.\ $m^* = 1.16262\,m_e$ against $1.16238$) and $\Gamma_{2'c}$ to
0.36\% ($m^* = 0.170440\,m_e$). The two routes here agree with each other to
$1\times10^{-7}$ on $\Gamma_{1v}$, and the tensors come out isotropic to
$2.9\times10^{-8}$ with nothing imposing cubic symmetry.

\begin{refuses}
\textbf{Refuses.} A band inside a \textbf{degenerate multiplet} has no tensor of
its own --- the eigensolver's arbitrary rotation within the manifold rotates it
--- so those bands come back as \code{nan} and what is reported instead is the
multiplet's \emph{summed} inverse mass, which is the trace and is invariant.
Silicon's threefold $\Gamma_{25'}$ is the case, and Elk's own output is the
evidence: it prints $-5.8938$, $-3.7690$, $-3.7690$ for three states whose
energies agree to $5\times10^{-8}$~Ha. A \textbf{spin spiral} is refused, since
a k-stencil moves both spheres. \code{delta} is the stencil's outermost
displacement in both routes; the $O(h^2)$ truncation is removed by a Richardson
step and then \emph{reported} in \code{truncation} rather than tuned away.
\end{refuses}

\subsection{\texorpdfstring{$\langle L\rangle$, $\langle S\rangle$ and $\langle J\rangle$ on each atom}{L, S and J on each atom}}

Where the orbital moment of a spin-orbit magnet actually sits. From one site
density matrix per atom and shell, built out of the same projection a projected
density of states uses,
\begin{equation*}
  \rho^a_{(ms),(m's')} = \sum_{n\mathbf k} w_{n\mathbf k}\,
     c^a_{ms,n\mathbf k}\, \overline{c^a_{m's',n\mathbf k}},
  \qquad c = \langle \phi | S | \psi \rangle,
\end{equation*}
and then $\langle L_i\rangle = \sum_s \mathrm{Tr}_m (L_i \rho_{ss})$ and
$\langle S_i \rangle = \tfrac12 \sum_m \mathrm{Tr}_s (\sigma_i \rho_{mm})$, with
$J = L + S$. The $L$ matrices are the complex-basis ones conjugated into the
\emph{real} spherical harmonics the code works in, using the same \code{rot\_ylm}
unitary the spin-orbit projectors are built from.

\begin{pycode}
from defumat import Calculator

calc = Calculator.from_file("ni.scf.in")   # noncolin, lspinorb, nosym
momenta = calc.get_angular_momenta()

print(momenta.table())      # what Elk writes to LSJ.OUT
momenta.atoms[0].l          # <L> on atom 1, cartesian, units of hbar
momenta.atoms[0].j          # <J> = <L> + <S>
momenta.total_orbital       # summed over the cell
\end{pycode}

\textbf{Validated by identities, not by another code's number} --- Elk
integrates over a muffin tin of a stated radius and this projects onto an
orbital set, so the two differ by definition the way Löwdin and muffin-tin
charges do. $\langle L\rangle$ is \emph{quenched to zero} without spin-orbit
coupling, measured at $1.7\times10^{-16}$ on silicon; with the coupling on,
fully-relativistic nickel gives $\langle L_z\rangle = 0.0364767\,\hbar$ against a
measured orbital moment of about $0.05\,\mu_B$, with
$|L|/|S| = 0.11665$ against an experimental $m_L/m_S \approx 0.1$ and $\langle
L\rangle \parallel \langle S\rangle$ (Hund's third rule). Driving the moment
along $z$, $x$ and $y$ gives $|\langle L\rangle| = 0.0364767$ in all three,
with a spread of 7.3e-11 over the three cubic axes, with nothing imposing that a magnitude is a scalar.

\textbf{A symmetry-reduced k-set works}, and what makes it work is that
$\langle L\rangle$ and $\langle S\rangle$ are averaged over the group as
\textbf{axial} vectors --- each operation carrying $\det(R)$ and the
time-reversal sign beside the rotation, the site charge carrying only the atom
permutation. The distinction is not a refinement: the \emph{polar} average of an
axial vector is exactly zero on any centrosymmetric crystal, so writing the
obvious one would return a clean and plausible nothing rather than a wrong
number.

\begin{refuses}
\textbf{Refuses.} A \textbf{fully-relativistic ultrasoft or
PAW} dataset, because the spinor overlap's off-diagonal spin blocks are
\code{qq\_so} and the projection here applies the scalar $S$ to each component;
a fully-relativistic \emph{norm-conserving} dataset has $S = 1$ and is exact. And
a \textbf{spin spiral}, whose two components do not share a set of atomic
orbitals. \code{pw.x}'s \code{lorbm} computes the \emph{cell's} orbital
magnetization and has no per-atom counterpart to check this against --- that
quantity is the next subsection, and the two are worth reading together.
\end{refuses}

% =====================================================================
\subsection{The orbital magnetization of the cell}

The other half of a magnet's moment, and the half no integral over the unit cell
can give. A Bloch state's current circulates through the whole crystal, so
$\int \mathbf r \times \mathbf j$ depends on where the cell is cut --- the same
difficulty the electric polarization has. What is well defined is a k-space
quantity (Thonhauser \emph{et al.}, PRL \textbf{95}, 137205 (2005); Ceresoli
\emph{et al.}, PRB \textbf{74}, 024408 (2006)):
\begin{equation*}
  \mathbf M = \frac{e}{2\hbar c}\,\mathrm{Im} \sum_n \int [d\mathbf k]\,
    \langle \partial_{\mathbf k} u_n | \times
      \left( H_{\mathbf k} + E_{n\mathbf k} - 2\mu \right)
    | \partial_{\mathbf k} u_n \rangle,
\end{equation*}
a \emph{local} circulation inside each cell plus an \emph{itinerant} one along
the boundary. Neither is separately measurable and their sum is. The derivative
is the covariant finite difference over a uniform mesh --- the dual states
$|w_m\rangle = \sum_n (M^{-1})_{nm} |u_{\mathbf k + \mathbf b, n}\rangle$ of the
neighbouring manifold --- so nothing is divided by $E_n - E_m$ and the answer is
invariant under any unitary mixing of the occupied bands.

\begin{pycode}
from defumat import Calculator

calc = Calculator.from_file("i-atom-soc.in")    # noncolin, lspinorb, gapped
orbm = calc.get_orbital_magnetization(divisions=(3, 3, 3))

orbm.lc, orbm.ic     # pw.x's two Kubo terms, Bohr magnetons per cell
orbm.total           # their sum, at mu = 0
orbm.chern           # the Chern vector; zero for an ordinary insulator
\end{pycode}

\textbf{Against \code{pw.x}'s \code{lorbm}} on the same $3\times3\times3$ grid,
an iodine atom in a twelve-bohr box: $M_{LC} = -0.5986388$ against
$-0.5986370$ and $M_{IC} = -0.5757995$ against $-0.5757987$, so the total is
$-1.1744384$ against $-1.1744357$. What is left at $2\times10^{-6}$ is the
difference between two separately converged densities: re-running at
\code{conv\_thr} of 1e-6, 1e-8 and 1e-12 moves $M_{LC}$ by 5e-6, 2e-7 and 2e-6
against the same reference. \textbf{The physics check is independent of both}:
iodine is $5s^2 5p^5$, one hole in the p shell, and Hund's rules give
$L = 1$ locked parallel to $S = 1/2$. The site projection of the previous
subsection measures $\langle L_z \rangle = 0.99977\,\hbar$ on the same run, and
the modern theory gives $1.1744\,\mu_B$ --- 17\% more, which is what a
projection onto atomic orbitals does not see: the circulation outside those
orbitals, and the nonlocal pseudopotential's own contribution to the velocity,
since $\langle L \rangle$ is $\mathbf r \times \mathbf p$ where this contracts
$H$.

\textbf{Two conventions are QE's and are worth knowing.} What both codes print
is $M(\mu = 0)$: \code{orbm\_kubo} imports the Fermi level and never uses it,
and the term it drops is $\mu$ times the Chern vector, which vanishes for any
topologically trivial manifold. And the two printed terms are \emph{not} the
papers' LC/IC split --- the Fortran prints $\mathrm{Im}\langle \partial u|H|
\partial u\rangle$ and $\mathrm{Im}\langle \partial u|E|\partial u\rangle$, which
differ from $(H-E)$ and $2(E-\mu)$ term by term and agree in the sum. This
follows the Fortran, so the two halves can be compared and not only the total.

\begin{refuses}
\textbf{Refuses.} \textbf{Ultrasoft and PAW}, which \code{pw.x} refuses in the
same place (\code{setup.f90}: `Orbital Magnetization not implemented with
USPP/PAW'): the overlap between neighbouring manifolds would need $q_{ij}(b)$
and the dual states would have to be dual in the $S$ metric. A run
\textbf{without spin-orbit coupling}, where the orbital moment is quenched
identically --- measured here as $10^{-9}$ with \code{soc\_scale = 0} in the
same relativistic dataset. \textbf{\code{nspin = 2}}, which is the same
statement: each channel is separately time-reversal symmetric. A
\textbf{metal}, or any manifold not separated from the band above it, which the
gap check refuses by name. A \textbf{spin spiral}. And a mesh with
\textbf{two divisions} along a direction that carries a derivative, where a
point's two neighbours are the same k-point and the central difference is an
alias; one division is allowed and sets that direction's derivative to zero,
which is what a slab means.
\end{refuses}

% =====================================================================
\subsection{The Fermi-surface nesting function}

How much of the Fermi surface maps onto itself when translated by $\mathbf q$:
\begin{equation*}
  N(\mathbf q) = \frac{1}{N_k}\sum_{\mathbf k} g(\mathbf k)\, g(\mathbf k + \mathbf q),
  \qquad
  g(\mathbf k) = g_s \sum_{s,n} \delta(\varepsilon_{sn\mathbf k} - E_F).
\end{equation*}
$g(\mathbf k)$ is the density of states \emph{at one k-point} --- a picture of
the Fermi surface on the grid --- so $N(\mathbf q)$ is the $\omega\to0$
imaginary part of the Lindhard function stripped of its matrix elements: the
\emph{geometric} half of an instability. Where it is large, a perturbation of
wavevector $\mathbf q$ connects many occupied states to many empty ones at no
energy cost, which is what softens a phonon, opens a charge-density-wave gap, or
picks a spin spiral's pitch.

\textbf{It is a convolution.} Elk writes an $O(N_q N_k)$ double loop with
$\mathbf k + \mathbf q$ folded back onto the grid; that fold makes the sum a
cyclic cross-correlation, so $N = \mathrm{ifftn}(|\mathrm{fftn}(g)|^2)/N_k$
gives every $\mathbf q$ in one transform. Measured on a $12^3$ aluminium grid
that is \textbf{0.001~s} here against Elk's 0.37~s. A symmetry-reduced wedge is
\emph{unfolded} onto the complete grid rather than refused, since
$\varepsilon_n(R\mathbf k) = \varepsilon_n(\mathbf k)$: 72 diagonalisations
instead of 1728.

\begin{pycode}
from defumat import Calculator

calc = Calculator.from_file("al.scf.in")
nest = calc.get_nesting(grid=(12, 12, 12))   # the grid is the convergence knob

nest.nesting        # (nq,) N(q), in (states/Ry/cell)^2
nest.qpoints        # (nq, 3) crystal coordinates, on the unshifted grid
nest.ratio          # N(q) / D(E_F)^2, dimensionless and 1 on average
nest.peak()         # the largest N(q) away from q = 0, and where it is
nest.along(2)       # (q_3, N) along one crystal axis, for a line plot
nest.as_grid()      # (n1, n2, n3), for a slice
nest.fermi_dos      # D(E_F) in states/Ry/cell
nest.elk_units      # the same array in NEST3D.OUT's normalisation
\end{pycode}

The functional entry point is \code{defumat.workflows.run\_nesting}, which
takes the converged density directly.

\textbf{Two identities come with the normalisation and both are worth reading.}
The mean of $N$ over the whole q-grid is exactly $D(E_F)^2$
(\code{sum\_rule} reports the residue; it is $3\times10^{-16}$ relative). And
$N(0) \ge N(\mathbf q)$ for every $\mathbf q$ by Cauchy--Schwarz, so
\textbf{the nesting peak is always a peak away from the origin} ---
\code{peak()} excludes $\mathbf q = 0$ for that reason, not as a plotting
convenience.

Validated against a closed form and against a quantity computed by an unrelated
route. For free electrons $N(q) = \Omega/4\pi^2 q$ below $2k_F$ and zero above:
the code reproduces the $1/q$ law to $10^{-4}$ and the hard cut to $10^{-16}$. A
half-filled hydrogen chain has $2k_F = \pi/c$, so it must nest at
$\mathbf q = (0,0,\tfrac12)$ --- it does, at \textbf{99.8\%} of the
Cauchy--Schwarz bound, and \code{relax\_spiral\_q} relaxes the corresponding
spin spiral to $0.500014$ by going downhill in a magnetic total energy, which
shares nothing with this. Against the vendored Elk binary on aluminium,
$N(0)/D(E_F)^2$ agrees to 1.5\%; the raw $N(0)$ differs by 4.7\%, which is the
all-electron Fermi surface underneath and is quoted rather than tuned away.

\begin{refuses}
\textbf{Refuses.} A constrained \code{tot\_magnetization}, which has one Fermi
level per channel, so $g(\mathbf k)$ is two different surfaces; a
\textbf{fixed-occupation} run, which never searches for a level, so
$\delta(\varepsilon - E_F)$ has no argument; and a \textbf{spin spiral}, because
the quantity is a statement about the state a spiral grows out of --- compute it
on the paramagnet and compare its peak with \code{relax\_spiral\_q}'s answer.
The delta defaults to a \textbf{Gaussian} even when the run used
Methfessel--Paxton or cold smearing: those go negative on the wings, so
$g(\mathbf k)$ can be negative and $N(\mathbf q)$ would have no sign at all.
Pass \code{smearing=} to have the run's own anyway.
\end{refuses}

\subsection{Magnetocrystalline anisotropy}

The energy it costs to point a magnet's moment one way rather than another ---
what makes a hard magnet hard and what pins a spiral into a plane. It is tenths
of a meV per atom against total energies of hundreds of Rydberg, so taking it as
a difference of two self-consistent total energies asks two SCF runs to agree in
their ninth digit. The \emph{force theorem} does not ask that. Converge the
magnet \textbf{without} spin-orbit coupling, rotate the converged density so its
magnetization points along $\mathbf n$, and diagonalise \textbf{once} with the
coupling switched on. At frozen density every term of the total energy except
the band energy is a functional of $\rho$ alone, so
\begin{equation*}
  E(\mathbf n_1) - E(\mathbf n_2)
    = E_{\text{band}}(\mathbf n_1) - E_{\text{band}}(\mathbf n_2),
  \qquad
  E_{\text{band}} = \sum_{n\mathbf k} w_{n\mathbf k}\, \varepsilon_{n\mathbf k},
\end{equation*}
exactly, and the whole calculation is one diagonalisation per direction.

\textbf{The cheaper thing gives zero.} Freezing the wavefunctions as well and
taking $\langle\psi|H_{\text{SOC}}|\psi\rangle$ once returns no anisotropy at
all: the coupling enters at first order as $\xi\,\langle \mathbf L\rangle\cdot
\mathbf n$, and the orbital moment of a scalar-relativistic collinear state is
quenched. The anisotropy is second order in the coupling, and what supplies it
is the repulsion between levels that a diagonalisation performs and an
expectation value does not. \code{frozen\_expectation} computes that vanishing
term anyway.

\textbf{Two pseudopotential files, one density.} The first leg runs with a
scalar-relativistic dataset and the one-shot leg with the fully-relativistic
dataset of the \emph{same generation}; only the density crosses, and a density
does not know which file made it. That is what makes an ultrasoft anisotropy
reachable, the alternative --- averaging one relativistic file's $j$ channels
back --- being refused for ultrasoft and PAW by \code{pw.x} itself.

\begin{pycode}
from defumat import Calculator

scalar = Calculator.from_file("scf.in")          # nspin = 2, Co.pbe-nd-rrkjus
spinor = Calculator.from_file("nscf.in")         # noncolin, lspinorb, rel- file
mae = scalar.get_anisotropy(spinor, directions="xyz")

print(mae.anisotropy_mev, mae.easy_axis)
print(mae.difference(0, 2))                      # in-plane minus out-of-plane
\end{pycode}

\code{directions} takes a list of vectors, an integer asking for that many
points spread evenly over the sphere, or one of \code{"cardinal"},
\code{"xz"} and \code{"xyz"}. \code{projected=True} adds the decomposition of
the band energy over atomic orbitals, resolved by atom, $\ell$, $m$ and spin,
which is what says \emph{which} orbitals supply the anisotropy; comparing its
total against $E_{\text{band}}$ measures the spilling.

\code{soc\_scale} scales the spin-orbit part of the nonlocal potential and the
overlap while keeping the same dataset and the same k-points. At $0$ the
anisotropy is exactly zero, which is the control for the whole calculation:
without the coupling the Hamiltonian is invariant under a global spin rotation,
so a band energy \emph{cannot} depend on where the moment points.

\textbf{A PAW dataset needs one file rather than two, and \code{becsum} beside
the density.} A PAW Hamiltonian's one-centre coefficients $D^{1}_{ij}$ are a
functional of \code{becsum} exactly as the grid potential is a functional of
$\rho$, so what the theorem freezes on this dataset is the \emph{pair}: the
potential has two representations and both are held fixed. \code{becsum} cannot
cross between the two files above --- a scalar-relativistic dataset carries one
projector per $(n,\ell)$ channel and its fully-relativistic partner one per
$(n,\ell,j)$, so the two arrays are indexed by different things --- and that is
why \code{pw.x} refuses the path outright rather than for a reason that could be
plumbed around. The way round it is \code{soc\_scale}: run the
\emph{self-consistent} leg on the fully-relativistic file with
\code{soc\_scale = 0} and the one-shot leg on the same file with the coupling
on, so that both legs share a projector set by construction. \code{becsum} is
then laid along the requested direction with the density, off the axis the
\emph{density} defines rather than its own, which is what an antiferromagnet
needs: its two sublattices point opposite ways while the cell has one frame to
turn. It is a projection onto that axis rather than a rigid rotation, exactly as
it is for the density, so a canted first leg hands over the collinear part of
its moment and not the canting.

\begin{pycode}
from defumat.scf import run_scf
from defumat.workflows.anisotropy import run_anisotropy

# one fully-relativistic PAW file, run twice: the coupling off for the
# self-consistent leg and on for the one-shot.
scf = run_scf(system.with_soc_scale(0.0), pseudos, conv_thr=1.0e-8)

mae = run_anisotropy(system, pseudos, scf.density,
                     directions=[(0, 0, 1), (1, 0, 0), (1, 1, 0)],
                     becsum=scf.becsum)
\end{pycode}

\noindent \code{Calculator.get\_anisotropy} passes \code{becsum} for you when
the first leg's projector set is the one the second leg has, and hands over
nothing when it is not, so the two-file route above is unaffected. With the
coupling off the identity holds on a PAW dataset as it does on an ultrasoft
one: $3.1\times10^{-10}$~meV over five directions on antiferromagnetic oxygen.
Leaving the rotation out --- carrying \code{becsum} but laying it along $z$ for
every direction --- gives \textbf{468~meV} on that same cell, whose answer is
exactly zero, which is why the identity is worth running.

\textbf{The torque: the same number from one angle.} Taking the anisotropy as
$E(\mathbf n_1)-E(\mathbf n_2)$ is $10^{-5}$~Ry out of $10^{2}$ --- seven digits
of cancellation. The torque takes it as a \emph{derivative} evaluated once,
where nothing cancels: for $E(\theta)=K_1\sin^2\theta$,
\begin{equation*}
  -\frac{\mathrm dF}{\mathrm d\theta} = -K_1\sin 2\theta,
  \qquad\text{so at }\theta=45^\circ\text{ the torque is }-K_1 .
\end{equation*}

\begin{pycode}
torque = scalar.get_torque(spinor)              # 45 degrees by default
print(torque.anisotropy_constant_mev)           # K1 in meV
\end{pycode}

\textbf{It differentiates the free energy, and for a smeared metal that is not
the band energy.} A Hellmann--Feynman derivative at frozen occupations gives
$\mathrm dF/\mathrm d\theta$ with $F=\sum w\varepsilon-TS$; $\sum w\varepsilon$
carries an extra $\sum(\mathrm dw/\mathrm d\theta)\varepsilon$ that the entropy
cancels. The difference is not small. On tetragonal cobalt, as the smearing
narrows:

\begin{center}
\begin{tabular}{lccc}
\code{degauss} (Ry) & band difference & free difference & torque \\
\hline
0.020 & 1.2353 & 0.5523 & 0.5523 \\
0.005 & 0.5434 & 0.5291 & 0.5291 \\
0.002 & 0.5308 & 0.5308 & 0.5308 \\
\end{tabular}
\end{center}

\noindent all in meV. The torque tracks the free-energy difference at every
width and both converge on $K_1=0.531$; at a production \code{degauss} of 0.02
the \emph{band}-energy difference reads 2.3 times that. Use
\code{MagneticAnisotropy.free\_energies} to compare like with like.

\begin{refuses}
\textbf{Refuses.} \textbf{PAW without \code{becsum}}, and \textbf{PAW on the
two-file route at all}: the one-centre coefficients are built from a
\code{becsum} that a run on the other file cannot supply, its projectors being
indexed by $(n,\ell,j)$ where the other's are indexed by $(n,\ell)$. The
refusal names the one-file route in its message, and a \code{becsum} whose
shape does not match this leg's projector count is refused rather than
reshaped. A \textbf{Hubbard $U$}, whose \code{ns} is not in the handoff either; a
potential-only \textbf{meta-GGA}, whose $\tau$ is not; a converged
\textbf{magnetic field} or constrained moment, whose energy sits outside the
reported total; and a \textbf{spin spiral}, which refuses spin-orbit coupling
permanently and so has no coupling to switch on. A direction other than the
system's own \code{angle1}/\code{angle2} needs \code{nosym}: a magnetic
noncollinear run reduces its k-grid with the \emph{magnetic} symmetry group,
which depends on where the moment points, so two directions would otherwise be
sampled on two different wedges. \code{soc\_scale} takes $0$ or $1$
and nothing in between: the coupling-free end of the overlap is spin-independent,
so the anisotropy vanishes identically there whatever else it is, but a blend of
it with the true overlap partway is not the overlap of any set of projectors and
$S$ stops being a usable metric.
\end{refuses}

\subsection{The relaxed anisotropy: total energies instead of band sums}

The section above buys its precision by freezing the density: every term of the
total energy except the band sum is a functional of $\rho$ alone, so it cancels
between two directions exactly. The price is that the density is not allowed to
respond to the spin-orbit field. \code{run\_relaxed\_anisotropy} pays the other
way round --- one whole self-consistent noncollinear run per direction, and a
difference of two \textbf{total} energies:
%
\[
  E_{\mathrm{MAE}} = E_{\mathrm{tot}}[\hat{\mathbf n}_1] - E_{\mathrm{tot}}[\hat{\mathbf n}_2],
  \qquad \text{each } E_{\mathrm{tot}} \text{ converged at its own } \hat{\mathbf n}.
\]
%
The extra energy the relaxation finds is variational, so it lowers \emph{both}
directions, and what survives in the difference is second order on a quantity
that is already second order in the coupling. Nothing bounds it in general,
which is why it is measured rather than argued about.

\begin{lstlisting}[language=Python]
from defumat import Calculator

soc = Calculator.from_file("ni-tetragonal-relaxed-mae-paw.in",
                           pseudo_dir="tests/data/pseudo")
mae = soc.get_relaxed_anisotropy(directions="xz")

print(mae.energies_mev)      # relative to the first direction
print(mae.easy_axis, mae.anisotropy_mev)
print(mae.converged, mae.drifts)   # read these BEFORE the number
\end{lstlisting}

There is \textbf{one} calculator here and no second leg. That is the whole
practical difference from the force theorem, and it is what lets this route
reach \textbf{PAW} and a Hubbard $U$: the theorem hands a density from a
scalar-relativistic run to a fully-relativistic one, and \code{ddd\_paw} and
\code{ns} are properties of \emph{wavefunctions} that cannot cross that gap. A
relaxed run hands nothing over --- each direction converges its own
\code{becsum} from its own states with one dataset throughout.

\textbf{Read \code{.converged} and \code{.drifts} first, and the second one is
the trap.} Nothing holds the moment in a relaxed run, which is the point of it;
so a direction that is not a stationary point of the anisotropy energy will turn
towards one that is, and its total energy then belongs to a state other than the
one asked for. A cardinal axis of a cubic or uniaxial crystal is stationary by
symmetry and does not move; an oblique direction on a strongly anisotropic
magnet need not be. A drift past one degree raises a warning naming the
direction and where the moment went.

\begin{refuses}
A \textbf{magnetic field} or a constrained moment, because the field's energy is
deliberately outside the reported total, so two directions' totals would differ
by a Zeeman term neither accounts for --- and because a constraint defeats the
drift check, the constraint being what holds the moment rather than the
anisotropy. A potential-only \textbf{meta-GGA}, which has no energy functional
at all, so its printed total is not the value of anything the run minimised. A
\textbf{spin spiral}, which refuses spin-orbit coupling permanently. A
\textbf{collinear} run, which has no direction to point. And \code{nosym} is
required rather than merely advised: a magnetic noncollinear run reduces its
k-grid with the \emph{magnetic} symmetry group, which depends on where the
moment points, and here that is worse than it is for the force theorem --- a
total energy carries the Ewald and Hartree terms too, so the contamination is
not confined to the bands.
\end{refuses}

\subsection{X-ray and magnetic structure factors}

What a diffractometer measures is not a density but its Fourier coefficients,
one per reflection:
\begin{equation*}
  F(\mathbf H) = \int_\Omega \rho(\mathbf r)\, e^{i\mathbf H\cdot\mathbf r}\,d^3r,
  \qquad
  F_j(\mathbf H) = \int_\Omega m_j(\mathbf r)\, e^{i\mathbf H\cdot\mathbf r}\,d^3r
\end{equation*}
for every reciprocal lattice vector with $|\mathbf H| \le \code{hmax}$ --- X-rays
scattering off the charge, neutrons off the magnetization. Elk's tasks 195 and
196; \code{pw.x} has neither. In a plane-wave code the integral \emph{is} an
array that already exists, up to the crystallographers' two conventions (the
positive phase, and no $1/\Omega$), so the whole step is one transform and a
gather: \textbf{6~ms} against the seconds its SCF cost. $F(0)$ is the electron
count of the cell and $F_j(0)$ its magnetic moment, which is what anchors both.

\begin{pycode}
from defumat import Calculator

calc = Calculator.from_file("si.scf.in")
factors = calc.get_structure_factors(hmax=5.0)

factors.charge               # (nh,) complex F(H), in electrons
factors.magnetization        # (nh, ncomp) in Bohr magnetons, or None
factors.vectors.indices      # (nh, 3) the (h k l) after the output transform
factors.vectors.length       # (nh,) |H| in 1/bohr
factors.vectors.multiplicity # the size of each star that was collapsed
factors.of((2, -2, -2))      # one reflection, by the index the table prints
print(factors.table(limit=10))

# Elk's wsfac: rebuild the density from the states in an energy window (Ry),
# which is what makes this a probe of bonding rather than of the whole charge.
bonding = calc.get_structure_factors(hmax=3.0, window=(-1.0, 0.2))
\end{pycode}

The functional entry point is \code{defumat.workflows.sfac.run\_structure\_factors}.
\code{transform="conventional"} labels the reflections on the conventional cubic
axes (Elk's \code{vhmat}), \code{reduce=False} keeps every member of every star,
and \code{core=True} adds a dataset's frozen core-correction charge.

\textbf{A pseudopotential density is valence-only, and the comparison with an
all-electron code is structured by reflection class rather than by one
tolerance.} On silicon, against the vendored Elk binary on the same cell and
k-grid: the origin is $8$ electrons here and $28$ there, and every
\emph{allowed} reflection is core-dominated in the same proportion --- $(111)$ is
$1.7495$ against $15.14$, and no tolerance makes that agree. What \emph{is}
comparable is a \textbf{forbidden} reflection, where the spherical part of every
atom cancels and the core with it: silicon's $(222)$ is \textbf{0.347406} here
against \textbf{0.33416}, 4\%, on a quantity whose entire value is the
aspherical bonding charge. That last statement is checked without any other
code, since this package's own SCF starting guess is a superposition of free
atoms: it gives $0$ at $(222)$ to $10^{-10}$. The space-group forbidden
reflections --- diamond's $(200)$, $(420)$ --- vanish to $10^{-9}$ here and
$10^{-13}$ there.

\textbf{Norm-conserving, ultrasoft and PAW all work, and the difference between
them is measurable.} $F(0) = 8.0000000000$ on each, which is the check that the
augmentation charge reached the transform at all --- for a soft dataset most of
the density near the nucleus \emph{is} that charge. The like-for-like reference
is not Elk's all-electron total but its \textbf{valence-only} run, which
\code{wsfac} above the 2p core produces: against it the augmentation charge is
worth a factor of two, $(111)$ coming out $-0.7\%$ for ultrasoft and PAW against
$-1.5\%$ norm-conserving, and $(004)$ $+31\%$ against $+61\%$. It is worth
\emph{nothing} on the forbidden $(222)$, which is 0.3474 on all three to
$3\times10^{-4}$: a silicon atom in diamond sits at $\bar 4 3m$, whose lowest
non-spherical invariant is $l = 3$, and an $s,p$ dataset's augmentation charge
carries multipoles only to $L = 2$. That is also what the 4\% against
all-electron is, by elimination --- not the k-grid, not either basis, not the
functional (0.07\%), and \textbf{not the core}, since Elk's own valence-only run
gives the identical 0.334165.

\textbf{Magnetic reflections are a low-$|\mathbf H|$ quantity and the boundary is
measured.} Ferromagnetic bcc iron against Elk's task 196, comparing the
normalized $F_{\mathrm{mag}}(\mathbf H)/F_{\mathrm{mag}}(0)$ because the moments
differ (2.2145 $\mu_B$ here, 2.0613 there): 5.8\% at the first reflection, 17\%
at the second, worse outwards. Raising \code{ecutwfc} from 30 to 45 moves those
by less than 0.1\%, so it is the pseudisation and not the basis --- iron's moment
is in the 3d shell, inside the radius the dataset smooths. The cheapest cell
carries the sharpest magnetic statement instead: an antiferromagnetic hydrogen
chain has charge factors $2.000000, 0, 0.8587, 0$ along the chain and magnetic
ones $0, 1.2254, 0, 0.3862$, so it scatters neutrons exactly where it scatters
no X-rays, and the zeros are exact.

\begin{refuses}
\textbf{Refuses.} An \code{hmax} past $\sqrt{\code{ecutrho}}$: the FFT box
carries a coefficient at every frequency out to its corner, and past the density's
own cutoff those are the aliased tail of $|\psi|^2$ rather than Fourier
components --- small, smooth and entirely plausible, which is why this is a guard
and not a check on the answer. An empty or inverted energy window, and
\code{core=True} on a dataset with no nonlinear core correction.
\textbf{Said rather than refused}, because the numbers are right and only a
claim about them would be wrong: these are \emph{valence} structure factors, and
a star's multiplicity applies to the charge, the magnetic members of a star
being related by $\det(R)\,R$ rather than being equal.
\end{refuses}

\subsection{Scanning-tunnelling microscopy images}

Tersoff and Hamann's result is that the current an s-wave tip draws at
$\mathbf r$ is the sample's local density of states there, at the energy the
bias selects. So an STM image is not a new calculation, it is the
\emph{density built from different occupations}:
\begin{equation*}
  \rho_{\mathrm{STM}}(\mathbf r)
   = \sum_{n\mathbf k} w_{\mathbf k}\,
     \tilde\delta(E - \varepsilon_{n\mathbf k})\,
     |\psi_{n\mathbf k}(\mathbf r)|^2 ,
\end{equation*}
with $\tilde\delta$ the same smeared delta a density of states is built from.
Elk's task 162 and \code{pw.x}'s \code{plot\_num = 5} both say exactly that.
Two energy selections: with no \code{bias} the weight is a delta at $E$ (Elk's,
a density of states per unit volume, in $1/(\mathrm{bohr}^3\,\mathrm{Ry})$),
and with one it is the window $[E, E+V]$ (QE's, a density in
electrons/bohr$^3$) --- a \emph{negative} $V$ imaging the filled states, which
is a real experiment's sign convention and the mode a semiconductor surface
needs.

\begin{pycode}
from defumat import Calculator

calc = Calculator.from_file("graphene.scf.in")

# constant height: the tip plane at crystal coordinate 0.85 along c
image = calc.get_stm(height=0.85, shape=(64, 64))

image.values        # (64, 64) the tunnelling density on the plane
image.coordinates   # (64, 64, 2) in-plane cartesian coordinates, bohr
image.extent()      # (x0, x1, y0, y1) for imshow
image.integral      # int rho_STM d3r = D(E_F): the sum rule
image.density       # the same field on the FFT grid

# filled states, 0.15 Ry below the Fermi level
filled = calc.get_stm(height=0.85, bias=-0.15)

# constant current: the corrugation, in bohr along the surface normal
scan = calc.get_stm(height=0.70, mode="constant-current",
                    current=1.0e-5, heights=(0.0, 6.0), nheights=80)
scan.heights_bohr   # (n1, n2), nan where the set-point is not crossed

# a magnetic tip: the image depends on which way it points
up = calc.get_stm(height=0.85, spin="up")            # collinear
along_x = calc.get_stm(height=0.85, spin=(1, 0, 0))  # noncollinear

# Elk's three corners, in crystal coordinates, for any other plane
side = calc.get_stm(plane=((0.15, 0, 0), (0.15, 0.3, 0), (0.15, 0, 1)))
\end{pycode}

The functional entry point is \code{defumat.workflows.stm.run\_stm}. A delta at
the Fermi level wants a denser k-grid than the SCF's, which \code{grid=}
re-solves the bands on and \textbf{recomputes the Fermi level for}; \code{smearing}
picks the delta (a Gaussian by default whatever the run used, because
Methfessel--Paxton and cold smearing go negative on their wings and a negative
tunnelling current is meaningless); and \code{band\_cutoff} is \code{stm.f90}'s
fixed truncation at three widths, off by default because it is an approximation.

\textbf{The plane is read out of the grid exactly.} A density is a finite sum
of plane waves, so its value at a point between grid points is that sum
evaluated there rather than a spline through neighbours. It matters for the one
thing an image is about: the corrugation in the vacuum is a small modulation on
a quantity falling by orders of magnitude across one grid spacing, and an
interpolant of it carries the grid's own periodicity as a false corrugation.
The three corners are in \emph{crystal} coordinates, Elk's \code{vclp2d}, and
the sampling excludes the far edge ($i/n$, not $i/(n-1)$) so a whole-cell image
tiles.

\textbf{Against \code{pp.x}: $6.7\times10^{-10}$} on every point of the
$15^3$ grid of QE's own fcc aluminium, reading its \code{filplot} text dump so
that nothing is interpolated on either side. Three of QE's conventions are
needed to make it the same sum: \code{stm.f90} divides by the volume and not by
\code{degauss}, so its image is the zero-bias one times the width; it uses the
run's own smearing; and it truncates the band sum at three widths. That
truncation is worth 0.4\% \emph{and the complete sum is the smaller one} ---
Marzari--Vanderbilt's delta is negative for $x > \sqrt2$, so the states QE drops
carry negative weight. The check that involves no other code is the sum rule
$\int\rho_{\mathrm{STM}}\,d^3r = D(E)$ against \code{compute\_dos}, which holds
to $10^{-10}$ on a scalar run and on a spinor one, and in its window form on an
ultrasoft one --- the augmentation charge follows the tunnelling weights, which
is why an ultrasoft or PAW image works here and \code{stm.f90}'s does not.

\textbf{A magnetic tip is the part neither code has.} It counts the states
whose spin is along its own moment, so with $P$ the tip's polarization
\begin{equation*}
  \rho_{\mathrm{STM}}(\mathbf r; \hat{\mathbf n})
   = \tfrac12\left[\rho(\mathbf r)
     + P\,\hat{\mathbf n}\cdot\mathbf m(\mathbf r)\right],
\end{equation*}
and $P = 1$ makes the image a genuine spin channel. On an antiferromagnetic
hydrogen chain the charge image is flat to six figures while the two spin
projections are each other's mirror ($\pm 2.58\%$, antisymmetric to $10^{-3}$),
and they add back to the charge to $10^{-14}$. On a chain whose four moments
each turn 90 degrees from the last, a tip along $x$ sees atoms 1 and 3 with
opposite sign and atoms 2 and 4 \textbf{not at all} ($10^{-16}$ against a
contrast of 0.063), while a tip along $z$ sees none of the four --- with nothing
imposing any of it.

\textbf{The physics case is graphite.} In AB-stacked bilayer graphene one
sublattice of the surface layer sits above an atom of the layer below and the
other above a hollow, so the two are inequivalent and only one is bright at the
Fermi level: 1.64$\times$ here, on two atoms of the same element carrying the
same charge. That is the textbook result that half the atoms of graphite are
invisible to an STM, and it is a statement about the density of states \emph{at
the Fermi level} rather than about the charge. The tunnelling density falls off
log-linearly into the vacuum ($R^2 = 0.9992$ over five orders) at a rate of at
least $2\sqrt{V_{\mathrm{vac}} - E_F}$, which is Tersoff--Hamann's own statement
and involves none of the assembly; the bound is one-sided because graphene's
states at $E_F$ sit at $K$, whose in-plane momentum steepens the decay
(measured 1.86, against 0.98 for $k_\parallel = 0$).

\begin{refuses}
\textbf{Refuses.} A spin spiral (its two spinor components live on different
plane-wave spheres, so $|\psi(\mathbf r)|^2$ is not the lattice-periodic object
this sum builds), a constrained \code{tot\_magnetization} (one Fermi level per
channel, and a tip sees one energy), an applied magnetic field (its energy is
outside the reported total, so the level the tip would be put at is not the
field-free one the image would be read as), a spin direction on a run with no
magnetization, and a \emph{transverse} direction on a collinear run --- there
$m_x$ is absent rather than zero, and projecting onto the axis would be a claim
the calculation cannot make. A constant-current scan longer than the lattice
period along the plane's normal, where the tip meets the periodic image of the
surface and the density rises again.
\textbf{Said rather than refused}: a constant-current image returns \code{nan}
at every point where the set-point is not crossed inside the scan, with a
warning saying how many; and a zero-bias image of a gapped run is identically
zero, the delta sitting in the gap, which is what \code{bias} is for.
\end{refuses}

\subsection{Tunnelling spectra: $\mathrm{d}I/\mathrm{d}V$ over an energy axis}

An image is the local density of states at one tip energy, which is one picture
at one bias. The other section of the same function is the curve at one place
over many biases, which is what a scanning tunnelling microscope measures when
it is used as a spectrometer, and what resolves a gap, a band edge or a state
sitting inside a gap --- none of which an image shows, since an image at a single
energy reports only how much weight there is and not where in energy it sits,
\begin{equation*}
  \frac{\mathrm{d}I}{\mathrm{d}V}(\mathbf r, V) \;\propto\;
  \rho(\mathbf r, E_F + eV) \;=\;
  \sum_{\mathbf k, n} w_{\mathbf k}\,
     \delta\!\left(E_F + eV - \varepsilon_{\mathbf k n}\right)
     \left|\psi_{\mathbf k n}(\mathbf r)\right|^2 ,
\end{equation*}
the same sum \code{run\_stm} builds, evaluated at every energy of an axis.

\textbf{This is the measurement a modulation is read with}, and it is why the
quantity is here rather than only in the unit cell: a charge density wave is a
\emph{gap} that opens in antiphase with the charge maxima, a spin density wave
moves the two spin channels in opposite directions from cell to cell, and a
domain wall carries a state that lives nowhere else. Each of the three is an
energy axis at every position, and an ultracell, the subject of \emph{A
modulation over many unit cells}, is the calculation that has them.

\textbf{The energy axis is nearly free, and that is the whole of the
implementation.} $\psi_{\mathbf k n}(\mathbf r)$ does not depend on the energy,
so the states are sampled at the tip points \emph{once} and each energy after
that is one matrix product against the per-state weight. On an eight-cell spin
density wave with a $96\times12$ map, one, twenty-one and forty-one energies cost
2.50, 2.56 and 2.61 s, against 2.47 s for a single image: forty-one energies are
4.4 per cent more than one, and 39 times less than the forty-one images that are
the only other way to get them. Neither \code{pw.x} nor Elk has an axis ---
\code{stm.f90} takes one scalar \code{sample\_bias} and \code{wfplot.f90} sets
one delta --- so both compute a spectrum by being run again at every bias, paying
the whole sum each time.

\begin{pycode}
import numpy as np
from defumat import Calculator

calc = Calculator.from_file("si-ultracell-mag.in")
field = lambda x: 0.02 * np.cos(2 * np.pi * x[..., 0] / 8)
wave = calc.get_ultracell(supercell=(8, 1, 1), kgrid=(1, 2, 2), nbnd=32,
                          magnetic_field=field)

# a spectrum above the middle of each unit cell: one tip point per cell,
# which is the line cut across the modulation
tips = np.stack([np.arange(8) + 0.5, np.full(8, 0.5), np.full(8, 0.35)], axis=-1)
bias = wave.fermi_energy + np.linspace(-0.06, 0.06, 41)
sts = calc.get_ultracell_sts(energies=bias, tip=tips, width=0.02, spin="up")

sts.values        # (41, 8) dI/dV, one column per cell
sts.bias_axis     # (41,) the bias in Ry, measured from the Fermi level
sts.current       # (41, 8) I(V), the same axis integrated outwards from V = 0
sts.integral      # (41,) D(E) per unit cell -- the *charge's* sum rule, which
                  # is what the weights give exactly whatever the tip is

# a whole map at every energy instead, and the ordinary unit-cell spectrum
maps = calc.get_ultracell_sts(energies=bias, height=0.35, shape=(48, 8))
plain = calc.get_sts(energies=bias, tip=tips[:1], width=0.02, grid=(4, 4, 4))
\end{pycode}

The entry points are \code{run\_sts} and \code{run\_ultracell\_sts}, or
\code{Calculator.get\_sts} and \code{Calculator.get\_ultracell\_sts}.

\textbf{Reading it.} \code{energies} is the axis in Rydberg and has no default:
the axis \emph{is} the measurement and a range has no natural width. What comes
back is \code{values}, of shape \code{(nE,) + } whatever the sampling was --- a
row of tip points, or a plane. \code{current} is $I(V)$ within this
approximation, a tip whose own density of states is flat and a barrier that does
not depend on the bias; the decay of each state into the vacuum \emph{is} in it,
the prefactor is not, and the experiment's sign convention is kept, so the filled
states below the Fermi level come back negative. A magnetic tip works exactly as
it does for an image, through \code{spin} and \code{polarization}, and on a spin
density wave it is what is wanted: an unpolarized tip cannot see the wave at
first order and sees its square instead.

\begin{refuses}
\textbf{Refuses.} Everything an image refuses, and two more of its own, both
because $\psi(\mathbf r)$ is sampled and squared here where an image sums into a
density. A \textbf{symmetry-reduced k-set}: \code{run\_stm} goes through the
density, which is symmetrised, so a wedge gives it the whole zone's answer, while
nothing symmetrises this one and a wedge would return the sum over that wedge
alone --- a smooth, plausible, wrong density of states everywhere off a symmetry
axis. On silicon's $4\times4\times4$ grid reduced to eight k-points the wedge sum
differs from the whole grid's by 98 per cent of the peak, while its \emph{integral}
is right to 1 per cent and the image on that same reduced set is right to 0.3 per
cent --- so a sum rule would not catch it and the image is genuinely unaffected.
Pass \code{grid} (which builds the \emph{complete} grid here, where an image's
reduces) or run \code{nosym}. A tip inside an \textbf{augmentation
sphere} of an ultrasoft or PAW dataset, where the pseudo-wavefunction is not the
true one and the norm is short by 3 to 9 per cent. On a dense bulk crystal that
leaves nowhere to put a tip at all --- SiC's radius is 2.4 bohr and its layers
are closer together than that --- which is the right answer rather than an
obstacle: a tip sits above a \emph{surface}, in vacuum, where a pseudo-wavefunction
is the true one. An ultracell is exempt from the
\emph{first} of these, requiring \code{nosym} already, and takes the second:
its image, its spectrum and its transmission all ask the guard on the unit
cell's own coordinates, every copy of the ultracell holding the same atoms at
the same heights.
Also refused: \code{bias} together with an energy axis, since the window is the
integral the axis already carries and \code{current} is it; and a
constant-current spectrum, where the tip moves as the bias is swept, which is a
different experiment and $n_{\rm heights}$ times the cost.
\textbf{Said rather than refused}: \code{current} raises on an axis whose step is
coarser than the smeared delta, because a trapezoid that steps over a level
returns a number of the wrong order rather than a slightly wrong one --- one
level of width $w$ on an axis of step $h$ integrates to $1.00004$ at $h = w$,
$1.14$ at $2w$ and exactly zero by $250w$, which is what any sane axis does
against the $10^{-5}$~Ry width a fixed-occupation run gets by default. A run with
no Fermi level takes the middle of the gap as the zero of bias, with a warning,
rather than 0~Ry.
\end{refuses}

\subsection{Vertical tunnelling transport through a two-dimensional material}

An STM image says what the tip \emph{sees}. This says what gets
\emph{through}: an electron enters at a point $\mathbf r$ above the material
and leaves into an infinite plane below it --- the substrate the sheet sits on
--- so what decides the current is the \emph{nonlocal} Green's function between
the two rather than the local density of states at the tip. A \textbf{point}
tip makes $\Gamma_t$ rank one, which collapses the Landauer--B\"uttiker trace
exactly:
\begin{equation*}
  T(\mathbf r; E) \;=\; \mathrm{Tr}\!\left[\Gamma_t G^r \Gamma_s G^a\right]
  \;\propto\; \int_{\mathrm{plane}}\! \mathrm{d}^2 r'\,
  \bigl|G(\mathbf r, \mathbf r'; E)\bigr|^2 .
\end{equation*}
The substrate is taken to be \emph{infinite and featureless}, so it is
invariant under every lateral lattice translation and conserves the lateral
momentum. The exit integral is then diagonal in $\mathbf k$, and everything
reduces to one quadratic form per k-point,
\begin{equation*}
  T(\mathbf r; E) = \sum_{\mathbf k} w_{\mathbf k}
    \sum_{nm} a_{n\mathbf k}(\mathbf r)\, a^{*}_{m\mathbf k}(\mathbf r)\,
    S_{\mathbf k}[n,m],
  \qquad
  S_{\mathbf k}[n,m] = \int_{\mathrm{plane}} \psi^{*}_{n\mathbf k}
     \psi_{m\mathbf k}\, \mathrm{d}^2 r' ,
\end{equation*}
with $a_{n\mathbf k} = \psi_{n\mathbf k}(\mathbf r)
\sqrt{\tilde\delta(E - \varepsilon_{n\mathbf k})/\eta}$. So \textbf{the sum
over $\mathbf k$ is incoherent and the interference is between bands at the
same $\mathbf k$}, which is what "the substrate is featureless" means
physically. $S_{\mathbf k}$ is a Gram matrix, so the transmission is
non-negative by construction and is blind to a rotation inside a degenerate
multiplet.

\textbf{Either electrode can be magnetic.} A spin-polarized substrate is a
$2\times2$ acceptance $P_s = (1 + P_s\,\hat{\mathbf n}_s\!\cdot\!\boldsymbol\sigma)/2$
sitting \emph{inside} $S_{\mathbf k}$, between two bands. A spin-polarized
\emph{tip} is the same matrix on the other side, $\Gamma_t = \gamma_t P_t$, and
it is not a weight but a contraction: with the spinor Green's function
$G_{ss'}(\mathbf r, \mathbf r') = \sum_n a_{n,s}(\mathbf r)\,
\psi^{*}_{n,s'}(\mathbf r')$ the Landauer trace is no longer taken over the
tip's spin, and the two spinor components stop adding and start interfering,
\begin{equation*}
  T(\mathbf r) = \int_{\rm plane}\!\mathrm{d}^2 r'\,
     \mathrm{Tr}_{\rm spin}\!\left[\Gamma_t\, G\, \Gamma_s\, G^\dagger\right]
  \;=\; \sum_{\mathbf k} w_{\mathbf k}\,
     \mathrm{Tr}_{\rm spin}\!\left[P_t\, M_{\mathbf k}(\mathbf r)\right],
  \qquad
  M_{\mathbf k}[s,s'] = \mathbf a_s^{\mathsf T} S_{\mathbf k}\, \mathbf a_{s'}^{*} .
\end{equation*}
$M$ is itself a Gram matrix, $A^{\mathsf T} S A^{*}$ for the $(n_{\rm bnd},2)$
amplitude block, so it is positive semi-definite and $\mathrm{Tr}[P_t M] \ge 0$:
the non-negativity survives. With \emph{both} leads polarized the map depends on
the angle between the two moments, which is a \textbf{tunnelling
magnetoresistance} image. \textbf{Note the asymmetry of the names}: here
\code{spin} is the \emph{substrate} and \code{tip\_spin} the \emph{tip}, which is
the opposite way round from \code{get\_stm}, where \code{spin} is the tip. The
substrate was the polarizer first and the name was kept.

\begin{pycode}
from defumat import Calculator

calc = Calculator.from_file("graphene.scf.in")

# tip plane above the sheet, substrate plane below it, atoms in between
t = calc.get_vertical_transport(exit_height=0.38, height=0.62,
                                shape=(48, 48), broadening=0.02,
                                grid=(6, 6, 1))

t.image            # (48, 48) the transmission map, arbitrary units
t.interference     # coherent minus incoherent: what no local picture gives
t.notes["channels"]  # open transmission channels the substrate offers
t.extent()         # (x0, x1, y0, y1) for imshow

# a transport dI/dV at one place: the tip sampling is paid once
import numpy as np
spectrum = calc.get_vertical_transport(
    exit_height=0.38, tip=np.array([[0.0, 0.0, 0.62]]),
    energies=calc.get_scf().fermi_energy + np.linspace(-0.2, 0.2, 41),
    broadening=0.02, grid=(6, 6, 1))

# a spin-polarized substrate: only on a magnetic run, and the two
# directions partition the unpolarized answer exactly
magnet = Calculator.from_file("h-sheet-lsda.scf.in")
up = magnet.get_vertical_transport(exit_height=0.15, height=0.85,
                                   spin="up", broadening=0.05)
down = magnet.get_vertical_transport(exit_height=0.15, height=0.85,
                                     spin="down", broadening=0.05)

# a magnetic tip on a noncollinear run, and both electrodes at once:
# the map then measures the angle between the two moments
spinor = Calculator.from_file("h-sheet-noncolin.scf.in")
valve = dict(exit_height=0.15, height=0.85, broadening=0.05,
             shape=(24, 24), spin=(1.0, 0.0, 0.0), polarization=0.9)
parallel = spinor.get_vertical_transport(tip_spin=(1.0, 0.0, 0.0),
                                         tip_polarization=0.9, **valve)
antiparallel = spinor.get_vertical_transport(tip_spin=(-1.0, 0.0, 0.0),
                                             tip_polarization=0.9, **valve)
tmr = antiparallel.image.mean() / parallel.image.mean()   # 5.29 on this cell

parallel.tip_spin           # (1.0, 0.0, 0.0), beside .spin for the substrate
parallel.tip_polarization   # 0.9
\end{pycode}

\textbf{The amplitude is the on-shell one and that is a measurement, not a
preference.} The literal Landauer denominator $1/(E - \varepsilon + i\eta)$
cannot be evaluated by a sum over states: the states far from $E$ build the
barrier's evanescent decay entirely by cancellation between themselves. On a
cell small enough to diagonalise completely (367 plane waves, 367 exact bands)
the running sum for one $G(\mathbf r,\mathbf r')$ wanders over more than an
order of magnitude and lands only at the complete basis, with a cancellation
ratio of \textbf{349}. Replacing the denominator by its modulus --- which is
Bardeen's golden rule, the sample visited on shell --- keeps the interference
between bands degenerate at the tip energy, which is the interference a real
experiment lets happen, and converges: to \textbf{3e-7} in the mean between 8
and 45 bands. \code{method="resolvent"} is reachable and warns.

\textbf{Validated by identities}, since neither \code{pw.x} nor Elk computes
this. (QE's \code{pwcond.x} is a Landauer transmission of a different geometry:
two semi-infinite crystalline leads with the current along one axis, one
conductance per energy, and no point contact and so no map. Elk has none.) Widening the substrate from a plane to the whole cell makes
$S_{\mathbf k}$ the identity by orthonormality, so the quantity becomes the
Tersoff--Hamann tunnelling density of states \emph{exactly} --- agreeing with
\code{run\_stm} to \textbf{4.6e-13} with no factor between them, which is the
check a normalisation error could not hide in. $S_{\mathbf k}$ swept along its
own normal integrates to the identity to 1.4e-15, and agrees with a real-space
quadrature to 5e-16. A substrate polarized along $\hat{\mathbf n}$ plus one
along $-\hat{\mathbf n}$ is a substrate that takes both, to round-off; one
across the moment takes exactly half. The same holds for the \emph{tip}, to
\textbf{3.1e-16}, with the substrate polarized along a third direction at the
same time, and a tip at $P=0$ is exactly half the unpolarized map (1.5e-16), the
convention a nonmagnetic tip has in the STM images. A \emph{magnetic} tip in the
whole-cell limit reproduces the spin-polarized tunnelling density of states to
\textbf{5.8e-13}, on a sheet whose moment is put in a generic direction: the
transposed index order of $M$ would return the answer for a tip at
$(n_x, -n_y, n_z)$, which is real, non-negative and --- on a moment lying along
an axis --- identical, so the check is only a check with $m_y \ne 0$. On that
cell the mirrored tip's image differs by a factor of two. Both electrodes
polarized at $P = 0.9$ give a parallel-to-antiparallel contrast of a factor
\textbf{5.3}, with a perpendicular tip at the mean of the two to 5e-6. The same cell run as $n_{\rm spin}=1$,
$n_{\rm spin}=2$ and as a spinor agrees to 1e-10 and 1e-5. And the physics:
graphene's two Dirac states at $K$ are degenerate partners, so the symmetry of
that point makes the substrate's overlap on the pair a multiple of the identity
and leaves nothing off-diagonal for the current to interfere through; its map
correlates with the STM image at \textbf{1.000000} with zero interference, while
an AB bilayer --- whose current must cross both layer-polarized sheets --- falls
to \textbf{0.16}, its coherent map a factor of 26 below the incoherent one.

\begin{refuses}
\textbf{Refuses.} A k-set with more than one division along the stacking axis:
the lateral momentum is conserved exactly and the perpendicular one is not, so
two $k_\perp$ at the same $k_\parallel$ interfere with a phase that depends on
where the exit plane sits. A \emph{symmetry-reduced} k-set, because a wedge sums
to the map symmetrised over the whole point group while only the subgroup that
leaves the exit plane in place belongs to this geometry --- a mirror through the
slab exchanges the tip side with the substrate side; \code{grid=} builds the
whole grid for that reason. A tilted exit plane, which has no closed-form
in-plane orthogonality to stand on. A tip or exit plane \emph{inside} an
augmentation sphere on an ultrasoft or PAW dataset --- in the vacuum a
pseudo-wavefunction is the true one, which is why those datasets otherwise need
nothing extra. A spin-selective substrate \emph{or a spin-polarized
tip} on a run with no magnetization, there being nothing for either moment to
couple to and every direction taking half of everything, which is the charge map
again; and a transverse direction on a collinear run for either lead, where
$m_x$ and $m_y$ are absent rather than zero. And, inherited from the STM images
unchanged: a spin spiral, a constrained \code{tot\_magnetization}, an applied
magnetic field.
A \code{bias=} window whose axis does not resolve \code{broadening}: the current
is a trapezoid over a smeared delta, so an axis coarser than one point per width
steps over the levels rather than integrating them and the answer is wrong by
orders. Measured on the hydrogen sheet at $\eta = 0.02$~Ry over a 0.24~Ry window,
against an axis eight times finer: 0.997 of it at $h = \eta$, \textbf{1.010 at
$2\eta$}, 0.298 at $4\eta$ and 0.121 at $12\eta$ --- not monotone, so a coarse
axis reads high before it collapses. The refusal names the \code{nenergies} that
would do.
\textbf{Said rather than refused}: a geometry whose atoms do not lie between the
two planes warns --- a cell is periodic, so with both planes on the same side
the electron tunnels through the vacuum and around the periodic image, which is
a real number and not this one.
\textbf{Not claimed}: a finite contact patch (which breaks the k-diagonality),
a tip with structure beyond an s-wave, and an absolute conductance --- the tip
and substrate couplings are unfixed prefactors, so what the result carries is
the map and its contrast.
\end{refuses}

\subsection{Which k-points the tunnelling current comes out of}

The section above asks \emph{where on the surface} the current goes. This asks
the conjugate question: \emph{which states} carry it. Replace the point tip by a
plane parallel to the surface and the real-space map collapses; what is left is
one number per k-point, and on a two-dimensional metal that number is the
answer to "which pocket of the Fermi surface does an electron actually tunnel
out of".

Integrating the point-tip map over one cell's worth of the tip plane turns
$\int a_{n\mathbf k}(\mathbf r)\, a^{*}_{m\mathbf k}(\mathbf r)\,\mathrm d^2 r$
into $g_{n\mathbf k} g_{m\mathbf k} S^{\rm tip}_{\mathbf k}[m,n]$, where
$S^{\rm tip}$ is the \emph{same} Gram matrix as the exit plane's, evaluated at
the tip's height. So
\begin{equation*}
  W(\mathbf k) \;=\; w_{\mathbf k}\,
    \mathrm{Tr}\!\left[ D_{\mathbf k}\, S^{\rm exit}_{\mathbf k}\,
                        D_{\mathbf k}\, S^{\rm tip}_{\mathbf k} \right],
  \qquad D_{\mathbf k} = \mathrm{diag}\bigl(g_{n\mathbf k}\bigr),
  \qquad \sum_{\mathbf k} W(\mathbf k) = \int_{\rm tip} T(\mathbf r)\,\mathrm d^2 r .
\end{equation*}
That last equality is a theorem, not an approximation: this is the same object
as the map, and nothing is sampled in real space, so it costs a small fraction
of one --- two Gram matrices and an $n_{\rm bnd}^2$ trace, against an
$n_k \times n_{\rm bnd} \times n_{\rm points}$ sampling.

\textbf{Both sides of it are Brillouin-zone integrals, so it holds mesh for
mesh.} Passing \code{grid=} re-solves the bands on the grid you name, which is
what the refusals below require; comparing the result against a quantity built
on the k-points the SCF converged on then compares two different integrals of
the same integrand. \code{run\_stm} is that case and has no \code{grid} of its
own, so a comparison against it takes no \code{grid} here either. On the
hydrogen sheet of the test suite, whose SCF runs a $4\times4\times1$ grid,
$\sum_{\mathbf k}$ \code{tersoff\_hamann} is the image's plane integral to
$1.2\times10^{-13}$ on the SCF's own k-points and a factor of $7.7$ out on a
$3\times3\times1$ one.

\textbf{Three columns come back together}, because they are the same trace with
one or both matrices replaced by the identity, and \emph{the physics is in their
ratio}. \code{weight} is $W(\mathbf k)$; \code{tersoff\_hamann} is its
$S^{\rm exit}\!\to\!1$ limit, which is the tip plane's local density of states
and so what an STM would see; \code{bare} is the further limit in which the tip
plane stops discriminating too, the plain Fermi surface
$\sum_n \tilde\delta(E - \varepsilon_{n\mathbf k})$. A fourth,
\code{incoherent}, is every channel tunnelling on its own, taken in the basis
that diagonalises $S^{\rm exit}$ inside each degenerate multiplet --- that
column is a \emph{diagonal} and so is not invariant under the rotation a
degenerate eigensolver is free in, where $W$ is a trace and needs no such care.

\textbf{Why the ratio is large.} A state at in-plane momentum
$\mathbf k_\parallel$ decays into the vacuum as
$\psi \sim e^{-\sqrt{\kappa_0^2 + k_\parallel^2}\, z}$, so a zone-corner pocket
is suppressed against a zone-centre one by a factor that grows exponentially
with the gap. That is also the \textbf{check that shares no machinery with
either code}: fit a decay constant to each pocket over a height sweep and they
must satisfy $\kappa_1^2 - \kappa_2^2 = |\mathbf k_1|^2 - |\mathbf k_2|^2$.
$\kappa$ is the \emph{amplitude} decay and $W$ is quadratic in $\psi$ on each
plane, so a sweep of the tip alone has $\mathrm{d}\ln W/\mathrm{d}z = -2\kappa$
and a sweep of both planes $-4\kappa$; halving on one side only makes the
identity come out four times too large, which is what it is for. A height sweep
costs \emph{one} band solve, not one per height --- the bands do not know where
the tip is.

\begin{pycode}
import numpy as np
from defumat import Calculator
from defumat.transport.momentum import decay_constants, decay_identity, pocket_mask

calc = Calculator.from_file("nbse2.in")          # a 1H-NbSe2 monolayer

# the two planes straddle the slab; a sweep of tip heights costs one band solve
run = calc.get_momentum_transport(exit_height=0.212, height=[0.677, 0.732, 0.788],
                                  grid=(6, 6, 1), broadening=0.006,
                                  smearing="gaussian")

run.weight            # (nheights, nk) -- the transmission per k-point
run.tersoff_hamann    # the same states, substrate made structureless
run.bare              # the plain Fermi surface
run.interference      # coherent minus incoherent
run.as_grid("weight") # reshaped to the Monkhorst-Pack grid, for a BZ image
run.per_kpoint()      # with the k-weight divided out, for another code

# the zone partition: closer to a corner than to Gamma, which on a hexagon is
# exactly 2/3 of the zone by area
corners = pocket_mask(run.kcartesian, calc.system.cell)
run.shares(corners, index=2)        # each column's share, and the suppression
run.reweighting(corners, index=2)   # (W_G/W_K)/(D_G/D_K): how the junction re-ranks

# the vacuum decay, and the identity that checks it
sweep = run.sweep("tersoff_hamann")
heights = np.array([1.5, 2.5, 3.5]) / 0.529177210903      # bohr
kappa = decay_constants(heights, sweep[:, [0, 14]])       # Gamma and K
decay_identity(kappa[1], kappa[0], run.kcartesian[14], run.kcartesian[0])
\end{pycode}

\textbf{The index order is the whole difficulty, and it needs two distinct
planes to see at all.} The elementwise-looking alternative
$\sum_{nm} g_n g_m S^{\rm exit}[n,m]\, S^{\rm tip}[n,m]$ is
$\mathrm{Tr}[D S^{\rm exit} D (S^{\rm tip})^{\mathsf T}]$: real, non-negative,
and \emph{exactly} equal to the right answer whenever either matrix is diagonal
--- so it agrees in the Tersoff--Hamann limit, agrees on a single-band metal,
and passes every check that does not involve two planes at once. On random
orthonormal bands it is \textbf{3\%} out, which is far above round-off and far
below anything a plot would show. Only the literal plane integral separates
them, and it does: the two routes agree to \textbf{1e-12}, one evaluating a
Green's function on a real-space quadrature of the tip plane and the other never
leaving reciprocal space.

\textbf{Validated against a second implementation, which is new here.} Neither
\code{pw.x} nor stock Elk computes this: Elk's \code{fermisurf.f90} (tasks
100/101) writes bands for a plotting program without weighting them by
anything, and its only tunnelling task is 162, the Tersoff--Hamann image. The
one other implementation known is elkpy's Elk patch (task 9007), and it was
written independently --- same on-shell Bardeen amplitude, same full Gram
matrices on both planes, same trace of a product, and the same index-order trap
recorded on its own side. On 1H-NbSe2 the two agree on the physics; the numbers
are in \code{PLAN.md} P77.

\begin{refuses}
\textbf{Refuses.} A \emph{magnetic plane tip}: a polarized tip contracts the two
spinor components through its own $2\times2$ projector, which the plane integral
does not collapse the way it collapses the trace over position, so it is left to
\code{get\_vertical\_transport}'s map rather than approximated here. The
substrate's own acceptance does go in, inside the Gram matrix, exactly as it
does there. A tilted tip plane, where the point-tip map allows one: the
closed-form Miller-index orthogonality both sides rest on needs a plane spanned
by two lattice vectors, and it is what makes this cheap.
\textbf{Inherited from the vertical transmission unchanged}: a k-set with more
than one division along the stacking axis, a symmetry-reduced k-set
(\code{grid=} builds the whole grid), a plane inside an augmentation sphere on
an ultrasoft or PAW dataset, gamma-only storage, a spin-selective substrate on a
run with no magnetization, and --- from the STM images --- a spin spiral, a
constrained \code{tot\_magnetization} and an applied magnetic field. Both planes
on the same side of the slab \textbf{warns} rather than refuses.
\textbf{Not claimed}: an absolute conductance. The lead couplings are unfixed
prefactors, so what this carries is the weights and their ratios --- which is
also why \code{shares} reports shares of a total rather than absolute weights.
\end{refuses}

\section{Things with no \texorpdfstring{\code{pw.x}}{pw.x} counterpart}

\textbf{Spin spirals} by the generalized Bloch theorem: a spiral of wavevector
$\mathbf q$ is a spinor whose two components live on \emph{different} plane-wave
spheres,
\begin{equation*}
\psi_{\mathbf k}(\mathbf r) = \begin{pmatrix}
  e^{i(\mathbf k - \mathbf q/2)\cdot\mathbf r}\,u^{\uparrow}(\mathbf r) \\[3pt]
  e^{i(\mathbf k + \mathbf q/2)\cdot\mathbf r}\,u^{\downarrow}(\mathbf r)
\end{pmatrix}
\end{equation*}
so in the rotated frame the density is lattice periodic and the SCF, the
functional and the mixer are untouched. On a norm-conserving dataset only two
terms carry $\mathbf q$ --- $|\mathbf k \mp \mathbf q/2 + \mathbf G|^2$ and
$v_{\mathrm{nl}}$ --- and $dE/d\mathbf q$ at frozen coefficients is
\code{jax.grad} of them.

\textbf{An ultrasoft or PAW dataset adds a third, and it is the one that used to
refuse the combination.} The transverse block of the density is
$\langle\psi^{\uparrow}|\ldots|\psi^{\downarrow}\rangle$, so it pairs projectors
at $\mathbf k + \mathbf q/2$ with projectors at $\mathbf k - \mathbf q/2$, and
the augmentation charge reaching it is the lattice sum
$\sum_{\mathbf R} e^{-i\mathbf q\cdot\mathbf R}\,
Q_{ij}(\mathbf r - \boldsymbol\tau_a - \mathbf R)$
rather than the periodic one. That sum is $e^{-i\mathbf q\cdot\mathbf r}$ times
a lattice-periodic function --- the same factor the smooth transverse density
carries, which is what makes the rotated frame work at all --- and the periodic
factor is the transform of
\begin{equation*}
Q_{ij}(\mathbf G - \mathbf q)\, e^{-i(\mathbf G - \mathbf q)\cdot\boldsymbol\tau_a},
\end{equation*}
the resident table displaced by $-\mathbf q$. The charge and $m_z$ blocks pair
projectors at the \emph{same} $\mathbf k$ and keep the resident table. Nothing
else changes: the displaced table is built by the same assembly, and at
$\mathbf q \to 0$ it is the resident one term for term.

\textbf{$dE/d\mathbf q$ on such a dataset differentiates three terms the
norm-conserving one does not have.} The displaced table is a function of
$\mathbf q$ exactly as $|\mathbf k \mp \mathbf q/2 + \mathbf G|^2$ and
$v_{\mathrm{nl}}$ are, so it is rebuilt inside the differentiated path rather
than held fixed; the orthonormality constraint the frozen-coefficient gradient
carries is $\langle\psi|S|\psi\rangle - 1$ rather than
$\langle\psi|\psi\rangle - 1$, and $S$ pairs each spinor component with the
projectors of its \emph{own} shifted sphere, so it moves with $\mathbf q$ and
its derivative is the spiral's Pulay term; and PAW's one-centre energy is part
of the total being differentiated. The last of those is the one a test would
have caught immediately --- without it the functional misses the whole of
$E_{\mathrm{PAW}}$ --- and the first two are the ones that would have been
silent, which is why each is measured separately below.

\textbf{What it agrees with, and what PAW's own floor is.} Against a central
difference of the \emph{re-converged} energy at a frozen plane-wave sphere, on
the oxygen chain at $q_3 = 0.3$, an ultrasoft dataset converges onto the finite
difference as $\delta^2$ and reaches $2.4\times10^{-7}$~Ry per unit $\mathbf q$
at $\delta = 0.0025$. PAW stops falling at $1.4\times10^{-6}$, and that floor is
\emph{not} the spiral's.

The sharpest way to see it needs no finite difference at all: at $q_3 = 0$ and
$q_3 = \tfrac12$ symmetry forces $dE/d\mathbf q$ to vanish exactly, so what
comes back there \emph{is} the gradient's own error --- $4$ to
$6\times10^{-10}$ at the zone centre for both datasets, and $8.5\times10^{-8}$
(ultrasoft) against $2.5\times10^{-6}$ (PAW) at the zone boundary. That residue
is insensitive to \code{conv\_thr} and to the mixing and falls by a factor of
\textbf{seven} when \code{ecutrho} is doubled, the finite-difference residual
falling by three over the same step since it also carries the difference's own
truncation. And the same cell with no spiral in it at all shows a gap of the
same size between the collinear and noncollinear paths, which \code{pw.x}
reproduces to every digit it prints. So it is what PAW's spherical terms cost
against the grid, inherited here rather than introduced.

Two consequences are worth knowing rather than discovering. The transverse pair
is augmented as one \emph{complex} field, because $Q_{ij}(\mathbf r)$ is real
and so the resident table satisfies
$Q_{ij}(-\mathbf G) = Q_{ij}(\mathbf G)^{*}$ while the displaced one does not;
and the identity against a supercell closes to $1.7\times10^{-9}$~Ry on an
ultrasoft dataset and $3.3\times10^{-7}$ on a PAW one --- the gap between those
two being PAW's own one-centre path, which shows an inconsistency of the same
size between this code's collinear and noncollinear routes on a cell with no
spiral in it at all, rather than anything the displacement does.

\begin{pycode}
from defumat.workflows.spiral import (run_spiral_scan, relax_spiral_q,
                                       heisenberg_exchange)

scan = run_spiral_scan(system, pseudos, wavevectors)   # E(q)
best = relax_spiral_q(system, pseudos)   # BFGS on the reciprocal metric
print(best.wavevector, best.scf.total_energy)

J = heisenberg_exchange(scan, system.cell, shells)   # fit E(q) for the J(R)

scan = run_spiral_scan(system, pseudos, wavevectors, gradients=True)
print(scan.relative)     # E(q) - E(q_0) in mRy, from the energies
print(scan.integrated)   # the same curve, from dE/dq along the path
\end{pycode}

\code{heisenberg\_exchange} fits
$E(\mathbf q) - E(0) = m^2 \sum_{\mathbf R} J(\mathbf R)\,[1 - \cos(\mathbf q\cdot\mathbf R)]$
over a set of lattice shells, which is how a spiral scan becomes a spin model.
\code{compute\_spiral\_gradient} gives $dE/d\mathbf q$ on its own.

\textbf{$E(\mathbf q)$ two ways.} With \code{gradients=True} the scan also takes
$dE/d\mathbf q$ at every point, and \code{scan.integrated} accumulates it along
the path,
\begin{equation*}
E(\mathbf q_n) - E(\mathbf q_0)
  = \sum_m \int \mathrm d\mathbf q\cdot\frac{\partial E}{\partial\mathbf q},
\end{equation*}
by the trapezoid rule on the segments between successive wavevectors (both
factors in lattice coordinates, so the contraction needs no metric and the path
may bend). The two curves are the same physics and the gap between them is the
useful part: the energies are computed on a plane-wave sphere \emph{rebuilt} at
every point and step by the Pulay error of a finite basis wherever a plane wave
crosses the cutoff, while the gradient is taken at a \emph{frozen} sphere and
does not see those steps. On the hydrogen chain of
\code{tests/data/qe/h-chain-spiral.in} at \code{ecutwfc = 25} the direct
energies rise, fall and rise again over a 13-point scan where the integrated
curve descends monotonically; the two agree to $0.005$~mRy at
\code{ecutwfc = 60} once the quadrature error is removed, against $0.130$ at
$25$.

Two error sources live in that gap and refining tells them apart: the trapezoid
rule's own $h^2$ shrinks when the $\mathbf q$ path is refined ($0.051 \to 0.016
\to 0.003$~mRy for $7 \to 13 \to 25$ points) and the basis noise does not
($0.139$, $0.138$, $0.137$). What the integrated route does \emph{not} buy is
worth stating, since the intuition that it might is a common one: it is not
cheaper (every point still needs its own SCF), it does not converge on a coarser
$k$-mesh (the gradient is the exact derivative of the \emph{same} fixed-mesh
energy, so it inherits the same Brillouin-zone error), and it needs a
\emph{tighter} \code{conv\_thr} rather than a looser one, an energy's error
being second order in the density's where a derivative's is first.

Also here: \textbf{per-atom magnetic fields}, the \textbf{topological
invariants} of the previous section, and \textbf{rank-$n$ tensor
symmetrisation} --- \code{symme.f90}'s \code{symmatrix3}/\code{symtensor3}
written at any rank.

\begin{refuses}
\textbf{Refuses.} A spiral with spin-orbit coupling, permanently --- it breaks
the theorem, and Elk refuses it too; a spiral with symmetry, until the spin
space group is written, so a spiral needs \code{nosym}; $dE/d\mathbf q$ on a
dataset whose augmentation charge is held as a \emph{radial} table rather than
on the dense $\mathbf G$ set --- the branch a large cell falls onto when
$n_h^2 n_{\mathrm{gm}}$ is over \code{DEFUMAT\_AUG\_MAX\_BYTES} --- because that
route reads $|\mathbf q|$ on the host to size its interpolation and running past
the end of the table is a NaN rather than a clamp, so a tracer cannot go through
it at all; and spiral relaxation under a magnetic field, whose energy is outside
the reported total, so the state is stationary for a different functional than
the one being differentiated.

\textbf{Warns rather than refuses}: $dE/d\mathbf q$ on an augmented dataset is
evaluated in \emph{one} pass over the whole $k$ axis whatever \code{k\_batch}
asks for, because the density carries $\mathbf q$ there and the Hartree energy
is quadratic in it, so a sum of per-chunk gradients is not the gradient. The
peak then scales with $n_k$ rather than with the chunk size, and it is the one
place $dE/d\mathbf q$ costs more than the SCF it follows: 11.4~GB on the
one-atom oxygen chain of \code{tests/data/qe/o-chain-spiral-us.in} at
\code{ecutrho = 200} with four $k$-points, against a few hundred megabytes for
the SCF.
\end{refuses}

% =====================================================================
\section{A modulation over many unit cells: the ultracell}

A spin density wave in chromium has a period of 21 unit cells. A charged
impurity in silicon is screened over about 20. A skyrmion is hundreds across.
Each of them is a slow \emph{envelope} on a crystal that is still atomically
periodic, and a supercell pays the same price per degree of freedom for the
envelope as for the atoms --- which is why that length scale is the one nobody
computes from first principles.

The ultracell separates the two. Take $N = n_1n_2n_3$ copies of the unit cell.
The reciprocal lattice vectors of that ultracell lying inside the unit cell's
Brillouin zone are $N$ wavevectors $\mathbf Q$, and everything allowed to be
long ranged is written
\begin{equation*}
\rho(\mathbf r) = \sum_{\mathbf Q} \rho_{\mathbf Q}(\mathbf r)\,
   e^{i\mathbf Q\cdot\mathbf r},
\qquad \rho_{\mathbf Q}\ \text{lattice periodic}.
\end{equation*}
The ultracell's Kohn-Sham states at an ultracell Brillouin zone point
$\mathbf k$ are then expanded in the \emph{unit cell's own} Kohn-Sham states at
the $N$ points that fold onto it,
\begin{equation*}
|\Psi\rangle = \sum_{\mathbf Q, n} c_{\mathbf Q n}\,
  |\psi_{\mathbf k + \mathbf Q, n}\rangle, \qquad n \le \code{nbnd},
\end{equation*}
and what is diagonalised is the $N\,\code{nbnd} \times N\,\code{nbnd}$ matrix
\begin{equation*}
H_{(\mathbf Q m),(\mathbf Q' n)} =
  \delta_{\mathbf Q \mathbf Q'}\delta_{mn}\,\varepsilon_{\mathbf k+\mathbf Q,n}
  + \langle \psi_{\mathbf k+\mathbf Q, m}|\,\delta V\,
    |\psi_{\mathbf k+\mathbf Q', n}\rangle ,
\end{equation*}
with $\delta V$ the ultracell potential \emph{minus} the unit cell's --- the
eigenvalue on the diagonal already carries the whole unmodulated potential, so
putting it in twice would be the one easy mistake here.

\textbf{Why this is cheap.} The unit-cell states at the folded k-points are
computed \textbf{once, before the self-consistent loop}, and frozen. Every
iteration after that builds a small dense matrix and diagonalises it; no
plane-wave Hamiltonian is touched again. The variational dimension per k-point
is $N\,\code{nbnd}$ where a real supercell's is $N\,\code{npw}$, and the ratio
$\code{npw}/\code{nbnd}$ is typically $10^2$ to $10^3$. It is a prefactor,
however, not a change of scaling: the dense solve is still
$O((N\,\code{nbnd})^3)$, so this reaches a few thousand cells along one
direction rather than a micron in three.

\textbf{What the approximation is, exactly.} As \code{nbnd} grows to the full
plane-wave count, the basis above spans precisely the ultracell's own
plane-wave basis. So the method is a \emph{variational truncation of the exact
$N$-cell supercell problem to \code{nbnd} bands per folded k-point}, and
nothing else. An eigenvalue's convergence carries no sign --- the Hamiltonian
moves with its own truncated density, so an ultracell eigenvalue is not an
upper bound on the supercell's, and neither is a density --- but it does
converge, which is the property the implementation is checked against: a
two-cell silicon ultracell under an applied modulation reproduces a real
four-atom supercell's induced density, and the disagreement falls monotonically
with \code{nbnd} (\code{tests/regression/test\_ultracell.py}).

\textbf{The total energy is the exception, and it does carry a sign.} What the
run reports as \code{total\_energy} is the Kohn-Sham free energy of the state
the loop converged to, and the bases are \emph{nested} in \code{nbnd}, so it is
a monotone non-increasing upper bound on the $N$-cell supercell's own energy:
$+8.15\times10^{-5}$, $+3.52\times10^{-6}$, $+3.62\times10^{-7}$ and
$+1.43\times10^{-7}$~Ry at \code{nbnd = 12, 24, 48, 64} on that same silicon
cell. Comparing against a supercell needs the \emph{same} FFT box on both
sides, since two discretisations of one functional differ by about
$10^{-6}$~Ry per cell here, which is larger than the top of that ladder; within
the ladder the box is the same one and the monotonicity is exact.

\begin{pycode}
from defumat import Calculator
import numpy as np

cal = Calculator.from_file("si.in")        # the two-atom unit cell

# a potential of one ultracell period along a_1, in Ry; x runs over [0, n_i)
modulation = lambda x: 0.05 * np.cos(2 * np.pi * x[..., 0] / 8)

ulr = cal.get_ultracell(supercell=(8, 1, 1), kgrid=(1, 2, 2),
                        nbnd=24, external=modulation)

print(ulr.iterations, ulr.accuracy)    # dr2 on the ultracell box, in Ry
print(ulr.density.shape)               # (nspin, 8*n1, n2, n3): the whole ultracell
print(np.abs(ulr.modulation).max())    # rho - tile(rho_0): the induced envelope
print(ulr.fermi_energy, ulr.band_energy)
\end{pycode}

\code{supercell} is three integers and the ultracell lattice vectors are derived
from them; \code{kgrid} samples the \emph{ultracell's} Brillouin zone, which is
$N$ times smaller than the unit cell's, so the folded set the frozen states are
computed on is \code{supercell * kgrid}. \code{external} is an applied potential
in Ry, given either as an array on the ultracell grid or as a function of
unit-cell crystal coordinates, which run over $[0, n_i)$ across the ultracell.
\code{ulr.density} is normalised \emph{per unit cell}, so it integrates to $N$
times the electron count over the ultracell and equals the tiled unit-cell
density exactly when nothing is modulated.

\textbf{\code{nbnd} is the convergence parameter and it is the only one.} It is
the size of the variational basis per folded k-point, so the default --- which
for an insulator is the occupied bands and nothing above them --- leaves the
envelope nothing to be built from. Pass it. \code{ulr.multiplet\_gap} reports
the smallest gap the truncation opens anywhere on the folded k-set: where that
is small the cut has fallen inside a degenerate multiplet, which member of it
survives is whatever the eigensolver returned, and the basis is arbitrary
there. \textbf{Whether that matters depends on where the cut is}, and the two
cases were measured rather than assumed. Low --- no more empty bands than
occupied ones --- the arbitrary rotation is among the very states the response
is built from, and \code{run\_ultracell} warns: \code{nbnd = 8} on silicon cuts
at $2.6\times10^{-14}$~Ry and does not converge at all. High in the empty
manifold a cut this degenerate is both common and harmless --- \code{nbnd = 32},
\code{48} and \code{80} each open one ($6\times10^{-15}$, $5\times10^{-12}$,
$4.7\times10^{-11}$~Ry) and
each converged in ten iterations to a monotonically improving answer --- so the
gap is reported there without a warning. Raise \code{nbnd} until the cut lands
in a gap, or until it lands high enough not to be one.

\textbf{Convergence.} The mixed quantity is the density on the ultracell box and
the convergence test is the Hartree energy of its residual --- QE's \code{dr2},
generalised to $|\mathbf G + \mathbf Q|$ --- so \code{conv\_thr} means here what
it means everywhere else in this package. Kerker preconditioning is on by
default and matters more here than in a unit cell: the smallest non-zero
$|\mathbf G+\mathbf Q|$ an ultracell carries is $N$ times smaller, so the
Hartree kernel there is $N^2$ larger and long-wavelength charge sloshing is the
whole difficulty. Elk's own example of the method runs at a mixing parameter of
$0.001$ for up to two thousand iterations; a two-cell silicon ultracell under an
applied modulation converges here in nine.

\textbf{The total energy, and what it is for.} \code{ulr.total\_energy} is the
Kohn-Sham free energy per unit cell, and it is what says whether a modulation is
the ground state rather than merely what a modulation does: the energy
\emph{gain} of a spin density wave over the uniform state is a difference of two
of these. Neither reference code has it --- Elk's \code{energyulr.f90} computes
the occupied eigenvalue sum alone and \code{pw.x} has no ultracell --- and the
eigenvalue sum is not the total: it double-counts the Hartree and
exchange-correlation energies exactly as an ordinary \code{eband} does, because
the frozen eigenvalue carries the whole self-consistent potential. The
assembly is therefore QE's own,
\[
  E = (\code{eband} + \code{deband}) + E_{\mathrm H} + E_{\mathrm{xc}}
      + E_{\mathrm{Ewald}} + \code{demet},
\]
with \code{deband} the integral of the density against the potential over the
ultracell, and \code{ulr.energy\_terms} carries the pieces. The Ewald term is the
\emph{unit cell's}, unchanged, because the atoms do not move.

\begin{pycode}
# the same run as above -- the energy comes with it and costs nothing extra
print(ulr.total_energy)                # Ry, per unit cell
print(ulr.energy_terms)                # the pieces, summing to it
print(ulr.energy_history[-3:])         # one per iteration, approaching from above
print(ulr.field_energy)                # None without a magnetic_field
\end{pycode}

\textbf{This is the one quantity in the method that converges with a sign.} An
ultracell eigenvalue is not an upper bound on the supercell's and neither is a
density, for the reason above --- the Hamiltonian moves with its own truncated
density. The energy is a bound because the bases are \emph{nested}: raising
\code{nbnd} keeps every function already in the basis, and the union over
$\mathbf Q$ of the folded k-point spheres is the supercell's own plane-wave
space. So \code{total\_energy} falls monotonically towards the real supercell's
from above, and on two-cell silicon under an applied modulation it does:
$+8.15\times10^{-5}$, $+3.52\times10^{-6}$, $+3.62\times10^{-7}$ and
$+1.43\times10^{-7}$~Ry at \code{nbnd = 12, 24, 48, 64}.

\textbf{Comparing against a supercell needs the same FFT box on both sides.}
Two discretisations of one functional differ by about $10^{-6}$~Ry per cell
here, which is larger than the top of that ladder, and a supercell picks its
dense grid from its own cutoff: at the default \code{ecutwfc = 12} it takes
$(32,15,15)$ where the tiled unit cell's is $(30,15,15)$, and the third rung
then comes out $1.49\times10^{-7}$~Ry \emph{below} the supercell. Choose a
cutoff at which the two grids agree --- \code{ecutwfc = 13} puts the unit cell on
$(18,18,18)$ and the two-cell supercell on $(36,18,18)$ --- and do it through
\code{ecutwfc} rather than \code{ecutrho}, since an ultracell refuses a double
grid. Within one ladder the box is the same on every rung and the monotonicity
is exact.

\textbf{An ultrasoft or PAW dataset.} The augmentation charge changes nothing
structural about the method, and saying why is shorter than the code. A matrix
element between two ultracell basis functions runs over the whole ultracell, so
it sums over the $N$ copies of every atom, and each copy's projectors carry the
Bloch phase $e^{i(\mathbf k_0+\mathbf Q)\cdot\mathbf R}$; what survives is
\[
  \sum_{\mathbf R} e^{i(\mathbf Q'-\mathbf Q)\cdot\mathbf R}\,
  Q^a_{ij}(\mathbf r - \boldsymbol\tau_a - \mathbf R),
\]
the ordinary augmentation charge \emph{displaced} by $\mathbf Q' - \mathbf Q$,
the ket's wavevector minus the bra's. That is the same object a spin spiral
already uses, which is its special case at $\mathbf Q = +\mathbf q/2$ and
$\mathbf Q' = -\mathbf q/2$. The overlap operator does \emph{not} enter: the
augmentation part of $\langle\psi_{\mathbf k_0+\mathbf Q}|S|
\psi_{\mathbf k_0+\mathbf Q'}\rangle$ carries the same lattice sum, and
$\mathbf Q - \mathbf Q'$ is a reciprocal vector of the ultracell, so that sum is
$N\delta_{\mathbf Q\mathbf Q'}$ --- the frozen states are exactly
$S$-orthonormal across $\mathbf Q$ and the eigenproblem stays an ordinary one.
The density gains the augmentation charge of every copy, built from a
\code{becsum} resolved by $\mathbf Q$-difference, and PAW gains its one-centre
terms per copy, the copies' occupations being the transform of that same
\code{becsum} over the same index.

\begin{pycode}
from defumat import Calculator
import numpy as np

cal = Calculator.from_file("si-ultracell-paw.in")   # PAW, ecutrho = 4 ecutwfc

ulr = cal.get_ultracell(supercell=(4, 1, 1), kgrid=(1, 2, 2), nbnd=32,
                        external=lambda x: 0.02 * np.cos(2 * np.pi * x[..., 0] / 4))

print(ulr.total_energy)                       # -89.18116879 Ry per unit cell
print(ulr.energy_terms["one_center_paw"])     # -67.17927909, the one-centre term
print(float(np.abs(ulr.modulation).max()))    # 1.7738e-04 e/bohr^3, the response
\end{pycode}

\code{ecutrho = 4 ecutwfc} is not what such a dataset wants, and
\code{run\_ultracell} says so with a warning: the run is then self-consistent
and comparable against a supercell at the same cutoffs while its absolute
energy is not converged in \code{ecutrho}. The dual the dataset was generated
for, 8 to 12 times \code{ecutwfc}, runs as well and is what an absolute number
should be taken at. One thing to know before comparing a run at a dual against
a real supercell: the supercell chooses its own FFT box, which at a dual is not
the tiled one --- $(54, 25, 25)$ against $(50, 25, 25)$ on this cell --- and it
then sits $2.20\times10^{-6}$ Ry per cell above its own unit cell with nothing
applied. Compare the energy of the modulation, each side against its own
unmodulated state, and that offset cancels. Both datasets were measured against a real
four-atom supercell under the same applied potential: the induced density's
error falls $1.18\times10^{-1}$, $1.56\times10^{-2}$, $3.72\times10^{-3}$
ultrasoft and $1.17\times10^{-1}$, $1.71\times10^{-2}$, $5.09\times10^{-3}$ PAW
at \code{nbnd = 12, 24, 48}, with the total above the supercell's at every rung
and falling. With no modulation applied the total reproduces the unit cell's own
SCF to $8.5\times10^{-12}$~Ry and $2.5\times10^{-11}$~Ry, the second including a
$-67.18$~Ry one-centre term.

\textbf{A spinor works the same way, with one line that is not the unit cell's.}
At \code{npol = 2} the four integrals are recombined into $2\times2$ spin blocks
and sandwiched between \code{fcoef} --- \code{newd\_so}, with a wavevector on it
--- and PAW's one-centre coefficients are added to the \emph{scalar} components
before that transform rather than after, which is the order the unit cell uses
and the one that puts them in the right spin structure. What differs is the lower
off-diagonal block. In a unit cell it is the conjugate transpose of the upper one,
because both are built from the same two real components; here every component is
a function of the $\mathbf Q$-difference and
$D_a(\mathbf Q_d)^{*} = D_a(-\mathbf Q_d)$, so that conjugate transpose belongs
to the \emph{opposite} difference. Both blocks are written out from the
components, and the matrix is Hermitian across the pair
$(\mathbf Q_d, -\mathbf Q_d)$ rather than within one entry.

With nothing applied the tiled total reproduces the unit cell's own SCF to
$2.8\times10^{-12}$~Ry ultrasoft, $8.5\times10^{-13}$~Ry PAW, and
$7.5\times10^{-13}$~Ry on fully relativistic platinum, which is the one cell
where \code{fcoef} is not the identity. Against a real supercell under the same
modulation the total sits above it at every rung and falls:
$+4.80\times10^{-5}$, $+4.48\times10^{-6}$, $+6.37\times10^{-7}$~Ry at
\code{nbnd = 24, 48, 96} for spinor PAW silicon, and $+3.09\times10^{-5}$,
$+1.55\times10^{-5}$, $+5.21\times10^{-6}$ at \code{nbnd = 16, 24, 40} for
platinum. \textbf{A spinor band holds one electron where a collinear band holds
two}, so the comparison against the collinear route is at half the band count,
and there the two agree to every digit printed --- $+6.37\times10^{-7}$ on both
at the last rung.

\begin{pycode}
from defumat import Calculator
import numpy as np

cal = Calculator.from_file("pt-ultracell-soc.in")   # fcc Pt, lspinorb, ultrasoft

ulr = cal.get_ultracell(supercell=(2, 1, 1), kgrid=(2, 2, 2), nbnd=40,
                        external=lambda x: 0.05 * np.cos(np.pi * x[..., 0]))

print(cal.system.npol, cal.system.nspin_mag)  # 2 1 -- a spinor carrying no moment
print(ulr.total_energy)                       # -69.34673356 Ry per unit cell
print(ulr.augmentation_residual)              # 2.8266e-07
\end{pycode}

\textbf{One number comes back with an augmented run that a norm-conserving one
does not report.} \code{ulr.augmentation\_residual} is how complex the augmented
density came back, relative to its own largest value, and it is a convergence
number in \code{ecutrho} rather than an error: the ultracell's reciprocal cutoff
set is the unit cell's dense sphere at \emph{every} $\mathbf Q$, and the negative
of $\mathbf G + \mathbf Q$ is that sphere \emph{displaced} rather than that
sphere, so a real function represented on it is not exactly real. It is
\textbf{exactly zero} at $N=1$, and on ultrasoft silicon it reads
$1.71\times10^{-4}$, $7.21\times10^{-5}$ and $2.96\times10^{-5}$ at
\code{ecutrho} of 64, 96 and 144~Ry. The Hartree term of a norm-conserving
ultracell carries the same truncation and does not report it.

\textbf{Two spin channels are in and are measured separately}, since everything
above runs one channel and would not see the other. The null at \code{nspin = 2}
reproduces the unit cell to $7.2\times10^{-13}$~Ry, and a uniform applied field at
$N=1$ --- where each channel gets its own $\delta V$ and therefore its own
$\int \delta V\,Q$ in $D_{ij}$, with opposite signs --- reproduces an ordinary
SCF's moment to $5.0\times10^{-3}$, $1.8\times10^{-3}$ and $4.9\times10^{-4}$
relative at \code{nbnd = 12, 24, 40}. What has \emph{not} been compared against a
supercell is a \emph{modulated} moment on an augmented dataset.

\textbf{An applied field's energy is not in the total}, which is QE's and Elk's
shared convention and this package's everywhere else: \code{ulr.field\_energy}
carries $-\int B\cdot m$ beside it. An applied \emph{scalar} potential is
treated the other way and its energy stays in the total, once, because that is
what \code{vltot} is --- which is also what makes a comparison against a
supercell carrying the same modulation a comparison of the same quantity.

\textbf{Under a field, the quantity that carries the sign is
\code{total\_energy + field\_energy}}, not the total on its own. What the
calculation minimises is the full energy with the Zeeman term in it, and the
reported total has that term taken out, so the total alone is a bound on
nothing --- measured on a partly-polarized hydrogen cell under a uniform
$0.02$~Ry field, it sits $7.4\times10^{-7}$~Ry above an ordinary SCF at
\code{nbnd = 12} and $9.7\times10^{-7}$ and $1.6\times10^{-6}$~Ry \emph{below}
it at 24 and 40, moving monotonically the wrong way, where the sum is
$+4.15\times10^{-6}$, $+2.47\times10^{-6}$ and $+5.15\times10^{-7}$ --- above at
every rung and falling. Add the field energy back whenever two magnetic states
are being compared.

\textbf{A spin density wave.} \code{nspin = 2} is in, and it is what the method
was published for: a modulation of the \emph{magnetization} over tens of cells
is the case a supercell cannot afford. Without spin-orbit coupling the ultracell
matrix is block diagonal in spin, so two channels cost two solves of the same
size rather than one of twice it; what the channels share is a \textbf{single
Fermi level} over all $2N\,\code{nbnd}$ states, and that sharing is the physics
--- it is what lets an electron cross from the minority channel in one cell to
the majority channel in the next, which is what a spin density wave is.

Nothing in a calculation breaks spin symmetry on its own, so a modulated moment
has to come from somewhere --- and the tiled state is an \emph{exact} fixed
point of this loop, so a magnetic ultracell given nothing converges in one
iteration and stays exactly as ferromagnetic as its unit cell. There are two
doors, and \textbf{they are different quantities rather than two routes to
one}. An applied field \emph{drives} a modulation, and what comes back is the
$\mathbf Q$-resolved spin susceptibility: a response. A \emph{seed} hands the
loop the texture as its initial condition and the loop keeps it: the ordered
state itself. The field comes first here because it is the honest measurement of
a susceptibility; the seed is under \emph{Seeding the texture} below.

\code{magnetic\_field} is the first knob: a collinear $B(\mathbf r)$ in Ry,
entering as $v_\uparrow \mathrel{-}= B$ and $v_\downarrow \mathrel{+}= B$, in
the same two forms \code{external} takes.

\begin{pycode}
from defumat import Calculator
import numpy as np

cal = Calculator.from_file("si-mag.in")     # the same cell at nspin = 2

# a field of one ultracell period along a_1, in Ry
field = lambda x: 0.05 * np.cos(2 * np.pi * x[..., 0] / 2)

ulr = cal.get_ultracell(supercell=(2, 1, 1), kgrid=(1, 2, 2),
                        nbnd=12, magnetic_field=field)

print(ulr.cell_moments())          # the moment of each cell: the wave itself
print(np.abs(ulr.magnetization).max())         # rho_up - rho_dw on the box
print(ulr.charge_accuracy, ulr.magnetic_accuracy)
\end{pycode}

\code{ulr.cell\_moments()} integrates $\rho_\uparrow - \rho_\downarrow$ over
each unit cell of the ultracell, which needs no sphere and no radius: the cells
tile the ultracell exactly, so the partition is the geometry. On the cell above
the two come back equal and opposite, which is the two-cell wave the field asked
for.

\subsection*{Seeding the texture}

\code{seed\_magnetization} is the second door. It is a factor
$s(\mathbf r)$ on the \emph{converged unit cell's own} magnetization,
\begin{equation*}
  \mathbf m(\mathbf r) \;\longmapsto\;
  |s(\mathbf r)|\; R\!\left(\hat{\mathbf e}_0 \to \hat s(\mathbf r)\right)
  \mathbf m_0(\mathbf r),
  \qquad |s| \le 1,
\end{equation*}
which reduces to $m \mapsto s\,m_0$ for a collinear run, where $s$ is a number.
That is \code{starting\_magnetization}'s role one level up --- a fraction of a
moment that is already there, with its range --- and it is a factor rather than
a magnetization because a caller knows the texture, one direction per cell, and
does not know the magnetization profile \emph{inside} a unit cell. $s = 1$, and
$s = \hat{\mathbf e}_0$ for a vector, are the identity. $|s| > 1$ is refused: the
reference's $|\mathbf m| \le n$ pointwise, so a factor above one asks for more
magnetization than the cell has and makes a channel density negative.

\textbf{Nothing is faded, and that is the design decision.} Elk seeds a random
\emph{field} (\code{rndbfcu}) and takes it away geometrically
(\code{reducebf}), and it has to: its self-consistency runs on the potential,
where a seed is a term in the functional and has to leave before the converged
state means anything --- which leaves the awkward criterion, a fixed point of a
field that is going to zero. Here the mixed quantity is the \emph{density}, so a
seed is an initial condition and not a term: the functional is the field-free one
from the first iteration onwards, and \code{conv\_thr} means exactly what it
means in every other run in this package. The price of that choice is the other
half of the same sentence --- a seed cannot \emph{hold} a texture the way a field
can, so a state whose basin the seed misses decays back to the tiled one, and the
loop says so by converging to it.

\begin{pycode}
from defumat import Calculator
import numpy as np

cal = Calculator.from_file("h-mag.in")   # partly polarized hydrogen, nspin = 2

# one period of the wave over two cells, as a factor on the cell's own moment
stagger = lambda x: np.cos(np.pi * x[..., 0])

ulr = cal.get_ultracell(supercell=(2, 1, 1), kgrid=(2, 2, 2), nbnd=24,
                        seed_magnetization=stagger)

print(ulr.iterations, ulr.converged)   # 8 True
print(ulr.total_energy)                # -0.93602520 Ry/cell: the staggered state
print(cal.get_scf().total_energy)      # -0.93300778: the ferromagnet it started in
print(ulr.cell_moments())              # [ 0.07381, -0.073736]
\end{pycode}

\textbf{What that buys, as a number.} The seeded run converges to a state
\textbf{3.085~mRy per cell below} the tiled ferromagnet --- a state the method
could not previously reach at all, since the ferromagnet is a fixed point and
nothing moves off it. And it is the right state: against the same two atoms run
as a real supercell through this package's own SCF, seeded the same way with
\code{starting\_magnetization} of opposite sign on two species, the wave's own
Fourier component agrees to $1.51\times10^{-3}$, $5.04\times10^{-4}$,
$1.74\times10^{-4}$ and $3.10\times10^{-5}$ relative at
\code{nbnd = 16, 24, 40, 64}, and the energy sits $+1.53\times10^{-4}$,
$+6.80\times10^{-5}$, $+2.32\times10^{-5}$ and $+5.57\times10^{-6}$~Ry above
the supercell's, above at every rung and falling. The supercell's amplitude is
$0.36546537\,\mu_{\mathrm B}$ per cell.

\textbf{The seed is evaluated pointwise, and a smooth seed is not a per-cell
sign.} $\cos(\pi x)$ changes sign halfway through each cell, so the
\emph{cell} moments it leaves are small ($0.0738$ against the atoms'
$0.365$) while the \emph{atoms} are cleanly opposite --- which is the physical
state, and the supercell agreement is what says so. Read a seeded wave off its
Fourier component, or off \code{cell\_moments()} only for a seed that is
constant within each cell, which is \code{floor(x)} in the callable. Turning
pointwise rather than in steps is also what makes a helix's pitch exact, below.

\textbf{Two refusals, and both would otherwise be silent.} A reference with
\textbf{no moment}: the seed scales a moment, so on a cell that converged
unpolarized it scales zero, every iteration is a no-op and the unpolarized state
comes back as a converged answer. And, for a vector seed, a \textbf{compensated}
reference --- an antiferromagnetic unit cell, whose net moment is a residue of
cancellation whose direction is round-off, so there is no direction to turn the
texture rigidly from. Turning such a cell sublattice by sublattice is not
written.

\textbf{An ultrasoft or PAW dataset is seeded in both of its mixed variables.}
Most of a transition metal's moment is inside the projector spheres rather than
on the grid, so a seed that reached the density alone would seed almost nothing
on a PAW magnet and would leave the first iteration's one-centre potential
disagreeing with its own density. \code{becsum} is therefore seeded from the
\emph{same} evaluated field, read at each atom's own grid point of the box, so
the spheres and the grid cannot disagree about the period. That half is checked
on constructed occupations rather than in a converged run
(\code{tests/unit/test\_ultracell\_seed.py}), because the copy index is a
lattice translation in C-order over the integer triple and a seed landing on the
wrong copy is a converged, plausible and wrong modulation that no tiled null can
see --- a tiled seed being the same on every copy.

\textbf{Read \code{magnetic\_accuracy}, not \code{accuracy}, on a magnetic
run.} \code{dr2} is the Hartree energy of the density residual, so its charge
half is weighted by $1/|\mathbf G+\mathbf Q|^2$ and its magnetization half by a
constant --- and in an ultracell the gap between them is wider than in a unit
cell, because the smallest $|\mathbf G + \mathbf Q|$ is $N$ times smaller and
the charge half's weight at the envelope's own wavevector is $N^2$ larger. A
converged total can therefore leave the moment still moving; the two halves are
reported apart for exactly that reason, and \code{ulr.history} carries both per
iteration.

\textbf{What a magnetic run is checked against.} At $N = 1$ a \emph{uniform}
applied field is the same physics as an ordinary SCF with \code{B\_field(3)},
which shares nothing with this code path --- it goes through
\code{add\_bfield.f90}'s expression inside a plane-wave SCF. On the silicon cell
at $B = 0.02$~Ry the reference moment is $0.71159883$ and the ultracell gives it
to $9.6\times10^{-3}$ relative at \code{nbnd = 12}, $3.2\times10^{-3}$ at 24 and
$7.7\times10^{-4}$ at 40: the same convergence in \code{nbnd} the charge shows,
now for the moment and against a route with no ultracell in it. Against a real
two-cell supercell of polarized hydrogen under a scalar modulation, the charge
and the magnetization agree to $1.7\times10^{-4}$ and $1.4\times10^{-4}$ at
\code{nbnd = 24}.

\textbf{A texture, not just an amplitude.} \code{nspin = 4} is in as well, and it buys
the modulation a collinear run cannot write down: one in which the moment's
\emph{direction} turns from cell to cell rather than its length growing and
shrinking along a fixed axis. A helix, a cycloid, a domain wall and a skyrmion
are all of that kind. \textbf{Spin-orbit coupling comes with it and costs
nothing}: it is a term in the Hamiltonian the frozen unit-cell states were
diagonalised with, and the ultracell only ever sees their eigenvalues and their
coefficients, so there is no spin-orbit term anywhere in the method.

What changes is the shape of one matrix. A collinear potential is diagonal in
spin, so the $2N\,\code{nbnd}$ problem is two blocks of half the size; a spinor
basis function has two components and the potential acts on it as the $2\times2$
matrix
%
\begin{equation}
V(\mathbf r) \;=\; v_0(\mathbf r)\,\mathbb 1 \;+\; \mathbf B(\mathbf r)\cdot\boldsymbol\sigma ,
\end{equation}
%
whose off-diagonal entries couple the components at every point --- and that
coupling is exactly what lets the magnetization turn. There is one block, one
Fermi level over all $N\,\code{nbnd}$ states, and each state holds \emph{one}
electron rather than two.

The applied field changes shape with the regime: \code{magnetic\_field} is a
scalar $B(\mathbf r)$ for \code{nspin = 2} and a \textbf{vector} one for
\code{nspin = 4}, either as a \code{(3, *box)} array or as a function of
unit-cell crystal coordinates returning three components per point. A field that
\emph{turns} is what drives a helix, and it cannot be written down at all in the
collinear form.

\begin{pycode}
from defumat import Calculator
import numpy as np

cal = Calculator.from_file("h-noncolin-ultracell.in")  # moments free to turn

# a field of 0.01 Ry turning once over the ultracell, in the xy plane
def turning(x):
    phase = 2 * np.pi * x[..., 0] / 4
    return 0.01 * np.stack([np.cos(phase), np.sin(phase),
                            np.zeros_like(phase)], axis=-1)

ulr = cal.get_ultracell(supercell=(4, 1, 1), kgrid=(1, 2, 2), nbnd=16,
                        magnetic_field=turning, mixing_beta=0.3)

print(ulr.cell_moments())         # (4, 3): one moment VECTOR per cell
print(ulr.magnetization.shape)    # (3, *box): m_x, m_y, m_z on the grid
\end{pycode}

\code{cell\_moments()} returns an $(N,3)$ array for a noncollinear run where it
returns $(N,)$ for a collinear one, and \code{magnetization} likewise becomes the
three cartesian components on the box. Neither needs a sphere or a radius: the
cells tile the ultracell exactly, so the partition is the geometry. On the cell
above the four moments come back at $27^\circ$, $62^\circ$, $-41^\circ$ and
$-48^\circ$ in the plane the field turns in, with nothing out of that plane ---
so they follow the field, but by less than the $90^\circ$ per cell it asks for.
That is not a convergence failure: a hydrogen lattice at $5.5$~bohr is a close-
packed ferromagnet whose inter-cell exchange is far larger than a $0.01$~Ry
field, and what the run is measuring is the stiffness of the magnet against a
long-wavelength twist. A texture that turns freely wants a crystal whose
modulated axis is weakly coupled.

\textbf{A noncollinear ultracell converges slowly when the field is weak, and
lowering \code{mixing\_beta} is the wrong response.} Without spin-orbit coupling
or anisotropy, turning \emph{every} moment together by the same angle costs no
energy, so the $\mathbf Q = 0$ transverse component of the magnetization has no
restoring force and the self-consistent solution sits in a nearly flat manifold
rather than at a point. The mixer has to \emph{traverse} that manifold to reach
the place the field pins, and \code{mixing\_beta} is how fast it travels.

On the four-cell hydrogen ultracell above, every run given 300 iterations: at a
$0.01$~Ry turning field every \code{beta} converges to the same state and the
best is $0.5$ (36 iterations, against 48 at $0.7$ and 115 at $0.2$). At
$0.002$~Ry --- a field five times weaker, pinning the direction five times less
--- \code{beta = 0.7} converges in \textbf{263} iterations and \code{beta =
0.3} does \emph{not converge in 300}, its residual sitting between $2\times
10^{-5}$ and $9\times10^{-5}$ for the whole run without ever leaving. The small
\code{beta} is not stable where the large one is unstable; it is \textbf{stuck}.

So a slow noncollinear run wants more iterations, or a stronger field to pin the
direction with --- not a smaller \code{mixing\_beta}. The run says so by name
when it fails to converge.

\textbf{What a noncollinear run is checked against.} At $N=1$ a uniform vector
field is the same physics as an ordinary SCF with \code{B\_field(1:3)}, which
goes through \code{add\_bfield.f90}'s expression inside a plane-wave SCF and
shares nothing with this code path --- and for a vector field that check reaches
the transverse components, which no collinear identity can. On a hydrogen cell
whose moment is free, that reference run says something clean on its own: the
converged moment's direction equals $\hat{\mathbf B}$ to \emph{six decimals},
which is what a magnet with no anisotropy does and what fixes the sign of the
coupling. Against it, at $\mathbf B = (0.005, 0.0005, 0)$~Ry, the ultracell
reproduces the moment to $1.9\times10^{-2}$ relative at \code{nbnd = 16},
$6.6\times10^{-3}$ at 40, $1.5\times10^{-3}$ at 64, $4.5\times10^{-4}$ at 96 and
$1.6\times10^{-4}$ at 128 --- monotone, and a factor of 116 across the range. What is converging is the
\emph{transverse} component: $m_x$ is right to $1.5\times10^{-3}$ at the first
rung, and $m_y$ walks $0.0601, 0.0634, 0.0691, 0.0694, 0.0732, 0.0740$ towards
the reference's $0.0743$.

\textbf{Expect to need more bands than a collinear run, by about a factor of
two.} A spinor state holds one electron and lives in a space of
$2\,\code{npwx}$, so a given \code{nbnd} is half as complete a basis as the same
number would be for one collinear channel. The ladder above is at 37 per cent of
the basis at \code{nbnd = 128}, which is where stage 1's charge ladder was at
\code{nbnd = 80} on silicon. Convergence is also slower the larger the
\emph{rotation} asked for: a field that turns the moment by 57 degrees converges
in \code{nbnd} several times more slowly than one that tilts it by 6, because a
change of the local spin frame has to be built out of empty states of the
opposite spin channel, which the exchange splitting puts far away.

\textbf{Seeding a texture, and which axis to write the seed about.}
\code{seed\_magnetization} takes a vector here, so the seed turns the moment
rather than scaling it, and a helix of pitch $N$ cells is the case to have in
mind. Such a helix is invariant under a translation by one cell followed by a
spin rotation of $360/N$ about its own axis, and the self-consistent map
commutes with both --- so that sector is closed and a run started in it stays in
it. \textbf{In the truncated problem it is closed only if the basis is}, and the
basis is the unit cell's own spinors: a rotation about the reference's own
magnetization is a \emph{phase} on each of them, since without spin-orbit
coupling they are eigenstates of
$\boldsymbol\sigma\cdot\hat{\mathbf e}_0$, and a rotation about any other axis is
not. There is therefore \emph{one} closed sector and its axis is
$\hat{\mathbf e}_0$, which makes the rule short: \textbf{write the seed about the
direction the reference's moment already points along}, which is to say converge
the unit cell with its moment along the axis the helix is to turn about.

\textbf{What that buys is the iteration count and the frame, not the pitch.} On
four cells of hydrogen at \code{nbnd = 16}, seeded about $z$ with the reference
along $z$, the run converges in \textbf{14} iterations and the pitch comes back
at $90.0001$, $90.0000$, $89.9999$ degrees. The same seed on a reference along
$(1,1,1)/\sqrt3$ takes \textbf{290} iterations and arrives at a helix turning
about $(1,1,1)/\sqrt3$ instead --- and it is the \emph{same state}: after one
global rotation the two runs' cell moments agree to $1.23\times10^{-3}$ on
moments of $0.272$, and their energies to $3.5\times10^{-9}$~Ry. A global spin
rotation costs nothing without spin-orbit coupling, so the unprotected run is
not wrong; it is traversing a flat manifold to reach the frame its own basis
prefers. That is what the 290 iterations are, and it is why the fix is the axis
rather than \code{mixing\_beta} --- the same lesson the weak-field case above
teaches, arriving from the other side. \textbf{The run says so}: a seed whose
directions do not sit on a cone about the reference's magnetization is warned
about by name, with that pair of iteration counts in the message.

\textbf{What the truncation costs instead is the canting}, and that is the half a
pitch check cannot see. The converged moments stand off the helix plane by a
\emph{uniform} angle --- which the closed sector allows, a cone being as
invariant under the pair as a flat helix is --- and that is a ferromagnetic
remnant at $\mathbf Q = 0$, towards the direction the basis was built around,
which the wave's own Fourier component is blind to. It falls with \code{nbnd}
where the pitch has nothing left to converge:

\begin{center}
\begin{tabular}{lllll}
\code{nbnd} & pitch error & cone & net moment & $E$ above the supercell \\
\hline
16 & $1.35\times10^{-4\,\circ}$ & $10.64^\circ$ & $0.0628$ & $+4.37\times10^{-4}$ \\
24 & $8.50\times10^{-5\,\circ}$ & $6.30^\circ$ & $0.0370$ & $+2.54\times10^{-4}$ \\
40 & $3.27\times10^{-5\,\circ}$ & $3.00^\circ$ & $0.0176$ & $+1.22\times10^{-4}$ \\
64 & $1.20\times10^{-5\,\circ}$ & $0.91^\circ$ & $0.0054$ & $+3.74\times10^{-5}$ \\
\end{tabular}
\end{center}

so reading the pitch alone at \code{nbnd = 16} would call the state exact while
18 per cent of its moment stood off the helix. Across that range the cone falls
by $11.65$ and the energy error by $11.68$, which is worth recording and is not
explained here.

\textbf{On a real magnet this is the number to watch, and it is not cosmetic.} A
seeded 15-cell helix on a fully relativistic PAW NiBr2 cell at \code{nbnd = 40}
came back with a pitch of $24.51^\circ$ against the $24.0^\circ$ asked for and
nothing measurable out of the plane, and with a 22.8 per cent remnant along the
reference direction that the pitch could not see: the moment's \emph{length} then
ran 9 per cent with the texture's own period, which reads as an amplitude wave and
is a uniform vector added to a helix. Removing it took the pitch to
$24.18^\circ$. What it cost is the harmonic content, which is usually the reason
for the calculation --- a coplanar helix of uniform length has no charge response
at $\mathbf q$ at all, and that run's charge came out at $\mathbf q$ over
$2\mathbf q$ by a factor of two, because a remnant along the reference feeds the
charge at $\mathbf q$ directly. So read \code{cell\_moments().mean(axis=0)} beside
the pitch, and treat it as a convergence parameter in \code{nbnd} rather than as
a property of the material. (Measured in another session and reported here.) The reference is the same four
atoms as a real supercell with the helix seeded per species, whose site moments
are flat to $10^{-5}$ and whose energy is $0.739$~mRy per cell \emph{below} the
ferromagnet the ultracell expands around --- which is what makes reaching it
worth a seed at all. At \code{nbnd = 8} the pitch still comes back to
$4.2\times10^{-3}$ degrees, the closure not needing convergence to hold, while
the run does not reach \code{conv\_thr} in 300 iterations and its energy sits
$+9.06\times10^{-4}$~Ry above the supercell.

\textbf{With \code{lspinorb} the closure fails for every axis}, because a
spin-orbit Hamiltonian's eigenstates are not eigenstates of
$\boldsymbol\sigma\cdot\hat{\mathbf n}$ for any $\hat{\mathbf n}$. The pitch is
then not protected and it is not supposed to be --- the anisotropy makes
different pitches inequivalent, so a texture that relaxes is physics rather than
error --- but the two cannot be separated by reading the pitch, and the
comparison to make is against a seeded supercell at the same $N$. \textbf{The
iteration count above may also not be charged there}, because spin-orbit coupling
\emph{gaps} the flat manifold and leaves nothing to wander along: a 15-cell PAW
NiBr2 helix seeded about $z$ with its reference along $x$ --- the expensive
arrangement --- converged in 54 iterations rather than anything like 290. One cell
is not a scaling, so read the warning as advice and not as a forecast.

\begin{pycode}
from defumat import Calculator
import numpy as np

cal = Calculator.from_file("h-noncolin-z.in")  # the moment along z, free to turn

# 90 degrees per cell, turning about the direction the moment already points in
def helix(x):
    phase = 2 * np.pi * x[..., 0] / 4
    return np.stack([np.cos(phase), np.sin(phase),
                     np.zeros_like(phase)], axis=-1)

ulr = cal.get_ultracell(supercell=(4, 1, 1), kgrid=(1, 2, 2), nbnd=24,
                        seed_magnetization=helix, mixing_beta=0.3)

m = ulr.cell_moments()
print(ulr.iterations)                              # 14
print(np.degrees(np.arctan2(m[:, 1], m[:, 0])))    # 35.32 125.32 -144.68 -54.68
print(np.linalg.norm(m, axis=1))                   # 0.33687 on all four
print(m.mean(axis=0))                              # [0, 0, 0.036926]: the remnant
\end{pycode}

The four directions are $90.0000$ degrees apart and the offset is not a defect:
the seed is evaluated pointwise, so the helix turns continuously across a cell
and a cell's own moment sits at the average of the phases inside it. The
\emph{differences} are the pitch.

\begin{refuses}
Gradient-corrected functionals: the gradient of the
ultracell density carries the envelope's own gradient, which a per-cell
evaluation drops (Elk takes that approximation silently; it is refused here).
The \emph{whole-cell} exit region of a transmission
(\code{run\_ultracell\_transport(exit\_region="volume")}) on an augmented
dataset, which is the one Gram matrix that is $\langle\psi|S|\psi\rangle$ ---
the transmission itself runs, as do the image and the spectrum, because all
three sample the state in the \emph{vacuum}, where a pseudo-wavefunction is the
true one, and all three refuse a plane inside an augmentation sphere as the unit
cell's own versions do.
Meta-GGA. A \code{seed\_magnetization} of $|s| > 1$, which asks for more
magnetization than the cell has and makes a channel density negative --- it is
a fraction of the reference's own moment, with
\code{starting\_magnetization}'s range. A seed on a reference with \textbf{no
moment}, since scaling zero leaves every iteration a no-op and returns the
unpolarized state as a converged answer; and a \emph{vector} seed on a
\textbf{compensated} reference, an antiferromagnetic unit cell whose net moment
is a residue of cancellation whose direction is round-off, so the texture has no
direction to be turned rigidly from --- turning one sublattice by sublattice is
not written. An applied \code{magnetic\_field} on a noncollinear cell that
\textbf{carries no magnetization} --- a spin-orbit run on a nonmagnetic
crystal, where the density and the potential have one component apiece and a
field has nothing to be added to; give the unit cell a small
\code{starting\_magnetization} so that it carries four, which it may then
converge back to nearly zero, and what the run measures is the Pauli
susceptibility. A ground state converged \textbf{under
a magnetic field or a constrained moment}, because the frozen eigenvalues would
carry the field the SCF ended with --- which \code{reducebf} and the
fixed-spin-moment scheme both change from the input --- while $\delta V$ is
rebuilt here from the density alone; apply the field through
\code{magnetic\_field} instead. DFT+U, the occupation matrix being per atom
where an ultracell has $N$ copies of each. Spin spirals, as a second modulation of the same states.
Gamma-only storage, the basis being built at $N$ points none of which is
$\Gamma$. A double grid is \emph{not} refused, and what remains of the refusal
that used to stand here is a warning firing the other way: \code{ecutrho}
defaults to four times \code{ecutwfc} here as it does in \code{pw.x}, so an
augmented input that simply leaves it out runs with the augmentation charge on
the wavefunction grid, which is the one thing such a dataset exists not to do. Symmetry is not used at all: a modulation breaks the
crystal's point group, so the run needs \code{nosym} and the whole k-grid. And a
non-integer ultracell is refused, where Elk's \code{avecu} is an independent
input with no consistency check against its Q-grid. An \textbf{unconverged}
unit-cell ground state, because its density is the one thing an ultracell holds
fixed and there is no later iteration in which a bad seed could work itself out.

Two things are \emph{not} refused and are limits on a claim rather than on a
run. The total energy's monotone bound rests on the free energy being minimised,
which Methfessel-Paxton smearing breaks --- its occupation function is
non-monotone, so the generalised entropy is not concave and the self-consistent
point is not a minimum; gaussian, Fermi-Dirac and Marzari-Vanderbilt are safe,
and the energy itself is computed either way. And a cell with more than one
self-consistent solution has no ladder at all: the energy then \emph{ranks}
solutions at one \code{nbnd} and does not certify a sequence across several,
since nothing says the same solution is found at each.
\end{refuses}

% =====================================================================
\section{What a modulation looks like to a tip}

A charge or spin density wave, a screened impurity, a domain wall: what all of
these are actually seen with is a scanning-tunnelling microscope, so an
ultracell that cannot be imaged is a calculation of something nobody measures
directly. This section is the join between the ultracell and the two things a
tip measures, the Tersoff-Hamann image and the vertical tunnelling
transmission. The third, the \emph{spectrum} at one place over many biases, is
the same relabelling read along an energy axis and has its own subsection under
\emph{Tunnelling spectra} --- on a modulation it is usually the one to reach
for, since a wave moves a band edge and an image at one bias cannot tell a level
that moved from a level that gained weight.

Nothing new is computed, and the reason is one observation about the basis. An
ultracell state is
\begin{equation*}
  \Psi_j(\mathbf r) = \sum_{\mathbf Q, n} a^j_{\mathbf Q n}\,
                      \psi_{\mathbf k_0 + \mathbf Q, n}(\mathbf r),
\end{equation*}
and in a plane-wave basis that sum is a relabelling rather than a transform: a
plane wave with unit-cell Miller index $\mathbf G$ belonging to $\mathbf Q$ is
the ultracell plane wave $\mathbf h = n\mathbf G + \mathbf q$, and no two of
them collide, because reducing $n(\mathbf G - \mathbf G') + (\mathbf q -
\mathbf q') = 0$ modulo $n$ forces $\mathbf q = \mathbf q'$ and the sphere sits
inside the dense grid. So the union over $\mathbf Q$ of the $N$ unit-cell
spheres \emph{is} the ultracell's own sphere at $\mathbf k_0$, the band index
contracts away, and what is left is an ordinary coefficient vector that every
routine written for a unit-cell wavefunction takes unchanged once it is handed
the ultracell's Miller indices, its cell and $\mathbf k_0$ in its reciprocal
coordinates. On a modulated two-cell silicon ultracell the state sampled from
that vector agrees with the box transform the density is built from to
$1.7\times10^{-16}$ against a peak of $0.14$, and it integrates to 1 over the
ultracell.

The image is then the ultracell density with a smeared delta at the tip energy
in place of the occupations, which is what a Tersoff-Hamann image is in Elk's
\code{wfplot.f90} and in QE's \code{stm.f90}, and the transmission is the same
contraction \code{run\_vertical\_transport} uses with the ultracell's geometry
in it. The entry points are \code{run\_ultracell\_stm} and
\code{run\_ultracell\_transport}, or \code{Calculator.get\_ultracell\_stm} and
\code{Calculator.get\_ultracell\_transport}.

\begin{pycode}
from defumat import Calculator
import numpy as np

cal = Calculator.from_file("si.in")        # the two-atom unit cell

# a potential of one ultracell period along a_1, in Ry
cal.get_ultracell(supercell=(8, 1, 1), kgrid=(1, 2, 2), nbnd=24,
                  external=lambda x: 0.05 * np.cos(2 * np.pi * x[..., 0] / 8))

img = cal.get_ultracell_stm(height=0.35, axis=2, shape=(96, 12),
                            bias=-1.5, width=0.01)
print(img.values.shape)        # (96, 12): the plane spans all eight cells
print(img.integral)            # 8.0039: the electrons in the window, per cell
\end{pycode}

\textbf{Two conventions, and the second is the one to know.} The plane is given
in \emph{unit-cell} crystal coordinates running over $[0, n_i)$ across the
ultracell, which is what \code{external} already takes, so \code{height = 0.35}
means what it means in a unit-cell run and a map spans the whole modulation by
default. And the image is a tunnelling density of states \emph{per unit cell},
in the same units a unit-cell image reports, so the two are directly
comparable: with nothing applied, the ultracell image is the unit cell's image
tiled, to $1.6\times10^{-7}$ of the peak at $N = 2$. The window above covers the
whole valence band, so the integral comes back as silicon's eight valence
electrons; a run with a Fermi level puts the tip at it instead, and one with
fixed occupations warns and takes the middle of the gap, which is
\code{stm.f90}'s own rule.

\textbf{A magnetic tip is what a spin density wave needs.} A collinear system is
invariant under flipping every spin together with the sign of $B$, so the charge
cannot respond at first order and a wave driven by \code{magnetic\_field} is
almost invisible in the charge \emph{density}: on an eight-cell silicon
ultracell carrying 0.12 Bohr magnetons the total density varies by
$1.9\times10^{-4}$ of its mean from cell to cell. What an unpolarized tip sees
is neither that small nor the wave itself, but the wave \emph{squared}, two
periods where the wave has one, at 37 per cent of its mean against a polarized
tip's 82, since the leading charge response is at twice the wavevector and the
tip energy is a smeared band edge rather than a count of electrons.
\code{spin} and \code{polarization} project the channel, exactly as they do for
a unit-cell image.

\textbf{The number.} The reference is the one the ultracell itself is checked
against: a real $N$-cell supercell under the same applied potential, run through
this package's own SCF and imaged the same way. The error on the induced
corrugation falls with \code{nbnd}, which is the whole content of "a variational
truncation to \code{nbnd} bands per folded k-point": 22.2, 5.2 and 2.0 per cent
at \code{nbnd} $= 12$, 24 and 48 on a two-cell silicon ultracell, where the
induced corrugation is itself 0.6 per cent of the image. Widening the exit
region of the transmission to the whole ultracell must give the image back
exactly, with no factor between them, and on a norm-conserving dataset does so
to $10^{-11}$; that identity is what holds the two normalisations against each
other, since the image carries its $N$ in the weights and the transmission
carries it in the volume the states are normalised over.

\textbf{An ultrasoft or PAW dataset needs nothing extra, and where the planes
are is the reason.} A tip point and an exit plane both sit in the vacuum, where
the augmentation charge is zero and a pseudo-wavefunction is the true one, so the
sampled amplitude and the plane's Gram matrix are exact as they stand --- which
is the same argument a unit-cell junction rests on, and it has always run on such
a dataset. The exception is the \emph{whole-cell} exit region: there the Gram
matrix is $\langle\Psi|S|\Psi\rangle$ rather than an integral over a plane, and
$S$ for an ultracell state spans $N$ unit-cell spheres with their own projectors,
so \code{exit\_region = "volume"} is refused by name. It is refused rather than
returned, and on purpose: the frozen basis is exactly $S$-orthonormal across
$\mathbf Q$, so that Gram matrix \emph{is} the identity by construction, and
handing the assembly the identity would make the diagnostic agree without looking
at the coefficients it exists to check.

\textbf{What such a run needs is a cell with vacuum in it, which a bulk crystal
is not.} Silicon's PAW dataset reaches 2.647~bohr where the interplanar spacing
along the stacking axis is 5.889~bohr with an atom every quarter of it, so the
widest clearance any plane can have is 2.208~bohr and the guard refuses every
height --- correctly, since a substrate and a tip belong either side of a
\emph{sheet}. On the square silicon sheet committed as
\code{ultracell-sheet-paw.in}, with the two planes at 0.15 and 0.85 and the atom
at 0.5, the unmodulated two-cell ultracell reproduces the unit cell's map tiled
to $6.2\times10^{-7}$ (ultrasoft) and $9.7\times10^{-7}$ (PAW) of the peak, and
under a $0.05\cos(\pi x_1)$~Ry modulation the map converges to a real two-cell
supercell's own transmission: $1.09\times10^{-2}$ at \code{nbnd = 12} and
$1.96\times10^{-3}$ at 24. The modulated rung is the one that discriminates,
because an unmodulated folded state is a single basis function whose
normalisation is the unit cell's, while a state that mixes $\mathbf Q$ makes
$\sum_{\mathbf Q}|a_{\mathbf Q}|^2 = 1$ the $S$-orthonormality claim itself.

\begin{pycode}
from defumat import Calculator

cal = Calculator.from_file("ultracell-sheet-paw.in")   # a PAW sheet in vacuum

cal.get_ultracell(supercell=(2, 1, 1), kgrid=(2, 4, 1), nbnd=12)

t = cal.get_ultracell_transport(exit_height=0.15, height=0.85, axis=2,
                                shape=(12, 4), broadening=0.05)
print(t.values.shape)            # (12, 4): the map spans both cells
print(float(t.values.max()))     # 4.1856e-05, in the transmission's own units
print(t.notes["hermiticity"])    # 0.0 -- the Gram matrices, not imposed
\end{pycode}

\textbf{Sampling the plane is exact and its set is the ultracell's own dense
sphere}, $|\mathbf G + \mathbf Q|^2 \le \code{ecutrho}$, rather than Elk's
choice of the unit cell's dense sphere at every $\mathbf Q$. Read back at the
box's own grid points, the sphere reproduces a modulated silicon density to
$3\times10^{-16}$ where the other set is out by $5\times10^{-9}$ at $N = 2$ and
$2\times10^{-8}$ at $N = 3$. The box itself aliases its outermost shell from
about $N = 5$, at $10^{-11}$, which is the ordinary $\code{ecutrho} = 4\,
\code{ecutwfc}$ corner aliasing one level out and is a property of the
ultracell's grid rather than of the sampling.

\begin{refuses}
Everything the ultracell itself refuses, since the run has to exist first. A
result whose states were dropped (\code{keep\_states = False}), which cannot be
rebuilt, the diagonalisation on the folded k-set being the expensive step of the
whole method. For the transmission, an ultracell that is more than one cell deep
along the stacking axis: the exit plane would then sit between two of the
stacked slabs, so the electron leaves through the material rather than out of
it, and two $\mathbf Q$ differing only along that axis share every in-plane
index the exit integral has to tell apart. The same axis must carry one k-point,
for the reason a unit-cell junction needs one. A denser k-set for the image, the
way \code{grid} re-solves the bands for a unit-cell image: on an ultracell that
is not post-processing but another run of the whole loop, so it is asked for
with \code{kgrid} in \code{get\_ultracell} instead. And a constrained
\code{tot\_magnetization}, whose two channels have their own Fermi levels where
a tip sees one energy. On an ultrasoft or PAW dataset: a plane inside an
\textbf{augmentation sphere}, which the guard reports the distance for, and the
transmission's \textbf{whole-cell} exit region, for the reason given above ---
the transmission itself, the image and the spectrum all run. An applied \code{magnetic\_field} is \textbf{not} refused
here, where a unit-cell image refuses it: there the field is something the
ground state was converged under and the reported total does not carry, while
here it is the modulation's own driver and the eigenvalues the tip energy is
read from are the eigenvalues under it.
\end{refuses}

% =====================================================================
\section{Continuing an all-electron ground state}

Elk is an all-electron LAPW code and this is a pseudopotential plane-wave one,
so their converged densities are different objects: Elk's carries the core
electrons and the nuclear cusp, and this one carries a smooth pseudo valence
density normalised to the valence electron count. What crosses between them is a
\emph{starting guess}, and the run's own SCF still happens on top of it.

Inside each muffin tin Elk stores the density as a real-spherical-harmonic
expansion on a logarithmic radial mesh, and everywhere else as a periodic
function on a uniform grid,
\begin{equation*}
  \rho(\mathbf r) =
  \begin{cases}
    \sum_{lm} f_{lm}(r)\,R_{lm}(\hat{\mathbf r}), & |\mathbf r - \mathbf R_a| < r_a,\\[4pt]
    \sum_{\mathbf G} \tilde\rho(\mathbf G)\, e^{i\mathbf G\cdot\mathbf r}, & \text{otherwise,}
  \end{cases}
\end{equation*}
so putting it on a plane-wave grid means deciding which region each grid point
is in, then interpolating radially or summing the Fourier series. The muffin-tin
part is written in pointwise and is never Fourier truncated, which is why the
nuclear cusp transfers exactly and why nothing is lost to aliasing.

\begin{pycode}
from defumat import Calculator

calc = Calculator.from_file("scf.in", pseudo_dir="pseudo")
seed = calc.get_elk_seed("elk_run/")     # STATE.OUT beside GEOMETRY.OUT
result = calc.get_scf(starting_density=seed)
\end{pycode}

The argument is an Elk \emph{run directory} rather than a file:
\code{STATE.OUT} carries neither the lattice vectors nor the atomic positions,
so \code{GEOMETRY.OUT} has to be beside it. And \code{elk.in} is not a
substitute, because Elk's default \code{tshift} moves the origin onto the
inversion centre and only \code{GEOMETRY.OUT} is written in that shifted frame.
Pass \code{report=[]} to get back a \code{SeedReport}, which is what the
reconstruction measured on the way: what it integrated to before renormalisation,
how many of Elk's reciprocal lattice vectors this run's density cutoff kept, and
what fraction of the norm sat outside it.

\code{ElkState} is the reader on its own, for looking at the state rather than
seeding with it. \code{read\_elk\_state} is the same thing as a function and
\code{read\_elk\_geometry} reads \code{GEOMETRY.OUT} alone, returning the
\code{ElkGeometry} that \code{state.geometry} also holds; all four are exported
from \code{defumat.io}.

\begin{pycode}
from defumat.io import ElkState

state = ElkState.read("elk_run/")
print(state.rmt, state.efermi)            # muffin-tin radii; e_F in Rydberg
print(state.muffin_tin_charges())         # the charge inside each sphere
print(state.interstitial_charge())        # and the charge outside all of them
rho = state.evaluate_at(points)           # the density at arbitrary points
\end{pycode}

The pieces of the reconstruction are exported too, from
\code{defumat.io.elk\_density}, because each is a separately checkable part of
somebody else's convention rather than an internal:
\code{real\_spherical\_harmonics} (the sign convention is Elk's and is not the
textbook one), \code{poly4} (the 4-point Lagrange radial window),
\code{spline\_weights} (the quadrature the radial integrals use, where Simpson's
rule is 9e-6 away), \code{characteristic\_function} (the Fourier-truncated step
that separates the two regions) and \code{interstitial\_coefficients}
(\code{rhoir} as a G-vector list). \code{density\_on} is what
\code{get\_elk\_seed} calls.

\textbf{What it is worth, measured.} On simple-cubic hydrogen at $a = 3.0$~bohr,
the reconstruction agrees with Elk's own evaluation of its own density at every
one of the 4096 points Elk printed, to $4.5\times10^{-11}$, which is the
precision the file was printed at rather than any error in the transfer. The
charge inside the sphere is $0.6125761996$ and outside it $0.3874238004$, both
exact to every digit Elk printed. And a run started from that density and a run
started from a superposition of atomic charges converge to the same state:
$-1.080181442650$~Ry either way, agreeing to $6\times10^{-14}$~Ry with converged
densities $9\times10^{-8}$ apart.

\textbf{It buys no iterations, and that is the honest result.} Both runs take
four. At a lower cutoff the seeded one takes \emph{one more}. For a hydrogen
atom a superposition of atomic charges is already almost the answer, and the
all-electron cusp inside the pseudisation radius is a place where Elk's density
is further from the pseudo one than the atomic guess is. The reader is worth
having because it bridges to an all-electron ground state, not because it is
faster.

\begin{refuses}
\textbf{Refuses.} A \textbf{spin-polarized} Elk state, and a
\textbf{spin spiral} with it --- a magnetization is a second field with a
transfer rule of its own, and a seed that carries the charge while silently
dropping the moment would start a magnetic run from a non-magnetic guess. A
\textbf{DFT+U} or fixed-tensor-moment state, whose occupation matrices are
trailing records nothing here reads. A state written by Elk older than
\textbf{2.0.0}, which Elk's own reader refuses. A directory with no
\code{GEOMETRY.OUT}, or one whose geometry is from a different run. A
calculation on a \textbf{different cell} --- a density is transferable between
identical lattices and there is no interpolation here that would make anything
else mean something. A calculation with more than one spin channel. And an
element with a \textbf{core}: the muffin-tin density includes the core states
and nothing in \code{STATE.OUT} separates them from the valence charge, so a run
whose valence electron count is far from what the state integrates to is refused
by name rather than scaled to fit. Hydrogen, with no core at all, is what this
path is for.

Passing an Elk density as a \textbf{fixed} density, with no SCF above it, is not
offered: the two codes' converged densities are different functions wherever
there is a core, and no splice of the two would be a solution of anything.
\end{refuses}

% =====================================================================
\section{Performance: dials, memory and GPUs}

\textbf{Two batching dials} control how much is in flight, and both defaults
follow the platform: on a CPU, QE's loop --- one $k$-point and one band at a
time, which is what a \emph{cache} wants --- and on a GPU the whole of both
axes, which is what a card wants. Either end is reachable everywhere:

\begin{shcode}
DEFUMAT_K_BATCH=all DEFUMAT_BAND_BATCH=all python3 run.py
\end{shcode}

or per call, \code{run\_scf(..., k\_batch=...)}. The chunk size changes the order
contributions are summed in and nothing else --- round-off, $\sim$1e-15 Ry.

\textbf{The band dial is worth twice as much in a spinor run}, because a band
in flight is then \emph{two} real-space boxes rather than one: a noncollinear
calculation with 403 bands on a $240\times54\times200$ grid puts
$403\times2\times2{,}592{,}000$ complex numbers --- 33 GB --- into one transform
if the whole block goes at once, and several such arrays are live inside a
single application of $H$. \code{DEFUMAT\_BAND\_BATCH=16} makes that 1.3 GB. The
answer is unchanged bit for bit, not merely to round-off, because chunking the
band axis reorders nothing: it is a map, not a sum.

\textbf{Why the default is per platform.} A GPU has no cache to fit in and
inverts the conclusion: on a ten-atom metal, \code{k=1,b=1} costs \textbf{4.5x}
what \code{k=all,b=all} does on the same card, sixteen-atom silicon runs at
\textbf{0.20x} --- an outright loss against four CPU cores --- and at
thirty-two atoms \code{b=1} did not finish at all. So a run on a device gets
both batched unless it says otherwise, and the plausible half-measure
\code{k=all,b=1} is never the default: it is the \emph{worst} setting
available, because it buys the batched mode's memory with the looped mode's
kernel launches. An explicit argument beats the environment variables, which
beat the platform.

\textbf{A third dial says where the wavefunctions are \emph{kept}}
(\code{wfc\_store}, \code{DEFUMAT\_WFC\_STORE}), which is a different question
from how many are in flight. The store is $n_{\rm spin} n_k n_{\rm bnd}
n_{\rm pwx} n_{\rm pol}$ complex --- 11.3~GB on a 45-atom slab at six
$k$-points, 45~GB at twenty-four --- and it stays resident all through the SCF
because it is the next iteration's starting guess. \qe{} does not keep it on the
device: \code{c\_bands.f90:get\_buffer} and its \code{save\_buffer} partner hold
one $k$-point's worth and \code{io\_level} says where the rest lives.
\code{wfc\_store='host'} is that buffer.

\emph{What it buys, and what it does not.} Parked, the store is off the
accelerator while the energy terms, the potential on the dense grid, the mixer's
history, the PAW one-centre terms and the checkpoint write are running ---
which is where the other large allocations are. It is \emph{on} the device for
the diagonalisation, for \code{becsum} and for the density, because it is an
input to all three, so this does \textbf{not} lower a peak that sits inside the
eigensolver. The default follows the platform for a measured reason: on a CPU
the host \emph{is} the device, and a round trip of the store runs at 1.3~GB/s
here, which at slab scale would be 18.6~s per iteration for nothing. On an
accelerator the same transfer is PCIe, around a second. So \code{device} on a
CPU, \code{host} on anything else, and the answer is bit-for-bit identical
either way --- a transfer is not arithmetic.

\textbf{A fourth dial says whether the \emph{projectors} are kept}
(\code{projectors}, \code{DEFUMAT\_PROJECTORS}), and it is the one that has been
measured to move a peak rather than a stage. $\langle k+G|\beta\rangle$ is
$n_k n_{\rm pwx} n_{\rm kb}$ complex and resident for the whole run: on the
45-atom slab above that is \textbf{13.96~GB}, larger than the wavefunctions
themselves. \qe{} does not hold it either --- \code{c\_bands.f90:111} calls
\code{init\_us\_2} \emph{inside} its $k$-loop --- and \code{projectors='rebuild'}
is that: only the phase-free core is kept, one column per \emph{species} channel
where the full array has one per \emph{atom} channel, and each $k$-point's is
formed on demand. The expression is the same one either way, so the result is
bit-for-bit identical.

\emph{Why it moves a peak where the store dial does not.} The other three dials
bound what is \emph{in flight} or move what is held between two points; this one
removes what is \emph{held}, so it sits underneath every stage. Measured on that
slab, two runs on the same card differing only in this flag: every stage bracket
falls by 13.96~GB, and the peak with them --- \textbf{79.14~GB to 65.18~GB},
with the converged energies identical to the last printed digit. \textbf{The
saving is the whole of $\langle k+G|\beta\rangle$}, not that array minus a core:
the core is built unconditionally and both routes pay for it, so what the dial
chooses between is the whole-$k$ array and one rebuilt chunk. The saving is
therefore $n_k/k_{\rm batch}$ times the chunk --- it grows with how many
$k$-points are \emph{not} in flight, and at a single $k$-point it is exactly
neutral, buying nothing and costing nothing. It is the ratio $n_{\rm kb}/n_{\rm
cs}$, atoms per species, that decides how much smaller the core is than the
array, which is why the core stays cheap on a large cell. Measured cost:
$n_{\rm at} n_{\rm pwx}$ complex exponentials each time a $k$-point's projectors
are used, which on that slab is \textbf{+1.5\% per SCF iteration} (46.1~s against
45.4~s) and 10~s more to compile. The default is \code{store}.

\begin{refuses}
Nothing is refused, but two things are worth knowing before switching. The
default is \code{store} and every validated number in this document was measured
with it. And a size estimate follows the dial only from \code{433abd7} onward:
an older \code{defumat size} reports the stored array for a rebuilt run and so
overstates the floor by the largest single line in its own table.
\end{refuses}

\textbf{Measured GPU speedups} (one H200 against one CPU core, per SCF
iteration): 2x on a cheap metal, 3--5x on norm-conserving silicon, 9--21x with
an augmentation charge or a magnetization, and \textbf{83--115x on spin-orbit}.
Energies agree to $\sim$1e-9 Ry --- the same agreements the CPU already has.

\textbf{Memory.} Spin-orbit is where it bites: a 20-atom bismuth cell with a
relativistic ultrasoft dataset peaks at \textbf{34.7~GB}, which fits a 141~GB
card and does not fit a 32~GB one. Everything else in that sweep sits under
2~GB.

\textbf{A response costs memory according to how it is differentiated}, which is
worth knowing before starting one. A \emph{forward} response --- the Sternheimer
solves behind $\epsilon$ and $Z^*$ --- costs 0.7--1.0~GB over the SCF that
produced the state, and a \emph{forward-over-reverse} one --- the dynamical
matrix, the Raman tensors --- 1--2~GB. The \emph{stress} is the exception and it
is a \emph{reverse} derivative through the radial transforms: \textbf{10.5~GB}
on eight-atom ultrasoft silicon, ten times any of the others, and an order more
than the SCF it came from. Measure your own with
\code{tools/gpu/phase5.py <case> -{}-property <name>}.

\textbf{The augmentation charge is the largest object in an ultrasoft or PAW
run on a large cell}, and there are two ways of keeping it. Below
\code{DEFUMAT\_AUG\_MAX\_BYTES} (default 2~GB) $Q_{ij}(\mathbf{G})$ is built on
the whole $G$ sphere once and kept, which is what every number in the accuracy
table below was measured with. Above it, \qe's scheme takes over: the
\emph{radial} transforms
%
\begin{equation*}
  Q^L_{nm}(q) \;=\; \frac{4\pi}{\Omega}\int_0^{r_{\mathrm{kkbeta}}}\!\!
  \mathrm{d}r\; j_L(qr)\,\bigl[r^2 Q^L_{nm}(r)\bigr]
\end{equation*}
%
are tabulated on a grid in $|q|$ at $\mathrm{d}q = 0.01$ and
$Q_{ij}(\mathbf{G})$ is rebuilt from them a block of $\mathbf{G}$ at a time,
every iteration, exactly as \code{addusdens} and \code{newd} do it. On a
45-atom slab with 21~\AA\ of vacuum that is \textbf{8.4~MB} against
\textbf{65~GB} for one relativistic nickel dataset.

\begin{shcode}
DEFUMAT_AUG_MAX_BYTES=off python3 run.py   # never tabulate
DEFUMAT_AUG_MAX_BYTES=2G  python3 run.py   # the default
DEFUMAT_AUG_CHUNK=8192    python3 run.py   # G vectors per block
\end{shcode}

\code{DEFUMAT\_AUG\_CHUNK} sizes the rebuild's working array,
$(n_h, n_h, \text{chunk})$, and defaults to putting it near 256~MB; like the
batching dials it is a loop bound over an exact sum and is invisible in a
result. The two schemes agree to \textbf{4e-10~Ry} on a total energy and
exactly at $\mathbf{G} = 0$. On eight-atom ultrasoft silicon the two
contractions cost about \textbf{4.5x} more from the table, and the whole run
comes out level, because the table is cheaper to build than the stored array
is.

\begin{refuses}
\textbf{Past the end of the table is NaN, not an extrapolation.} The table
reaches $2\sqrt{\code{ecutrho}}$ --- \qe's own \code{cell\_factor} --- so that a
strained cell, whose $G$ vectors move while the $G$ \emph{set} does not, still
lands on it. A strain large enough to take $|\mathbf{G}|$ past that gives NaN
rather than a clamped value, because a clamped $Q^L(q)$ would be smooth, of the
right order and wrong, and would show up only as a bad stress.
\end{refuses}

\textbf{A PAW run's other large object is its one-centre quadrature, and what
it costs is a \emph{derivative} rather than a run.} Each atom's sphere carries
$(n_{\mathrm{spin}}, n_x, \text{mesh})$ radial fields, and inside a force or a
stress every atom's reverse pass is live at once: \textbf{272~MB per atom} for a
relativistic nickel dataset with a magnetization, so 4.1~GB for a fifteen-atom
sublattice and no ceiling. \code{DEFUMAT\_PAW\_ATOM\_BATCH} is how many spheres
are in flight together --- \code{all} for one \code{vmap}, an integer for
\code{paw\_onecenter.f90}'s own loop over atoms:

\begin{shcode}
DEFUMAT_PAW_ATOM_BATCH=1   python3 run.py   # one sphere at a time
DEFUMAT_PAW_ATOM_BATCH=all python3 run.py   # the whole sublattice at once
\end{shcode}

Left alone it chunks above three atoms per species and not below, because that
is where the crossover was measured: chunked, the tape stops growing with the
sublattice altogether (4.1~GB $\to$ 0.80~GB at fifteen nickel atoms) and the run
gets \emph{faster} --- a ten-atom PAW cell's force is 0.505~s against 0.418~s
and its SCF iteration 1.40~s against 1.28~s. On two atoms it is the other way,
0.761~s against 1.286~s, which is what the threshold is for. Like the other
batching dials it is a loop bound over an exact sum and is invisible in a result
beyond round-off.

\textbf{Precision.} Every dtype comes from \code{config.Precision}; x64 is
enabled before any array exists, and all validation is float64. float32 is a
performance mode and \textbf{never one a correctness claim is made in}.

% =====================================================================
\section{Accuracy summary}

Against \qe\ on the same input:

\begin{center}
\small
\begin{tabular}{@{}lr@{}}
\toprule
\textbf{quantity} & \textbf{agreement} \\
\midrule
total energy, norm-conserving LDA & $\sim$1e-9 Ry \\
ultrasoft, PAW & $\le$3e-9 Ry \\
PBE / revPBE / PBEsol & $\le$6e-9 Ry \\
metals, every smearing and tetrahedra & 2.5e-8 Ry \\
collinear spin & 1.2e-9 Ry \\
noncollinear magnetism & 2.8e-9 Ry \\
spin-orbit coupling & $\le$1.3e-8 Ry \\
DFT+U & $\le$6.7e-9 Ry \\
\midrule
band structure & 0.0002 eV \\
forces & $\le$2e-5 Ry/bohr \\
stress (13 cases) & $\le$2.7e-7 Ry/bohr$^3$ \\
relaxed geometry & 1e-6 bohr \\
\midrule
dielectric constant (NC, US, PAW) & $\le$1.2e-4 \\
Born effective charges (NC) & every printed digit \\
Born effective charges (US, PAW) & $\le$4.9e-5 \\
phonons at $\Gamma$, silicon & 0.05 cm$^{-1}$ \\
phonons at $\Gamma$, silicon (US, PAW) & 0.02--0.03 cm$^{-1}$ \\
phonons at $\Gamma$, aluminium (metal) & 0.002 cm$^{-1}$ \\
Raman/IR spectra vs \code{dynmat.x} & every printed digit \\
\bottomrule
\end{tabular}
\end{center}

\textbf{One case does not close, and is recorded rather than explained.}
\code{bi10-soc} --- ten bismuth atoms with a fully-relativistic dataset --- sits
\textbf{1.9e-4 Ry} from QE on a total of $-1477.737$, which is 1.3e-7 relative.
Both codes choose the same symmetry operations, the same $k$-point, the same
grid and the same $G$-vectors, and both converge. It is in the test suite at
that tolerance, not hidden.

% =====================================================================
\section*{Where to go next}
\addcontentsline{toc}{section}{Where to go next}

\begin{itemize}\itemsep2pt
\item \code{README.md} --- the short tour and a first calculation
\item \code{notebooks/} --- 26 worked examples on concrete systems, executed and
  committed, each with a \code{.md} export beside it
\item \code{PERFORMANCE.md} --- the running performance log, CPU and GPU
\item \code{PLAN.md} --- the architecture, the phase breakdown, the trap catalogue
\item \code{GPU.md} --- the GPU roadmap
\end{itemize}

\end{document}
